\chapter{Proofs}
\label{ch:proofs}

\section{Sequents and deduction graphs}\label{section:sequents}

A \emph{sequent} is a pair $(L, B)$, where $L$ is a list of formulas
called \emph{assumptions} and $B$ is a single formula called the
\emph{assertion}.  A sequent is represented as
\[
  A_1,\, A_2,\, \ldots,\, A_n \;\Longrightarrow\; B
\]
where $A_1, \ldots, A_n$ are the assumptions (or hypotheses or the context) and $B$
is the assertion. In the course of a proof, many sequents may be created.

A \emph{deduction graph} is a directed acyclic graph of
\emph{sequent nodes} connected by \emph{inference nodes}.  An inference
node records which rule was applied, which sequent nodes are its
hypotheses, and which sequent node is its conclusion.

A sequent node is \emph{grounded} when it has an inference node all of
whose hypothesis nodes are themselves grounded.  Leaf inferences (with no
hypothesis nodes) are immediately grounded.  Grounding propagates upward.
A proof is \emph{complete} when the root sequent node is grounded.

\begin{remark}
A \emph{proof state} (\texttt{<proof-state>}) is not the same thing as a
deduction graph.  It is a thin wrapper that bundles three items together: the
deduction graph itself, the \emph{root} sequent node (the original goal,
fixed at \texttt{start-proof} time and checked by \texttt{qed}), and the
\emph{focus} sequent node (the subgoal currently being worked on).  The
deduction graph is the mathematical object that grows as rules are applied;
the proof state is the navigational view into it that the interactive session
and the \texttt{cmd-*} commands operate on.  In the interactive session,
\texttt{*ps*} always holds the current proof state.
\end{remark}

%% =========================================================
\section{Kernel operations}\label{sec:kernel-boundary}

A proof in VNB is carried out on a deduction graph, and the graph is
changed only by \emph{kernel operations}.  A kernel operation is applied
to one sequent node $N$.  It adds one inference node: the conclusion is
$N$, the rule recorded is the name of the operation, and the hypotheses
are sequent nodes that the operation builds from $N$ in a way fixed by
the operation.  Nothing else changes the graph.

There are \VNBkernelopcount{} kernel operations.  The list below is
produced from the list the system itself checks each time it starts
(Section~\ref{sec:kernel-checked}).

\VNBkerneloptable

The last \VNBkernelopcountoracle{} are \emph{oracles}.  An oracle runs a
decision procedure, written in Scheme, on the assertion of $N$ and some
of its assumptions.  If the procedure succeeds, the operation adds an
inference node, with no hypotheses except in the case of
\texttt{arith-simplify}, which posts one.  The graph records the name of
the oracle; for \texttt{ineq} and \texttt{sos} it also records the
certificate the procedure found.  The inference is written only if the
oracle's checker accepts it (Section~\ref{sec:inference-checked}).

The kernel operations, together with the axioms
(Section~\ref{sec:proof-debt}), are everything VNB assumes.  A result
reported \texttt{modulo 0} rests on these and on nothing else.

Every other proof command is a \emph{tactic}.  A tactic is a Scheme
procedure that calls kernel operations, possibly many, possibly after a
search.  It has no other access to the graph.  A tactic can fail to find
a proof.  It cannot add an inference that no kernel operation would add.
Section~\ref{sec:tactics-ref} describes the tactics, and
Section~\ref{sec:command-uses} lists, for each command, the kernel
operations it can reach.

\begin{remark}
One command may record different operations according to the form of the
assertion: \texttt{di} records \texttt{and-intro}, \texttt{implies-intro},
\texttt{forall-intro} and so on.  Two commands may record the same
operation.  The name recorded in the graph is that of the operation, not
of the command.
\end{remark}

\subsection{What the operations do}\label{sec:kernel-ops-detail}

%% Inference figures.  \VNBfig{operation}{hypotheses}{conclusion}, and
%% \VNBfigS with the hypotheses stacked one per row (separated by \\); the side
%% conditions follow in a VNBcond list.  Defined here because this is the only
%% section that uses them.
\providecommand{\VNBfig}[3]{\[\frac{\displaystyle #2}{\displaystyle #3}\;\;\mbox{\texttt{#1}}\]}
\providecommand{\VNBnohyp}{\phantom{\Gamma}}
\providecommand{\FV}{\mathrm{FV}}
\providecommand{\aeq}{=_{\alpha}}
\providecommand{\VNBfigS}[3]{\[\frac{\displaystyle\begin{array}{c}#2\end{array}}{\displaystyle #3}\;\;\mbox{\texttt{#1}}\]}
\newenvironment{VNBcond}{\sloppy\begin{description}[font=\normalfont\itshape,leftmargin=0.85cm,itemsep=0.5mm,topsep=0.5mm]}{\end{description}}

In this section $N$ is the node the operation is applied to, with
sequent $\Gamma \Longrightarrow G$.  ``Hypotheses'' are the sequents of
the new hypothesis nodes.  $B[x{:=}t]$ is $B$ with $t$
substituted for the free occurrences of $x$, bound variables being
renamed where needed to avoid capture.  $\FV(E)$ is the set of variables
free in $E$, and $\FV(\Gamma)$ the union over the formulas of $\Gamma$.
An operation that \emph{declines} leaves the graph unchanged.  Posting a
sequent that the graph already contains, up to renaming of bound
variables, returns the existing node, so a new hypothesis may already be
grounded.

An operation that takes an index or a pattern records it beside its
name: the graph holds \texttt{(cartesian-elim 2)}, \texttt{(macete
\textit{name} $L$ $R$)}, \texttt{(ineq \textit{certificate})}.  The name under which the inference is counted, and under which
it appears in the list above, is the leading symbol.  The arguments are
data about one application.

\emph{Definedness.}  VNB terms may fail to denote, and the sign of
denotation is $t = t$: strict equality holds only between defined terms
(Chapter~\ref{sec:language}).  Two operations, \texttt{forall-elim}
and \texttt{reflexivity}, therefore consult a \emph{definedness
certificate}, one test applied to a term and the assumptions of $N$.
$\Gamma$ certifies $t$ when $t$ is a variable or an atomic constant, a
numeral or a ground arithmetic term, a term that $\Gamma$ types
($t \in X$) or equates ($t = u$), a term occurring outside every binder
in an assumption of the form $\in$, $=$, $\leq$ or $<$, a total class
constructor applied to certified arguments, a separation, comprehension
or indexed union whose domain is certified, an accessor of a structure
$\Gamma$ carries, an applied structure operation whose arguments
$\Gamma$ types in its declared domain, or $f(a)$ with $f \in
\mathrm{FUN}(D, \cdot)$ and $a \in D$ in $\Gamma$.  A term whose head is
\textsc{choice} or \textsc{iota}, and an application of an untyped
function, are never certified.  The policy is
\texttt{docs/definedness-instantiation-2026-09-18.md}.

\subsubsection*{Reading the figures}

Each operation is stated below as an inference figure.  The sequents above
the line are the hypotheses, the sequent below it is the sequent of $N$,
and the labelled conditions under the figure are the side conditions.  A
figure is the relation that the operation's \emph{checker} tests before
the inference is written (Section~\ref{sec:inference-checked}): an
inference of that form that meets the conditions is accepted, whichever
procedure produced it, and any other is refused.  The operation itself
produces one instance of its figure.  Where it has a choice to make
(which occurrences to rewrite, what to call a new variable), the prose
after the figure says which choice it makes.  A condition labelled
\emph{Not checked} is imposed by the operation but not tested by its
checker.

The figures share five conventions.
%% the conventions: header of rule-checkers-logic.scm, lines 14-27;
%% chk-set=? :82, chk-ctx+? :100, chk-ctx-+? :103, chk-arity? :133,
%% chk-match-under :167, chk-adds-nothing? :126
\begin{itemize}
\item \emph{Contexts are sets.}  $A \aeq B$ means that $A$ and $B$ are
  the same formula up to renaming of bound variables.  Two contexts are
  equal when every formula of each is $\aeq$ to a formula of the other.
  $\Gamma, A$ is $\Gamma \cup \{A\}$, so it is $\Gamma$ itself when
  $\Gamma$ already contains $A$; $\Gamma - F$ is $\Gamma$ with every
  formula $\aeq F$ removed.  The order of a context carries no meaning.
\item \emph{Arity.}  The number of hypotheses is the number drawn.  A
  figure with two sequents above the line refuses an inference with one
  or with three.  Where the number depends on the formula (one hypothesis
  for each component of a tuple), the figure says so.
\item \emph{Assumption.}  ``$F \in \Gamma$'' means that some member of
  $\Gamma$ is $\aeq F$.  $F$ is a member of the context, never a
  subformula of one.  The operation is told which member by its argument;
  the checker is not, and searches $\Gamma$ for one that fits.
\item \emph{Matching.}  A term that occurs in a hypothesis but not in the
  conclusion (the $t$ of \texttt{forall-elim}, the witness of
  \texttt{forsome-intro}, the eigenvariables of \texttt{forall-intro}
  and \texttt{forsome-elim}) is not recorded.  The checker recovers it by
  one-way matching: the conclusion's formula, with a hole in place of the
  variable, is matched against the hypothesis' formula.  A hole is never
  filled by a term that contains a variable bound inside the pattern.
\item \emph{Weakening case.}  \texttt{forall-elim}, \texttt{forsome-elim},
  \texttt{theorem-assumption}, \texttt{if-true} and \texttt{if-false}
  also accept a hypothesis whose assertion is $G$ and whose context is
  included in $\Gamma$.  That is what the operation produces when the
  formula it adds is already in $\Gamma$, and the inference is then an
  instance of \texttt{weakening}.
\end{itemize}

\subsection{Logic, equality, induction and the structure eliminators}

\subsubsection*{Direct inference (command \texttt{di})}

The command records one of five operations, according to the form of $G$
(\texttt{pi-direct-inference!}).  On any other form it declines.
%% pi-direct-inference!, primitive-inferences.scm:75

\begin{center}
\begin{tabular}{@{}ll@{}}
\toprule
$G$ & operation \\
\midrule
$A \wedge B$ & \texttt{and-intro} \\
$A \Rightarrow B$ & \texttt{implies-intro} \\
$A \Leftrightarrow B$ & \texttt{iff-intro} \\
$\neg A$ & \texttt{not-intro} \\
$\forall x.\, B$ & \texttt{forall-intro} \\
\bottomrule
\end{tabular}
\end{center}

%% rule-checkers-logic.scm:347 and-intro, :386 implies-intro, :412 iff-intro, :399 not-intro
\VNBfig{and-intro}{\Gamma \Longrightarrow A \qquad \Gamma \Longrightarrow B}{\Gamma \Longrightarrow A \wedge B}
\VNBfig{implies-intro}{\Gamma, A \Longrightarrow B}{\Gamma \Longrightarrow A \Rightarrow B}
\VNBfig{iff-intro}{\Gamma, A \Longrightarrow B \qquad \Gamma, B \Longrightarrow A}{\Gamma \Longrightarrow A \Leftrightarrow B}
\VNBfig{not-intro}{\Gamma, A \Longrightarrow \textsc{falsity}}{\Gamma \Longrightarrow \neg A}

To prove $\forall x.\, B$ one proves $B$ for an $x$ about which nothing
is assumed.  In a sequent that means a variable that does not occur free
in any assumption: if $\Gamma$ said something about $x$, a proof of $B$
from $\Gamma$ would be a proof about that particular $x$, not about every
$x$.  A variable introduced under this condition is called an
\emph{eigenvariable}\index{eigenvariable}, and the condition is the \emph{eigenvariable
condition}.  \texttt{forall-intro} removes a chain of leading universal
quantifiers in one step, and collects on the way the guards $y \in A$
that restrict them.  Write the assertion as $G_0 = G$ and, for
$1 \leq i \leq k$, $G_{i-1} = \forall x_i.\, C_i$.
%% chk-forall-strip :441, chk-forall-eigen-ok? :491, chk-forall-intro-at :514,
%% forall-intro checker :555, all rule-checkers-logic.scm
\VNBfig{forall-intro}{\Gamma, g_1, \ldots, g_m \Longrightarrow G_k}{\Gamma \Longrightarrow G_0}
\begin{VNBcond}
\item[Chain] $1 \leq k$, and for each $i$ the formula $C_i[x_i{:=}y_i]$
  is either $(y_i \in A_i) \Rightarrow G_i$, in which case $y_i \in A_i$
  is \emph{collected}, or else it is $G_i$.  $g_1, \ldots, g_m$ are the
  collected formulas.
\item[Eigenvariable] Each $y_i$ is a variable, and
  $y_i \notin \FV(\Gamma)$, $y_i \notin \FV(G_{i-1})$ and
  $y_i \notin V_{i-1}$, where $V_0 = \emptyset$ and $V_i$ is
  $V_{i-1} \cup \FV(y_i \in A_i)$ when step $i$ collected a guard and
  $V_{i-1} \cup \{y_i\}$ otherwise.
\item[Matching] The $y_i$ are read off the hypothesis.  The checker
  tries $k$ equal to the whole length of the chain first and then each
  shorter length, since a shorter chain is also a sound inference.
\end{VNBcond}
In words: an eigenvariable is free neither in the context, nor in the
formula still to be generalised at that step, nor in a guard collected
before it (domain included), and the eigenvariables are pairwise
distinct.

The operation (\texttt{peel-foralls-raw}) always takes the whole chain.
It keeps $y_i = x_i$ when $x_i$ is free neither in $\Gamma$ nor in a
guard already collected, and otherwise takes a new variable, spelled as
$x_i$ followed by an underscore and a number.
%% peel-foralls-raw, primitive-inferences.scm:41
A guard is collected only when it has the form $y_i \in A_i$ for the
variable just bound.  Any other antecedent stays in the assertion:
$\forall x.\, P(x) \Rightarrow C$ yields the hypothesis
$\Gamma \Longrightarrow P(x) \Rightarrow C$, and a second direct inference
gives $\Gamma, P(x) \Longrightarrow C$.

\subsubsection*{\texttt{or-intro-left}, \texttt{or-intro-right}
  (commands \texttt{oi-l}, \texttt{oi-r})}

%% rule-checkers-logic.scm:362, :374
\VNBfig{or-intro-left}{\Gamma \Longrightarrow A}{\Gamma \Longrightarrow A \vee B}
\VNBfig{or-intro-right}{\Gamma \Longrightarrow B}{\Gamma \Longrightarrow A \vee B}

\subsubsection*{\texttt{forsome-intro}: supplying a witness (command \texttt{ew})}

%% rule-checkers-logic.scm:574; pi-exists-witness!, primitive-inferences.scm:357
\VNBfig{forsome-intro}{\Gamma \Longrightarrow B[x{:=}t]}{\Gamma \Longrightarrow \exists x.\, B}
\begin{VNBcond}
\item[Matching] $t$ is read off the hypothesis.
\item[Not checked] Well-formedness of $B[x{:=}t]$.  The operation
  checks it and raises an error, recording nothing, on a malformed
  witness.
\end{VNBcond}
The argument of the operation is $t$.  No definedness obligation is
posted for $t$, so the operation is not the mirror of
\texttt{forall-elim} in that respect.

\subsubsection*{\texttt{truth-intro}}

%% rule-checkers-logic.scm:588
\VNBfig{truth-intro}{\VNBnohyp}{\Gamma \Longrightarrow \textsc{truth}}

This operation is unreachable.  No proof command calls the procedure
that records it, and a \textsc{truth} assertion is already closed by
\texttt{assumption} and by \texttt{reflexivity}, both of which accept it.
It is listed because the code is present and a future command could
reach it.

\subsubsection*{Antecedent inference (command \texttt{ai})}

The argument is an assumption $F$ of $N$.  The command records one of
five operations according to the form of $F$
(\texttt{pi-antecedent-inference!}).
%% pi-antecedent-inference!, primitive-inferences.scm:132
\begin{center}
\begin{tabular}{@{}ll@{}}
\toprule
$F$ & operation \\
\midrule
$A \wedge B$ & \texttt{and-elim} \\
$A \vee B$ & \texttt{or-elim} \\
$\neg A$, with $A$ in $\Gamma$ & \texttt{not-elim} \\
$\exists x.\, B$ & \texttt{forsome-elim} \\
$A \Leftrightarrow B$ & \texttt{iff-elim} \\
\bottomrule
\end{tabular}
\end{center}
The command declines when $F$ is not an assumption of $N$, when $F$ is
$\neg A$ and $A$ is not in $\Gamma$, and when $F$ has any other form: an
implication, a universal formula or an atomic formula.  An implication in
the context is used by \texttt{detach} and \texttt{backchain}, a universal
formula by \texttt{forall-elim}.

%% rule-checkers-logic.scm:597 and-elim, :615 or-elim, :633 not-elim, :645 iff-elim
\VNBfig{and-elim}{(\Gamma - F), A, B \Longrightarrow G}{\Gamma \Longrightarrow G}
\begin{VNBcond}
\item[Assumption] $F = A \wedge B$, $F \in \Gamma$.
\end{VNBcond}
\VNBfig{or-elim}{(\Gamma - F), A \Longrightarrow G \qquad (\Gamma - F), B \Longrightarrow G}{\Gamma \Longrightarrow G}
\begin{VNBcond}
\item[Assumption] $F = A \vee B$, $F \in \Gamma$.
\end{VNBcond}
\VNBfig{not-elim}{\VNBnohyp}{\Gamma \Longrightarrow G}
\begin{VNBcond}
\item[Assumption] $\neg A \in \Gamma$ and $A \in \Gamma$.
\end{VNBcond}
\VNBfig{iff-elim}{(\Gamma - F), A \Rightarrow B, B \Rightarrow A \Longrightarrow G}{\Gamma \Longrightarrow G}
\begin{VNBcond}
\item[Assumption] $F = A \Leftrightarrow B$, $F \in \Gamma$.
\end{VNBcond}

%% forsome-elim checker, rule-checkers-logic.scm:665
\VNBfig{forsome-elim}{(\Gamma - F), B[x{:=}y] \Longrightarrow G}{\Gamma \Longrightarrow G}
\begin{VNBcond}
\item[Assumption] $F = \exists x.\, B$, $F \in \Gamma$.
\item[Eigenvariable] $y$ is a variable, and $y \notin \FV(G)$,
  $y \notin \FV(\Gamma - F)$, $y \notin \FV(F)$.  When $x$ is not free
  in $B$ there is no $y$ to test.
\item[Matching] $y$ is read off the hypothesis.
\end{VNBcond}
Nothing else is required: in particular nothing about the other
assumptions of $N$.  The operation names $y$ by the counter
(\texttt{fresh-var}): $x$ followed by an underscore and $n$, with $n$ the
next value for which that variable is free in none of $B$, $G$ and
$\Gamma - F$.  It never reuses $x$ itself.  The weakening case of the
conventions arises here only when $x$ is not free in $B$ and $B$ is
already in $\Gamma - F$.
%% pi-antecedent-inference! FORSOME arm, primitive-inferences.scm:169;
%% fresh-var, expressions.scm:701

\subsubsection*{\texttt{forall-elim}: instantiation of a universal assumption
  (commands \texttt{inst}, \texttt{inst+}; used by \texttt{fact})}

%% forall-elim checker, rule-checkers-logic.scm:702; pi-instantiate!, primitive-inferences.scm:282
\VNBfig{forall-elim}{\Gamma, B[x{:=}t] \Longrightarrow G \qquad \bigl(\Gamma \Longrightarrow t = t\bigr)}{\Gamma \Longrightarrow G}
\begin{VNBcond}
\item[Assumption] $\forall x.\, B \in \Gamma$.  It stays in the context.
\item[Matching] $t$ is read off the first hypothesis.
\item[Definedness] The second hypothesis is present exactly when
  $\Gamma$ does not certify $t$ (the certificate at the head of this
  section).
\item[Not checked] The checker accepts the second hypothesis present or
  absent.  When it is present, the checker requires it to be
  $\Gamma \Longrightarrow u = u$ for some term $u$, and does not require
  $u$ to be $t$.  Definedness of $t$ is therefore enforced by the
  operation alone.
\end{VNBcond}
The arguments are the universal assumption and $t$.  The second
hypothesis is there because a term may be undefined ($f(a)$ for $a$
outside the domain of $f$, a description with no referent), and a
universal formula speaks only of defined values.  $B[x{:=}t]$ is checked
for well-formedness; a malformed term raises an error and leaves the
graph unchanged.  The same operation is recorded, once per term, by the
procedure behind \texttt{apply-thm}, which cites a theorem and then
instantiates it at a list of terms (\texttt{pi-spec!}).

\subsubsection*{\texttt{assumption}: the goal is already in the context
  (command \texttt{ass})}

%% rule-checkers-logic.scm:738
\VNBfig{assumption}{\VNBnohyp}{\Gamma \Longrightarrow G}
\begin{VNBcond}
\item[Assumption] $G \in \Gamma$, or $G$ is \textsc{truth}.
\end{VNBcond}

\subsubsection*{\texttt{theorem-assumption}: citing a theorem or an axiom
  (command \texttt{ta}; used by \texttt{fact})}

%% rule-checkers-logic.scm:750; chk-installed-theorem-name :329
\VNBfig{theorem-assumption}{\Gamma, S \Longrightarrow G}{\Gamma \Longrightarrow G}
\begin{VNBcond}
\item[Installed] $S \aeq$ the formula installed under some name in the
  theorem table.
\end{VNBcond}
The argument of the operation is a name, and $S$ is the formula installed
under it; the operation declines when there is none.  The name is not
recorded with the inference, so the checker looks the statement up in
the table.

This is the only way a formula enters a proof without being proved in
it.  $S$ is an axiom (there are \VNBaxiomcount{} of them, listed in the
catalog), a definition, a proven theorem, or an asserted support.  The
bill printed when a theorem is installed
(Section~\ref{sec:proof-debt}) lists the asserted supports reached
through this operation, directly or through the theorems cited.

\subsubsection*{\texttt{cut}: introducing a lemma (commands \texttt{cut},
  \texttt{have!})}

%% rule-checkers-logic.scm:767; pi-cut!, primitive-inferences.scm:223
\VNBfig{cut}{\Gamma \Longrightarrow A \qquad \Gamma, A \Longrightarrow G}{\Gamma \Longrightarrow G}

The argument is $A$.  It is checked for well-formedness first, and a
malformed lemma raises an error before anything is recorded.

\begin{remark}
The operation never declines.  When $\Gamma$ already contains $A$, the
second hypothesis is $\Gamma \Longrightarrow G$, which is the sequent of
$N$ itself, so the inference has $N$ among its own hypotheses and
grounds nothing.  The same happens when $A$ is $G$.  A command that cuts
a formula already present therefore appears to succeed and leaves the
proof where it was.
\end{remark}

\subsubsection*{\texttt{weakening}: dropping assumptions (command \texttt{wk})}

%% rule-checkers-logic.scm:781; pi-weaken!, primitive-inferences.scm:208
\VNBfig{weakening}{\Delta \Longrightarrow G}{\Gamma \Longrightarrow G}
\begin{VNBcond}
\item[Inclusion] $\Delta \subseteq \Gamma$.
\end{VNBcond}
The operation takes one assumption $F$ of $N$ and posts
$\Delta = \Gamma - F$; it declines when $F$ is not an assumption of $N$.

\subsubsection*{\texttt{detach}: forward modus ponens (command \texttt{detach!})}

%% rule-checkers-logic.scm:790; pi-detach!, primitive-inferences.scm:442
\VNBfig{detach}{\Gamma, B \Longrightarrow G}{\Gamma \Longrightarrow G}
\begin{VNBcond}
\item[Assumption] $A \Rightarrow B \in \Gamma$ and $A \in \Gamma$.
\end{VNBcond}
The argument is the implication.  Both it and its antecedent stay.

\subsubsection*{\texttt{backchain}: reducing the goal to an antecedent
  (commands \texttt{bc}, \texttt{bc*})}

%% rule-checkers-logic.scm:807; pi-backchain!, primitive-inferences.scm:419
\VNBfig{backchain}{\Gamma \Longrightarrow A}{\Gamma \Longrightarrow G}
\begin{VNBcond}
\item[Assumption] $A \Rightarrow G \in \Gamma$.
\end{VNBcond}

\subsubsection*{\texttt{proof-by-contradiction} (command \texttt{pbc})}

%% rule-checkers-logic.scm:824
\VNBfig{proof-by-contradiction}{\Gamma, \neg G \Longrightarrow \textsc{falsity}}{\Gamma \Longrightarrow G}

There are no arguments, and the operation never declines.

\subsubsection*{\texttt{contraposition}}

%% rule-checkers-logic.scm:835
\VNBfig{contraposition}{\Gamma \Longrightarrow \neg B}{\Gamma \Longrightarrow \neg A}
\begin{VNBcond}
\item[Assumption] $A \Rightarrow B \in \Gamma$.
\end{VNBcond}
This operation is unreachable: no proof command calls the procedure that
records it.

\subsubsection*{\texttt{eq-subst}: rewriting the assertion by a context
  equation (command \texttt{subst})}

%% rule-checkers-logic.scm:857; chk-replaced? :241
\VNBfig{eq-subst}{\Gamma \Longrightarrow G'}{\Gamma \Longrightarrow G}
\begin{VNBcond}
\item[Assumption] $\Gamma$ contains $s = t$, $t = s$, $s \equiv t$ or
  $t \equiv s$.
\item[Replacement] $G' \not\aeq G$, and $G'$ is $G$ with some
  occurrences of $s$ replaced by $t$, or with some occurrences of $t$
  replaced by $s$.  A binder whose bound variable the replacement renamed
  is matched with its renamed form.
\item[Not checked] That no replaced occurrence lies in the scope of a
  binder whose variable is free in $s$ or in $t$.  The checker compares
  $G$ and $G'$ position by position and does not test it.
\end{VNBcond}
The argument is the equation $s = t$ (or $s \equiv t$).  The operation
replaces \emph{every} occurrence of $s$ by $t$, in operator position as
well as in argument position, and declines when nothing changes.  It
leaves alone a subterm below a binder whose bound variable occurs free in
$s$ or in $t$, since rewriting there would capture; when $s$ is a
variable it replaces the free occurrences only.  $G'$ is checked for
well-formedness.
%% pi-eq-subst! :563, replace-term :496, primitive-inferences.scm

\begin{remark}
The direction of the rewrite is fixed by the argument, not by the
orientation in which $\Gamma$ happens to hold the equation.  To use a
context equation from right to left, name it the other way round.
\end{remark}

\subsubsection*{\texttt{reflexivity} and \texttt{quasi-reflexivity}
  (commands \texttt{rfl}, \texttt{qrfl})}

%% rule-checkers-logic.scm:879, :890
\VNBfig{reflexivity}{\VNBnohyp}{\Gamma \Longrightarrow t = t}
\begin{VNBcond}
\item[Definedness] $\Gamma$ certifies $t$.  The operation also accepts
  the assertion \textsc{truth}.
\item[Not checked] The definedness condition.  The checker accepts
  $\Gamma \Longrightarrow u = u$ for any term $u$.
\end{VNBcond}
\VNBfig{quasi-reflexivity}{\VNBnohyp}{\Gamma \Longrightarrow t \equiv t}

\texttt{quasi-reflexivity} asks nothing further: quasi-equality holds
between two undefined terms as well.

\subsubsection*{\texttt{if-true}, \texttt{if-false}
  (commands \texttt{if-true}, \texttt{if-false})}

%% chk-if-rule, rule-checkers-logic.scm:900-933
\VNBfig{if-true}{\Gamma \Longrightarrow p \qquad \Gamma, \mathrm{IF}(p,a,b) = a \Longrightarrow G}{\Gamma \Longrightarrow G}
\VNBfig{if-false}{\Gamma \Longrightarrow \neg p \qquad \Gamma, \mathrm{IF}(p,a,b) = b \Longrightarrow G}{\Gamma \Longrightarrow G}

The argument is the conditional term $\mathrm{IF}(p, a, b)$.  It need not
occur in $G$.  The assertion is not rewritten: the branch equation
arrives as an assumption, and it takes an \texttt{eq-subst} to use it.
An argument that is not a conditional term raises an error rather than
declining.

\subsubsection*{\texttt{cartesian-intro}, \texttt{cartesian-elim}
  (commands \texttt{ci}, \texttt{ce})}

%% rule-checkers-logic.scm:939, :963
\VNBfig{cartesian-intro}{\Gamma \Longrightarrow a_1 \in A_1 \quad \cdots \quad \Gamma \Longrightarrow a_n \in A_n}{\Gamma \Longrightarrow [a_1, \ldots, a_n] \in A_1 \times \cdots \times A_n}
\begin{VNBcond}
\item[Arity] The literal tuple and the literal product have the same
  length $n$, and there is one hypothesis per component, in order.
\end{VNBcond}
\VNBfig{cartesian-elim $k$}{\Gamma, a_k \in A_k \Longrightarrow G}{\Gamma \Longrightarrow G}
\begin{VNBcond}
\item[Assumption] $[a_1, \ldots, a_n] \in A_1 \times \cdots \times A_n
  \in \Gamma$, with tuple and product of the same length.  It stays.
\item[Index] $k$ is recorded with the name, and $1 \leq k \leq n$.
\end{VNBcond}

\subsubsection*{\texttt{tuples-intro}, \texttt{tuples-elim}
  (commands \texttt{ti}, \texttt{te})}

$\mathrm{TUPLES}(A)$ is the class of finite tuples over $A$.
%% rule-checkers-logic.scm:991, :1012
\VNBfig{tuples-intro}{\Gamma \Longrightarrow a_1 \in A \quad \cdots \quad \Gamma \Longrightarrow a_n \in A}{\Gamma \Longrightarrow [a_1, \ldots, a_n] \in \mathrm{TUPLES}(A)}
\begin{VNBcond}
\item[Arity] One hypothesis per component; none for the empty tuple.
\end{VNBcond}
\VNBfig{tuples-elim $k$}{\Gamma, a_k \in A \Longrightarrow G}{\Gamma \Longrightarrow G}
\begin{VNBcond}
\item[Assumption] $[a_1, \ldots, a_n] \in \mathrm{TUPLES}(A) \in \Gamma$.
\item[Index] $k$ is recorded with the name, and $1 \leq k \leq n$.
\end{VNBcond}

\subsubsection*{\texttt{union-intro}, \texttt{union-elim}
  (commands \texttt{ui}, \texttt{ue})}

$\mathrm{UNION}$ takes $n$ arguments.
%% rule-checkers-logic.scm:1038, :1060
\VNBfig{union-intro $k$}{\Gamma \Longrightarrow x \in A_k}{\Gamma \Longrightarrow x \in A_1 \cup \cdots \cup A_n}
\begin{VNBcond}
\item[Index] $k$ is recorded with the name, and $1 \leq k \leq n$.
\end{VNBcond}
\VNBfig{union-elim}{(\Gamma - F), x \in A_1 \Longrightarrow G \quad \cdots \quad (\Gamma - F), x \in A_n \Longrightarrow G}{\Gamma \Longrightarrow G}
\begin{VNBcond}
\item[Assumption] $F = x \in A_1 \cup \cdots \cup A_n$, $F \in \Gamma$.
  It is consumed.
\item[Arity] One hypothesis per set, in order.
\end{VNBcond}

\subsubsection*{\texttt{intersection-intro}, \texttt{intersection-elim}
  (commands \texttt{ii}, \texttt{ie})}

%% rule-checkers-logic.scm:1079, :1098
\VNBfig{intersection-intro}{\Gamma \Longrightarrow x \in A_1 \quad \cdots \quad \Gamma \Longrightarrow x \in A_n}{\Gamma \Longrightarrow x \in A_1 \cap \cdots \cap A_n}
\VNBfig{intersection-elim $k$}{\Gamma, x \in A_k \Longrightarrow G}{\Gamma \Longrightarrow G}
\begin{VNBcond}
\item[Assumption] $x \in A_1 \cap \cdots \cap A_n \in \Gamma$.  It stays.
\item[Index] $k$ is recorded with the name, and $1 \leq k \leq n$.
\end{VNBcond}

\subsubsection*{\texttt{nth-reduce}, \texttt{length-reduce},
  \texttt{functoid-beta} (commands \texttt{nth-r}, \texttt{len-r},
  \texttt{beta})}

The three have the same figure and differ in the contraction.
%% chk-contracts? rule-checkers-logic.scm:284; checkers :1137, :1153, :1174
\VNBfig{}{\Gamma \Longrightarrow G'}{\Gamma \Longrightarrow G}
\begin{VNBcond}
\item[Contraction] $G' \not\aeq G$, and $G'$ is obtained from $G$ by
  contracting redexes, where a redex that a contraction creates may be
  contracted in its turn.  The redexes are, respectively:
  $\mathrm{NTH}(k, [t_1, \ldots, t_n])$ with $1 \leq k \leq n$, which
  contracts to $t_k$; $\mathrm{LENGTH}([t_1, \ldots, t_n])$, which
  contracts to $n$; and a functoid applied to as many arguments as it
  binds, which contracts to its body with the arguments substituted in
  parallel for the bound variables.
\end{VNBcond}
None takes an argument.  Each operation contracts every redex of its
kind in $G$ and declines when $G'$ is $G$.  \texttt{nth-reduce} raises an
error on an index outside $1 \leq k \leq n$.  Nothing is asked of the
entries: a literal tuple has $n$ places whether or not the terms in them
denote.

\subsubsection*{\texttt{nn-induction} (command \texttt{ni})}

%% rule-checkers-logic.scm:1191
\VNBfig{nn-induction}{\Gamma \Longrightarrow B[n{:=}0] \qquad \Gamma \Longrightarrow \forall n.\, n \in \mathbb{N} \Rightarrow (B \Rightarrow B[n{:=}\mathrm{succ}\,n])}{\Gamma \Longrightarrow \forall n.\, n \in \mathbb{N} \Rightarrow B}
\begin{VNBcond}
\item[Shape] The guard is $n \in \mathbb{N}$ on the variable of the
  outermost quantifier.
\end{VNBcond}
The operation has no argument and declines on any other shape.  In
particular the induction variable must be the outermost one: a statement
that quantifies over something else first is not of this shape.

\subsubsection*{\texttt{transfinite-induction} (command \texttt{tfi})}

%% rule-checkers-logic.scm:1247; chk-tfi-beta-ok? :1241
\VNBfig{transfinite-induction}{\Gamma \Longrightarrow \forall \alpha.\, \bigl(\alpha \in \mathrm{ORD} \wedge \forall \beta.\, \beta <_{\mathrm{ORD}} \alpha \Rightarrow P[\alpha{:=}\beta]\bigr) \Rightarrow P}{\Gamma \Longrightarrow \forall \alpha.\, \alpha \in \mathrm{ORD} \Rightarrow P}
\begin{VNBcond}
\item[Freshness] $\beta$ is a variable, $\beta \neq \alpha$,
  $\beta \notin \FV(P)$, $\beta \notin \FV(\Gamma)$.
\item[Matching] $\beta$ is read off the hypothesis.
\end{VNBcond}
This is strong induction: the hypothesis of the step is $P$ at every
smaller ordinal.

\begin{remark}
The operation \texttt{transfinite-induction} and the axiom of the same
name in \texttt{structure-library/ordinals.scm} are two different
objects.  One is code in the kernel, the other a formula in the catalog.
\end{remark}

\subsubsection*{\texttt{transfinite-induction-3cases} (command \texttt{tfi3})}

%% rule-checkers-logic.scm:1274
\[
\frac{\displaystyle
  \begin{array}{l}
  \Gamma \Longrightarrow P[\alpha{:=}0] \\[2pt]
  \Gamma \Longrightarrow \forall \alpha.\,
    (\alpha \in \mathrm{ORD} \wedge P) \Rightarrow P[\alpha{:=}\mathrm{succ}_{\mathrm{ORD}}\,\alpha] \\[2pt]
  \Gamma \Longrightarrow \forall \alpha.\,
    \bigl(\mathrm{LIMIT\text{-}ORD}(\alpha) \wedge
          \forall \beta.\, \beta <_{\mathrm{ORD}} \alpha \Rightarrow P[\alpha{:=}\beta]\bigr)
    \Rightarrow P
  \end{array}}
 {\displaystyle \Gamma \Longrightarrow \forall \alpha.\, \alpha \in \mathrm{ORD} \Rightarrow P}
\]
\begin{VNBcond}
\item[Freshness] As for \texttt{transfinite-induction}, with $\beta$ read
  off the third hypothesis.
\end{VNBcond}
The three hypotheses are the base, successor and limit cases, in that
order.

\subsection{Set formation, descriptions and functions}

The three binders of this group carry a formula or a class expression in
their body, with a variable of the binder free in it.  A first-order
axiom would need a variable ranging over formulas, which VNB does not
have, so each is characterised by operations instead.
$\mathrm{SEP}(x, A, p)$ is $\{x \in A \mid p\}$,
$\mathrm{COMP}(x, p)$ is $\{x \mid p\}$, and
$\mathrm{BIG\text{-}UNION}(z, A, b)$ is $\bigcup_{z \in A} b$.

\subsubsection*{\texttt{sep-sethood}, \texttt{sep-mem-intro},
  \texttt{sep-mem-elim}
  (commands \texttt{sep-set}, \texttt{sep-mi}, \texttt{sep-me})}

%% rcs-sep-sethood :107, rcs-sep-mem-intro :160, rcs-sep-mem-elim :167, rule-checkers-schema.scm
\VNBfig{sep-sethood}{\Gamma \Longrightarrow A \in \mathrm{SET}}{\Gamma \Longrightarrow \mathrm{SEP}(x,A,p) \in \mathrm{SET}}
\VNBfig{sep-mem-intro}{\Gamma \Longrightarrow y \in A \qquad \Gamma \Longrightarrow p[x{:=}y]}{\Gamma \Longrightarrow y \in \mathrm{SEP}(x,A,p)}
\VNBfig{sep-mem-elim}{(\Gamma - F), y \in A, p[x{:=}y] \Longrightarrow G}{\Gamma \Longrightarrow G}
\begin{VNBcond}
\item[Assumption] $F = y \in \mathrm{SEP}(x,A,p)$, $F \in \Gamma$.  It
  is consumed.
\end{VNBcond}

\subsubsection*{\texttt{comp-mem-intro}, \texttt{comp-mem-elim}
  (commands \texttt{comp-mi}, \texttt{comp-me})}

A comprehension is in general a proper class, and its members are sets,
so the domain condition of separation is replaced by membership in
$\mathrm{SET}$.
%% rcs-comp-mem-intro :177, rcs-comp-mem-elim :182, rule-checkers-schema.scm
\VNBfig{comp-mem-intro}{\Gamma \Longrightarrow y \in \mathrm{SET} \qquad \Gamma \Longrightarrow p[x{:=}y]}{\Gamma \Longrightarrow y \in \mathrm{COMP}(x,p)}
\VNBfig{comp-mem-elim}{(\Gamma - F), y \in \mathrm{SET}, p[x{:=}y] \Longrightarrow G}{\Gamma \Longrightarrow G}
\begin{VNBcond}
\item[Assumption] $F = y \in \mathrm{COMP}(x,p)$, $F \in \Gamma$.  It is
  consumed.
\end{VNBcond}
There is no sethood operation for comprehension.  A comprehension is a
set only when it is separately shown to be bounded by one.

\subsubsection*{\texttt{big-union-sethood}, \texttt{big-union-mem-intro},
  \texttt{big-union-mem-elim}
  (commands \texttt{bu-set}, \texttt{bu-mi}, \texttt{bu-me})}

Write $U$ for $\mathrm{BIG\text{-}UNION}(z, A, b)$.
%% rcs-big-union-sethood :194, rcs-big-union-mem-intro :234,
%% rcs-big-union-mem-elim :260, rule-checkers-schema.scm
\VNBfigS{big-union-sethood}{\Gamma \Longrightarrow A \in \mathrm{SET} \\ \Gamma \Longrightarrow \forall v.\, v \in A \Rightarrow b[z{:=}v] \in \mathrm{SET}}{\Gamma \Longrightarrow U \in \mathrm{SET}}
\begin{VNBcond}
\item[Freshness] $v = z$, or $v \notin \FV(A)$.
\item[Not checked] That $v \notin \FV(b)$ when $v \neq z$.  Without it
  the quantifier of the second hypothesis captures a free $v$ of $b$.
  The operation takes $v$ new, free in none of $b$, $A$, $G$ and
  $\Gamma$ (\texttt{pi-big-union-sethood!}, primitive-inferences.scm).
\end{VNBcond}
%% pi-big-union-sethood!, primitive-inferences.scm:1651
\VNBfig{big-union-mem-intro}{\Gamma \Longrightarrow w \in A \qquad \Gamma \Longrightarrow x \in b[z{:=}w]}{\Gamma \Longrightarrow x \in U}
\begin{VNBcond}
\item[Matching] $w$ is read off the first hypothesis.
\end{VNBcond}
\VNBfig{big-union-mem-elim}{(\Gamma - F), e \in A, x \in b[z{:=}e] \Longrightarrow G}{\Gamma \Longrightarrow G}
\begin{VNBcond}
\item[Assumption] $F = x \in U$, $F \in \Gamma$.  It is consumed.
\item[Eigenvariable] $e$ is a variable free in none of $G$, $A$, $b$,
  $x$ and $\Gamma - F$.
\end{VNBcond}
\texttt{big-union-mem-intro} takes the witness $w$ as its argument, and
checks $x \in b[z{:=}w]$ for well-formedness.

\subsubsection*{\texttt{iota-def}, \texttt{iota-in-elim}
  (commands \texttt{iota-d}, \texttt{iota-e})}

$\mathrm{IOTA}(x, p)$ is the unique $x$ satisfying $p$, and denotes
nothing when there is no such $x$ or more than one.  Write $\iota$ for
$\mathrm{IOTA}(x, p)$.  Both operations take the term as their argument;
it need not occur in $G$.  Both grant the defining property
$p[x{:=}\iota]$ as an assumption, and they differ in what pays for it.
%% rcs-iota-def :308, rcs-iota-in-elim :372, rcs--context-defines? :350, rule-checkers-schema.scm
\VNBfigS{iota-def}{\Gamma \Longrightarrow \exists x.\, \bigl(p \wedge \forall y.\, p[x{:=}y] \Rightarrow x = y\bigr) \\ \Gamma, p[x{:=}\iota] \Longrightarrow G}{\Gamma \Longrightarrow G}
\begin{VNBcond}
\item[Freshness] $y$ is a variable, $y \neq x$, $y \notin \FV(p)$.
\end{VNBcond}
\VNBfig{iota-in-elim}{\Gamma, p[x{:=}\iota] \Longrightarrow G}{\Gamma \Longrightarrow G}
\begin{VNBcond}
\item[Denotation] $\Gamma$ contains $\iota \in X$ for some $X$, or an
  equation $\iota = u$ or $u = \iota$.
\end{VNBcond}
\texttt{iota-def} pays with a proof of existence and uniqueness.
\texttt{iota-in-elim} pays with the context: the two forms of the
Denotation condition are the first two clauses of the definedness
certificate, and the operation declines without one of them.

\subsubsection*{\texttt{lambda-type} (command \texttt{lam-t})}

$\mathrm{VNB\text{-}LAMBDA}(\sigma, A, b)$ is the set of pairs
$\{(x, b) : x \in A\}$, with $\sigma$ the binder, a variable or a list of
variables.  The term carries its domain.
%% rcs-lambda-type :437, rcs--lambda-typing-ok? :414, rule-checkers-schema.scm
\VNBfig{lambda-type}{\Gamma \Longrightarrow \forall v.\, v \in A \Rightarrow b[x{:=}v] \in B \qquad \Gamma \Longrightarrow A \in \mathrm{SET}}{\Gamma \Longrightarrow \mathrm{VNB\text{-}LAMBDA}(x, A, b) \in \mathrm{FUN}(A, B)}
\begin{VNBcond}
\item[Domain] The domain the term declares is $\aeq$ the domain of the
  $\mathrm{FUN}$.
\item[Freshness] $v = x$, or $v \notin \FV(b) \cup \FV(B)$.
\item[Binder list] For $\sigma = [x_1, \ldots, x_n]$ the domain is
  $A_1 \times \cdots \times A_n$ with the same $n$, and the first
  hypothesis is the $n$-fold guarded universal
  $\forall v_1.\, v_1 \in A_1 \Rightarrow \cdots \Rightarrow
   \forall v_n.\, v_n \in A_n \Rightarrow b[x_1{:=}v_1, \ldots, x_n{:=}v_n] \in B$,
  each $v_i$ fresh as above and the $v_i$ pairwise distinct.
\end{VNBcond}
The operation has no argument, and declines unless the Domain and Binder
list conditions hold.  Neither condition is decoration.  Without the
first, one term types into $\mathrm{FUN}(A, B)$ for every $A$.  Without
the sethood hypothesis, a function on a proper class, which is itself a
proper class, would be certified into a $\mathrm{FUN}$.

\subsubsection*{\texttt{lambda-beta}, \texttt{lambda-beta-hyp}
  (commands \texttt{lam-b}, \texttt{lam-b-h})}

A \emph{redex} is a term
\[
  \mathrm{VNB\text{-}LAMBDA}(\sigma, A, b)(u_1, \ldots, u_n),
\]
with $n = 1$ for a single binder and $n$ the length of $\sigma$ for a
binder list.  It contracts to $b$ with the $u_i$ substituted in parallel
for the variables of $\sigma$.  Write $o_1, \ldots, o_m$ for the
\emph{obligations}, $m \geq 0$.
%% rcs-lambda-beta :854, rcs-lambda-beta-hyp :881, rcs--beta-match/fuel :701,
%% rcs--beta-licences-ok? :816, rcs--beta-licensed? :775, rule-checkers-schema.scm
\VNBfigS{lambda-beta}{\Gamma \Longrightarrow G' \\ \Gamma \Longrightarrow o_1 \quad \cdots \quad \Gamma \Longrightarrow o_m}{\Gamma \Longrightarrow G}
\VNBfigS{lambda-beta-hyp}{(\Gamma - F), F' \Longrightarrow G \\ \Gamma \Longrightarrow o_1 \quad \cdots \quad \Gamma \Longrightarrow o_m}{\Gamma \Longrightarrow G}
\begin{VNBcond}
\item[Contraction] $G'$ (respectively $F'$, for some $F \in \Gamma$) is
  reached from $G$ (respectively $F$) by a sequence of contractions,
  each taken at a position of the formula reached so far.  A redex that
  a contraction creates may be contracted in its turn.
\item[Licence] Each redex contracted at a position $\pi$ is licensed
  there or owed.  It is \emph{licensed} when $u \in A$ (for $n = 1$) or
  $[u_1, \ldots, u_n] \in A$, or, when $A$ is a product
  $A_1 \times \cdots \times A_n$, every $u_i \in A_i$, is among the
  formulas known at $\pi$: the assumptions of $\Gamma$ in which no
  variable bound between the root and $\pi$ occurs free, together with
  the guards of the guarded universals that enclose $\pi$ and the domain
  memberships of the $\mathrm{VNB\text{-}LAMBDA}$, $\mathrm{SEP}$ and
  $\mathrm{BIG\text{-}UNION}$ binders that enclose it.  It is
  \emph{owed} when that membership is one of the $o_j$ and contains no
  variable bound between the root and $\pi$.  The argument may be taken
  as written or with the redexes inside it contracted; in the second
  case those inner redexes count as contracted and need their own
  licence.
\item[Obligations] Every $o_j$ is owed by some redex.
\item[Bound] The checker's search takes at most 400 contractions, and
  refuses when that runs out.
\end{VNBcond}
\texttt{lambda-beta} takes no argument; \texttt{lambda-beta-hyp} takes
the assumption $F$.  Each reduces every redex it can license or owe, in
one bottom-up pass, and declines when the reduction changes nothing.
The obligations are posted after the reduced formula, one per redex
whose licence is not evident.  A member of $\mathrm{FUN}(A, B)$ is
defined exactly on $A$, so reducing off the domain without the
obligation would make a term denote that the library says does not.

\begin{remark}
The obligation is posted in the context of $N$.  A redex under a binder
that has not yet been introduced therefore owes a membership in which the
bound variable is free, and that leaf cannot be closed.  Introducing the
binder and typing the argument first is not a preference but the
condition under which the obligation is provable.
\end{remark}

\subsection{Rewriting}

\subsubsection*{\texttt{macete}: rewriting the assertion by a theorem
  (commands \texttt{mac}, \texttt{macm})}

A \emph{macete} is a rewrite rule read off an installed theorem of the
form
\[
  \forall x_1 \ldots x_n.\; C_1 \Rightarrow (\cdots (C_k \Rightarrow L = R)),
\]
where $=$ may also be quasi-equality or $\Leftrightarrow$, and the
conditions $C_i$ may be absent.  A name may also denote a rule that is
not a theorem: a functoid $F$ with parameters $\vec p$ and body $e$
denotes $F(\vec p) \to e$, a structure accessor with slot $k$ denotes
$\mathrm{ACC}(s) \to \mathrm{NTH}(k, s)$, and a functor projection
\texttt{F@ACC} denotes the corresponding component of $F$'s body.
%% rkw--macete-definition/compute, rule-checkers-rewrite.scm:593
The \emph{local context} of a position in a formula is $\Gamma$, with
conjunctions split into their conjuncts, extended on the way down:
entering $B$ in $A \Rightarrow B$ adds $A$; entering either conjunct of a
conjunction adds the other; entering either disjunct of $A \vee B$ adds
the negation of the other; entering the scope of a binder
($\forall$, $\exists$, $\mathrm{IOTA}$, $\mathrm{COMP}$, the body of
$\mathrm{SEP}$, $\mathrm{BIG\text{-}UNION}$ or $\mathrm{VNB\text{-}LAMBDA}$,
a functoid) removes every formula in which the bound variable occurs
free (Monk~\cite{Monk1988}).
%% rkw--descend, rule-checkers-rewrite.scm:471
%% rkw--check-macete, rule-checkers-rewrite.scm:687; rkw--condition-discharged? :390
\VNBfig{macete $N$ $L$ $R$}{\Gamma \Longrightarrow G' \qquad \Gamma \Longrightarrow c_1 \quad \cdots \quad \Gamma \Longrightarrow c_m}{\Gamma \Longrightarrow G}
\begin{VNBcond}
\item[Licence] $N$ names an installed theorem, a functoid, an accessor
  or a functor projection, and denotes the rule
  $(C_1, \ldots, C_k;\ L_0 \to R_0)$.
\item[Tag] The recorded $L$ and $R$ are $\sigma L_0$ and $\sigma R_0$
  for one substitution $\sigma$.  (Not tested for the two variadic rules,
  \texttt{union-decompose} and \texttt{intersection-decompose}.)
\item[Rewrite] $G'$ is $G$ with some subterms of the form $\tau L_0$
  replaced by $\tau R_0$, $\tau$ a substitution of the rule's variables
  found by matching.  A pattern $\mathrm{succ}(P)$ matches a positive
  integer literal $m$ by matching $P$ against $m-1$.
\item[Conditions] At each rewritten position every $\tau C_i$ is
  \textsc{truth}, or is in the local context of the position, or is a
  closed arithmetic sentence that evaluates to true, or is one of the
  $c_j$.
\end{VNBcond}
The operation replaces each subterm that matches $L_0$ and whose
conditions hold in the local context; a match whose conditions do not
all hold is left alone, and the rewriter descends into its subterms
instead.  It declines when $G'$ is $G$.  Its local context is the one
above, except that inside a conjunction $A \wedge B$ it extends only the
context of $B$, by $A$.  Hypotheses $c_j$ appear only when the command
posts the conditions it could not discharge as subgoals.  The name of
the theorem and the two patterns are recorded with the operation, so
that the inference can be checked against the rule it claims to apply
(Section~\ref{sec:inference-checked}).

\subsubsection*{\texttt{macete-hyp}: rewriting an assumption
  (commands \texttt{mac-h}, \texttt{mac-h*})}

%% rkw--check-macete-hyp, rule-checkers-rewrite.scm:741
\VNBfig{macete-hyp $N$ $L$ $R$}{(\Gamma - H), H' \Longrightarrow G \qquad \Gamma \Longrightarrow c_1 \quad \cdots \quad \Gamma \Longrightarrow c_m}{\Gamma \Longrightarrow G}
\begin{VNBcond}
\item[Assumption] $H \in \Gamma$.  Exactly one formula of $\Gamma$ is
  missing from the first hypothesis, and at most one formula is new in
  it (none when $H'$ was already in $\Gamma$).
\item[Licence, Tag, Rewrite, Conditions] As for \texttt{macete}, with
  $H$ and $H'$ in place of $G$ and $G'$, and the local context started
  from all of $\Gamma$, $H$ included.
\item[Not checked] The two conditions below on which the operation
  declines.  Neither is needed for soundness: every rule a name denotes
  replaces a term by an equal or quasi-equal one, or a formula by an
  equivalent one, under its conditions.
\end{VNBcond}
The arguments are the name and $H$.  The rewrite is made whether or not
the conditions hold, and one hypothesis $\Gamma \Longrightarrow c_j$ is
posted for each instantiated condition that the local context does not
contain.  Besides declining when the rewrite changes nothing, the
operation declines on two conditions read off the named theorem.  The
theorem's core must be an equivalence, that is an $\Leftrightarrow$ or
one of the two equalities, since replacing an assumption by a
consequence of it would lose information.  And no variable of $R$ or of
a condition may be left undetermined by the match, since such a variable
would be conjured free into the context.

\subsubsection*{\texttt{cartesian-decompose} (command \texttt{mac})}

%% rkw--check-cartesian-decompose :870, rkw--cartesian-shape? :819, rule-checkers-rewrite.scm
\VNBfig{cartesian-decompose}{\Gamma \Longrightarrow G'}{\Gamma \Longrightarrow G}
\begin{VNBcond}
\item[Rewrite] $G'$ is $G$ with some subformulas
  $x \in A_1 \times \cdots \times A_n$, $n \geq 1$, replaced by
  \[
    \exists v_1.\, v_1 \in A_1 \wedge
    \bigl(\cdots \wedge
    \exists v_n.\, v_n \in A_n \wedge x = [v_1, \ldots, v_n] \bigr).
  \]
\item[Freshness] The $v_i$ are pairwise distinct, and none is free in
  $x$ or in any $A_j$.
\end{VNBcond}
No argument beyond the name.  The operation rewrites every such
subformula, with new $v_i$, and declines when $G'$ is $G$.  This is the
membership condition of the product, and it is not a formula in the
catalog: the fresh witnesses cannot be written as a fixed template, so it
is code.  It is on the primitive shelf, and contributes nothing to a
bill.

\subsubsection*{\texttt{tuple-equality-decompose} (command \texttt{mac})}

%% rkw--check-tuple-equality-decompose :902, rkw--tuple-conjunction :883, rule-checkers-rewrite.scm
\VNBfig{tuple-equality-decompose}{\Gamma \Longrightarrow G'}{\Gamma \Longrightarrow G}
\begin{VNBcond}
\item[Rewrite] $G'$ is $G$ with some subformulas
  $[a_1, \ldots, a_n] = [b_1, \ldots, b_n]$, literal tuples of the same
  length, replaced by $a_1 = b_1 \wedge (a_2 = b_2 \wedge \cdots)$; by
  $a_1 = b_1$ when $n = 1$, and by \textsc{truth} when $n = 0$.
\end{VNBcond}
No argument beyond the name.  The operation rewrites every such
subformula and declines when $G'$ is $G$.  Tuples of different lengths
are not rewritten.  That two of them cannot be equal is a separate fact
and not this operation's business.

\subsection{Oracles}

An oracle decides a class of assertions by computation.  Six of the
seven, on success, record an inference with no hypotheses, so $N$ is
grounded at once; \texttt{arith-simplify} posts one hypothesis.
\texttt{ineq} and \texttt{sos} record, with the name, the certificate the
computation produced, and print it; the others record the name alone.
Each inference is checked before it is written, as described in
Section~\ref{sec:inference-checked}.  The figures below state what the
checker of each oracle tests.
%% ineq tag: pi-ineq!, structure-library/ineq-oracle.scm:308;
%% sos tag: structure-library/sos-oracle.scm:120; arith-eval.scm:382-388

\subsubsection*{\texttt{arith-ground}, \texttt{arith-forsome},
  \texttt{arith-simplify} (command \texttt{arith})}

One procedure records these three, trying them in this order.
%% rco--check-arith-ground :242, rco--check-arith-forsome :265,
%% rco--check-arith-simplify :298, rule-checkers-oracle.scm
\VNBfig{arith-ground}{\VNBnohyp}{\Gamma \Longrightarrow G}
\begin{VNBcond}
\item[Evaluation] $G$ is closed, and exact evaluation gives true.
\end{VNBcond}
\VNBfig{arith-forsome}{\VNBnohyp}{\Gamma \Longrightarrow \exists x.\, x \in C \wedge e = x}
\begin{VNBcond}
\item[Evaluation] $e$ evaluates exactly to a number $v$, and the body
  evaluates to true at $x = v$.  The equation may also be written
  $x = e$.
\end{VNBcond}
\VNBfig{arith-simplify}{\Gamma \Longrightarrow G'}{\Gamma \Longrightarrow G}
\begin{VNBcond}
\item[Evaluation] $G'$ is $G$ with some ground subterms replaced by
  the numbers to which they evaluate exactly.
\end{VNBcond}

\texttt{arith-ground} fires when $G$ evaluates to true.  The evaluator
reads $=$ and $\leq$ between terms it can reduce to exact numbers,
membership in $\mathbb{N}$, $\mathbb{Z}$, $\mathbb{Q}$, $\mathbb{R}$ and
$\mathbb{C}$, and the connectives $\neg$, $\wedge$, $\vee$,
$\Rightarrow$.  Everything else is undecided: a free variable, a
quantifier, a biconditional, a strict $<$, or a value that is not exact.
Arithmetic is exact throughout; a computed floating-point value is
refused rather than rounded.

\texttt{arith-forsome} fires on the shape of its figure when $e$ reduces
to a number and that number lies in $C$.

\texttt{arith-simplify} is the fallback.  When replacing the ground
numeric subterms of $G$ by their values changes $G$, it posts the single
hypothesis $\Gamma \Longrightarrow G'$.

The command declines when none of the three applies.

\subsubsection*{\texttt{ring-simplify}: identities of rings
  (commands \texttt{rs}, \texttt{simp})}

%% rco--check-rs :939, rco--crs-number-ok? :881, rco--num-gens-ok? :846, rule-checkers-oracle.scm
\VNBfig{ring-simplify}{\VNBnohyp}{\Gamma \Longrightarrow \forall x_1.\, x_1 \in D_1 \Rightarrow \cdots \forall x_k.\, x_k \in D_k \Rightarrow e_1 = e_2}
\begin{VNBcond}
\item[Normal form] $e_1$ and $e_2$ have the same normal form in the
  free associative ring $\mathbb{Z}[\textit{generators}]$.
\item[Typing] Each $D_i$ is one of $\mathbb{N}$, $\mathbb{Z}$,
  $\mathbb{Q}$, $\mathbb{R}$, $\mathbb{C}$, and every generator of
  $e_1$ or $e_2$ is one of the $x_i$ or is typed in one of these domains
  by an assumption $g \in D$ of $\Gamma$.
\end{VNBcond}
No arguments.  A \emph{generator} is a maximal subterm that ground
evaluation cannot reduce to a number, and multiplication does not
commute.  The operation declines unless both conditions hold.  The
generators counted are those of the two sides as written, including any
that cancel.  Otherwise $x - x = 0$ would close for an $x$ of which
nothing is known, and in VNB a membership is also what says that a term
denotes.

\subsubsection*{\texttt{comm-ring-simplify}: identities of commutative rings
  (command \texttt{crs})}

%% rco--check-crs :929, rco--crs-structure-ok? :894, rule-checkers-oracle.scm
\VNBfig{comm-ring-simplify}{\VNBnohyp}{\Gamma \Longrightarrow \forall x_1.\, x_1 \in D_1 \Rightarrow \cdots \forall x_k.\, x_k \in D_k \Rightarrow e_1 = e_2}
\begin{VNBcond}
\item[Normal form] $e_1$ and $e_2$ have the same normal form in the
  free commutative ring $\mathbb{Z}[\textit{generators}]$, literal
  powers being expanded first.
\item[Typing] Either as for \texttt{ring-simplify}; or the operations
  are those of a structure $R$, $(\mathrm{ADD}\ R)$, $(\mathrm{MUL}\ R)$,
  $(\mathrm{NEG}\ R)$, $(\mathrm{ZERO}\ R)$, $(\mathrm{ONE}\ R)$,
  $\mathrm{IS\text{-}COMMUTATIVE\text{-}RING}(R)$ is a peeled guard or is
  in $\Gamma$, and every generator is a variable bound by a guard
  $x_i \in \mathrm{CARR}(R)$ or is typed in $\mathrm{CARR}(R)$ by an
  assumption.
\end{VNBcond}
No arguments.  The operation tries the number reading first and the
structure reading second.

An existential assertion is outside this operation and every other
oracle except \texttt{arith-forsome}'s one shape: none of them supplies a
witness.  For $\exists a \in \mathbb{Z}.\, (x+y)^3 = x^3 + y \cdot a$ the
witness is given with \texttt{forsome-intro} (command \texttt{ew}), and
the equation that remains is an identity within reach of
\texttt{comm-ring-simplify}.

\subsubsection*{\texttt{ineq}: an oracle for linear arithmetic over the reals
  (command \texttt{ineq})}

%% rco--check-ineq :556, rco--farkas-ok? :515, rco--ineq-constraint :494, rule-checkers-oracle.scm
\VNBfig{ineq $c$}{\VNBnohyp}{\Gamma \Longrightarrow \forall x_1.\, x_1 \in \mathbb{R} \Rightarrow \cdots \forall x_k.\, x_k \in \mathbb{R} \Rightarrow G_0}
\begin{VNBcond}
\item[Certificate] $c$ is a list of pairs $(\mathit{id}\ .\ m)$,
  each $m$ an exact nonnegative rational, where $\mathit{id}$ is
  \texttt{goal}, the negation of $G_0$; or \texttt{h}$I$, the $I$-th
  assumption of $N$ counted from one; or \texttt{h}$I$\texttt{-rev}, the
  $I$-th assumption, an equation, read from right to left.  When $G_0$
  is an equation, $c$ is a pair of two such lists, one for each
  direction.
\item[Premises] Each assumption the certificate names is an equation or
  a $\leq$ or $<$ inequality.
\item[Absurdity] The sum of the named constraints, each written
  $\ell \leq 0$ or $\ell < 0$ and multiplied by its $m$, has every atom's
  coefficient zero, and its constant is $> 0$, or $\geq 0$ when a strict
  constraint enters with positive multiplier.
\item[Typing] Every atom of a constraint that enters with a positive
  multiplier is certified real: $\Gamma$ contains $a \in \mathbb{R}$, or $a$ is one of the
  $x_i$, or $a$ is $|t|$ with every atom of $t$ certified.
\item[Eigenvariable] No $x_i$ is free in a premise the certificate
  names.
\end{VNBcond}
The arguments of the operation are the positions in $\Gamma$ of the
assumptions to be used as premises, counted from one.  When the
procedure succeeds $N$ is grounded, and otherwise the operation
declines.

The assertion and each named premise are read as linear equations or
inequalities in their \emph{atoms}, the maximal subterms that are not
sums, differences, or products by a numeral.  A named premise that does
not read this way is left out, which can only make the procedure prove
less; an index outside the range of $\Gamma$ is an error and abandons the
call.  The operation requires the Typing condition of every atom, an
atom that cancels included, and declines when a variable bound by one of
the leading quantifiers occurs free in \emph{any} assumption of $N$.
Fourier--Motzkin elimination is run on the premises and the negation of
the assertion.  The procedure succeeds when the elimination reaches
$0 < 0$.  The nonnegative combination of premises that produces it, a
\emph{Farkas certificate}, is printed and is recorded in the graph with
the name of the operation.

\begin{remark}
The condition on bound variables is the eigenvariable condition of
\texttt{forall-intro}.  Until 2026-09-20 the oracle did not impose it,
and from $x \in \mathbb{R},\; x \leq 0$ it closed
$\forall x \in \mathbb{R}.\; x \leq 0$.  The defect was found when the
checker of Section~\ref{sec:inference-checked} was written for this
operation, and no proof in the library had used it.
\end{remark}

\subsubsection*{\texttt{sos}: nonnegativity by a sum of squares
  (command \texttt{sos})}

%% rco--check-sos :956, rule-checkers-oracle.scm; tag built at structure-library/sos-oracle.scm:120
\VNBfig{sos $c$}{\VNBnohyp}{\Gamma \Longrightarrow \forall x_1.\, x_1 \in \mathbb{R} \Rightarrow \cdots \forall x_k.\, x_k \in \mathbb{R} \Rightarrow a \leq b}
\begin{VNBcond}
\item[Certificate] $c$ is a list of pairs $(c_i\ .\ \lambda_i)$, each
  $\lambda_i$ an exact nonnegative rational.
\item[Identity] $b - a = \lambda_1 c_1^2 + \cdots + \lambda_n c_n^2$ as
  polynomials in the free commutative ring over the generators.
\item[Typing] Every generator of $a$, $b$ and the $c_i$ is certified
  real, as for \texttt{ineq}.
\end{VNBcond}
The assertion may also be written $b \geq a$.  The arguments of the
operation are the terms $c_1, \ldots, c_n$ to be squared.  An exact
linear program decides whether weights $\lambda_i \geq 0$ exist; on
success the pairs $(c_i\ .\ \lambda_i)$ are printed and recorded in the
graph with the name.

The oracle verifies a certificate, it does not search for one: the
squares are supplied by the caller.  A strict inequality is not accepted,
since a square may be zero.

\subsection{Which command reaches which operation}\label{sec:command-uses}

A tactic has no access to the graph except through the operations above,
so the operations one tactic can reach are the bound on what it can do.
The table below gives that bound for each proof command: the operations
recorded by any kernel procedure reachable from the command's own code,
whether or not the library happens to exercise it.  It is produced from
\texttt{reference/kernel-map-static.sexp}, the result of reading the
source of every loaded file with the Scheme reader.  One thing is left
out of every row: after any command that changes a proof, the system tries
to close the definedness obligations that instantiation has left
(Section~\ref{sec:kernel-ops-detail}, \texttt{forall-elim}), and that step
can reach twenty kernel procedures of its own.  It is the same for every
command and is stated once in \texttt{reference/KERNEL-MAP.md}.

A command listed as \texttt{none} records nothing.  Those are the
navigation, search and session commands.

%% The table and its counts, generated by gen-tactic-uses.py from
%% reference/kernel-map-static.sexp and the command registry.
%% GENERATED by gen-tactic-uses.py from reference/kernel-map-static.sexp and
%% tactics-help.scm (*tactic-kind*).  Do not edit.
\newcommand{\VNBtacticusescount}{105}
\newcommand{\VNBtacticusesmax}{64}
\newcommand{\VNBtacticusestable}{%
\begin{longtable}{@{}p{0.17\textwidth}p{0.79\textwidth}@{}}
\toprule
command & operations it can reach \\
\midrule
\endfirsthead
\toprule
command & operations it can reach \\
\midrule
\endhead
\texttt{ai} & \texttt{and-elim} \texttt{forsome-elim} \texttt{iff-elim} \texttt{not-elim} \texttt{or-elim}\\
\texttt{apply-thm} & \texttt{forall-elim} \texttt{theorem-assumption}\\
\texttt{arith} & \texttt{arith-forsome} \texttt{arith-ground} \texttt{arith-simplify} \texttt{cartesian-decompose} \texttt{eq-subst} \texttt{macete} \texttt{tuple-equality-decompose}\\
\texttt{arith--decide} & \texttt{arith-forsome} \texttt{arith-ground} \texttt{arith-simplify}\\
\texttt{ass} & \texttt{assumption}\\
\texttt{ass-all} & \texttt{assumption}\\
\texttt{backup-one} & none\\
\texttt{bc} & \texttt{backchain}\\
\texttt{bc*} & \texttt{assumption} \texttt{backchain} \texttt{cut} \texttt{forall-elim} \texttt{theorem-assumption}\\
\texttt{bc*--attempt} & \texttt{assumption} \texttt{backchain} \texttt{cut} \texttt{forall-elim} \texttt{theorem-assumption}\\
\texttt{beta} & \texttt{functoid-beta}\\
\texttt{bu-me} & \texttt{big-union-mem-elim}\\
\texttt{bu-mi} & \texttt{big-union-mem-intro}\\
\texttt{bu-set} & \texttt{big-union-sethood}\\
\texttt{calc} & \emph{all 64 reachable operations}\\
\texttt{ce} & \texttt{cartesian-elim}\\
\texttt{choose} & \texttt{and-elim} \texttt{and-intro} \texttt{arith-forsome} \texttt{arith-ground} \texttt{arith-simplify} \texttt{assumption} \texttt{cartesian-decompose} \texttt{cut} \texttt{detach} \texttt{eq-subst} \texttt{forall-elim} \texttt{forall-intro} \texttt{forsome-elim} \texttt{forsome-intro} \texttt{iff-elim} \texttt{iff-intro} \texttt{implies-intro} \texttt{macete} \texttt{not-elim} \texttt{not-intro} \texttt{or-elim} \texttt{theorem-assumption} \texttt{tuple-equality-decompose}\\
\texttt{choose-pos} & \texttt{and-elim} \texttt{and-intro} \texttt{arith-forsome} \texttt{arith-ground} \texttt{arith-simplify} \texttt{assumption} \texttt{cartesian-decompose} \texttt{cut} \texttt{detach} \texttt{eq-subst} \texttt{forall-elim} \texttt{forall-intro} \texttt{forsome-elim} \texttt{forsome-intro} \texttt{iff-elim} \texttt{iff-intro} \texttt{implies-intro} \texttt{macete} \texttt{not-elim} \texttt{not-intro} \texttt{or-elim} \texttt{theorem-assumption} \texttt{tuple-equality-decompose}\\
\texttt{ci} & \texttt{cartesian-intro}\\
\texttt{comp-me} & \texttt{comp-mem-elim}\\
\texttt{comp-mi} & \texttt{comp-mem-intro}\\
\texttt{contra} & \texttt{and-elim} \texttt{and-intro} \texttt{arith-forsome} \texttt{arith-ground} \texttt{arith-simplify} \texttt{assumption} \texttt{cartesian-decompose} \texttt{cut} \texttt{detach} \texttt{eq-subst} \texttt{forall-elim} \texttt{forall-intro} \texttt{forsome-elim} \texttt{iff-elim} \texttt{iff-intro} \texttt{implies-intro} \texttt{ineq} \texttt{macete} \texttt{not-elim} \texttt{not-intro} \texttt{or-elim} \texttt{theorem-assumption} \texttt{tuple-equality-decompose}\\
\texttt{crs} & \texttt{comm-ring-simplify}\\
\texttt{cut} & \texttt{cut}\\
\texttt{detach!} & \texttt{detach}\\
\texttt{di} & \texttt{and-intro} \texttt{forall-intro} \texttt{iff-intro} \texttt{implies-intro} \texttt{not-intro}\\
\texttt{dial} & none\\
\texttt{dial-wff} & none\\
\texttt{dk-focus!} & none\\
\texttt{eps-part} & \texttt{and-elim} \texttt{and-intro} \texttt{arith-forsome} \texttt{arith-ground} \texttt{arith-simplify} \texttt{assumption} \texttt{cartesian-decompose} \texttt{cut} \texttt{detach} \texttt{eq-subst} \texttt{forall-elim} \texttt{forall-intro} \texttt{forsome-elim} \texttt{iff-elim} \texttt{iff-intro} \texttt{implies-intro} \texttt{macete} \texttt{not-elim} \texttt{not-intro} \texttt{or-elim} \texttt{theorem-assumption} \texttt{tuple-equality-decompose}\\
\texttt{ew} & \texttt{forsome-intro}\\
\texttt{fact} & \texttt{arith-forsome} \texttt{arith-ground} \texttt{arith-simplify} \texttt{cut} \texttt{detach} \texttt{forall-elim} \texttt{theorem-assumption}\\
\texttt{focus} & none\\
\texttt{focus-id} & none\\
\texttt{goal-status} & none\\
\texttt{grind} & \texttt{and-elim} \texttt{and-intro} \texttt{forall-intro} \texttt{forsome-elim} \texttt{iff-elim} \texttt{iff-intro} \texttt{implies-intro} \texttt{macete-hyp} \texttt{not-elim} \texttt{not-intro} \texttt{or-elim}\\
\texttt{grind-and-mp} & \texttt{and-elim} \texttt{and-intro} \texttt{arith-forsome} \texttt{arith-ground} \texttt{arith-simplify} \texttt{assumption} \texttt{cut} \texttt{detach} \texttt{forall-elim} \texttt{forall-intro} \texttt{forsome-elim} \texttt{iff-elim} \texttt{iff-intro} \texttt{implies-intro} \texttt{macete-hyp} \texttt{not-elim} \texttt{not-intro} \texttt{or-elim}\\
\texttt{have!} & \texttt{and-intro} \texttt{arith-forsome} \texttt{arith-ground} \texttt{arith-simplify} \texttt{assumption} \texttt{cartesian-decompose} \texttt{cut} \texttt{detach} \texttt{eq-subst} \texttt{forall-elim} \texttt{forall-intro} \texttt{iff-intro} \texttt{implies-intro} \texttt{macete} \texttt{not-intro} \texttt{theorem-assumption} \texttt{tuple-equality-decompose}\\
\texttt{ie} & \texttt{intersection-elim}\\
\texttt{if-false} & \texttt{if-false}\\
\texttt{if-true} & \texttt{if-true}\\
\texttt{ii} & \texttt{intersection-intro}\\
\texttt{in-rr} & \texttt{and-intro} \texttt{arith-forsome} \texttt{arith-ground} \texttt{arith-simplify} \texttt{assumption} \texttt{cartesian-decompose} \texttt{cartesian-intro} \texttt{cut} \texttt{detach} \texttt{eq-subst} \texttt{forall-elim} \texttt{forall-intro} \texttt{iff-intro} \texttt{implies-intro} \texttt{macete} \texttt{not-intro} \texttt{theorem-assumption} \texttt{tuple-equality-decompose}\\
\texttt{in-rr--refocus!} & none\\
\texttt{ineq} & \texttt{ineq}\\
\texttt{inst} & \texttt{forall-elim}\\
\texttt{inst+} & \texttt{arith-forsome} \texttt{arith-ground} \texttt{arith-simplify} \texttt{cut} \texttt{detach} \texttt{forall-elim}\\
\texttt{iota-d} & \texttt{iota-def}\\
\texttt{iota-e} & \texttt{iota-in-elim}\\
\texttt{lam-b} & \texttt{lambda-beta}\\
\texttt{lam-b-h} & \texttt{lambda-beta-hyp}\\
\texttt{lam-t} & \texttt{lambda-type}\\
\texttt{len-r} & \texttt{length-reduce}\\
\texttt{mac} & \texttt{cartesian-decompose} \texttt{macete} \texttt{tuple-equality-decompose}\\
\texttt{mac-h} & \texttt{macete-hyp}\\
\texttt{mac-h*} & \texttt{and-elim} \texttt{forsome-elim} \texttt{iff-elim} \texttt{macete-hyp} \texttt{not-elim} \texttt{or-elim}\\
\texttt{macm} & \texttt{cartesian-decompose} \texttt{macete} \texttt{tuple-equality-decompose}\\
\texttt{minimize!} & \texttt{and-elim} \texttt{and-intro} \texttt{arith-forsome} \texttt{arith-ground} \texttt{arith-simplify} \texttt{assumption} \texttt{cartesian-decompose} \texttt{cut} \texttt{detach} \texttt{eq-subst} \texttt{forall-elim} \texttt{forall-intro} \texttt{forsome-elim} \texttt{forsome-intro} \texttt{iff-elim} \texttt{iff-intro} \texttt{implies-intro} \texttt{macete} \texttt{not-elim} \texttt{not-intro} \texttt{or-elim} \texttt{reflexivity} \texttt{sep-mem-elim} \texttt{sep-mem-intro} \texttt{theorem-assumption} \texttt{tuple-equality-decompose}\\
\texttt{mp} & \texttt{arith-forsome} \texttt{arith-ground} \texttt{arith-simplify} \texttt{assumption} \texttt{cut} \texttt{detach} \texttt{forall-elim}\\
\texttt{ni} & \texttt{nn-induction}\\
\texttt{nth-r} & \texttt{nth-reduce}\\
\texttt{obtain} & \emph{all 64 reachable operations}\\
\texttt{obtain-at} & \texttt{and-elim} \texttt{and-intro} \texttt{arith-forsome} \texttt{arith-ground} \texttt{arith-simplify} \texttt{assumption} \texttt{cartesian-decompose} \texttt{cut} \texttt{detach} \texttt{eq-subst} \texttt{forall-elim} \texttt{forall-intro} \texttt{forsome-elim} \texttt{iff-elim} \texttt{iff-intro} \texttt{implies-intro} \texttt{macete} \texttt{not-elim} \texttt{not-intro} \texttt{or-elim} \texttt{theorem-assumption} \texttt{tuple-equality-decompose}\\
\texttt{oi-l} & \texttt{or-intro-left}\\
\texttt{oi-r} & \texttt{or-intro-right}\\
\texttt{orelse} & none\\
\texttt{pbc} & \texttt{proof-by-contradiction}\\
\texttt{prep} & none\\
\texttt{prop} & \texttt{and-elim} \texttt{and-intro} \texttt{arith-forsome} \texttt{arith-ground} \texttt{arith-simplify} \texttt{assumption} \texttt{cartesian-decompose} \texttt{cut} \texttt{detach} \texttt{eq-subst} \texttt{forall-elim} \texttt{forall-intro} \texttt{forsome-elim} \texttt{iff-elim} \texttt{iff-intro} \texttt{implies-intro} \texttt{macete} \texttt{not-elim} \texttt{not-intro} \texttt{or-elim} \texttt{or-intro-left} \texttt{or-intro-right} \texttt{proof-by-contradiction} \texttt{theorem-assumption} \texttt{tuple-equality-decompose}\\
\texttt{push-not-h} & \texttt{and-elim} \texttt{and-intro} \texttt{arith-forsome} \texttt{arith-ground} \texttt{arith-simplify} \texttt{assumption} \texttt{cartesian-decompose} \texttt{cut} \texttt{detach} \texttt{eq-subst} \texttt{forall-elim} \texttt{forall-intro} \texttt{forsome-elim} \texttt{forsome-intro} \texttt{iff-elim} \texttt{iff-intro} \texttt{implies-intro} \texttt{macete} \texttt{not-elim} \texttt{not-intro} \texttt{or-elim} \texttt{or-intro-left} \texttt{or-intro-right} \texttt{proof-by-contradiction} \texttt{theorem-assumption} \texttt{tuple-equality-decompose}\\
\texttt{qed} & none\\
\texttt{qrfl} & \texttt{quasi-reflexivity}\\
\texttt{quietly} & none\\
\texttt{repeat} & none\\
\texttt{replay-proof} & \emph{all 64 reachable operations}\\
\texttt{rfl} & \texttt{reflexivity}\\
\texttt{rs} & \texttt{ring-simplify}\\
\texttt{save-proof} & none\\
\texttt{scout} & \emph{all 64 reachable operations}\\
\texttt{scout-run} & none\\
\texttt{scout-show} & \emph{all 64 reachable operations}\\
\texttt{sep-me} & \texttt{sep-mem-elim}\\
\texttt{sep-mi} & \texttt{sep-mem-intro}\\
\texttt{sep-set} & \texttt{sep-sethood}\\
\texttt{show} & none\\
\texttt{simp} & \texttt{comm-ring-simplify} \texttt{cut} \texttt{eq-subst}\\
\texttt{slot} & \texttt{cartesian-decompose} \texttt{macete} \texttt{tuple-equality-decompose}\\
\texttt{slot-h} & \texttt{macete-hyp}\\
\texttt{sos} & \texttt{sos}\\
\texttt{sp} & none\\
\texttt{subst} & \texttt{eq-subst}\\
\texttt{supply} & \emph{all 64 reachable operations}\\
\texttt{ta} & \texttt{theorem-assumption}\\
\texttt{te} & \texttt{tuples-elim}\\
\texttt{tfi} & \texttt{transfinite-induction}\\
\texttt{tfi3} & \texttt{transfinite-induction-3cases}\\
\texttt{ti} & \texttt{tuples-intro}\\
\texttt{ue} & \texttt{union-elim}\\
\texttt{ui} & \texttt{union-intro}\\
\texttt{undo} & none\\
\texttt{use-at} & \texttt{and-intro} \texttt{arith-forsome} \texttt{arith-ground} \texttt{arith-simplify} \texttt{assumption} \texttt{cartesian-decompose} \texttt{cut} \texttt{detach} \texttt{eq-subst} \texttt{forall-elim} \texttt{forall-intro} \texttt{iff-intro} \texttt{implies-intro} \texttt{macete} \texttt{not-intro} \texttt{theorem-assumption} \texttt{tuple-equality-decompose}\\
\texttt{use-em} & \texttt{and-elim} \texttt{and-intro} \texttt{arith-forsome} \texttt{arith-ground} \texttt{arith-simplify} \texttt{assumption} \texttt{cartesian-decompose} \texttt{cut} \texttt{detach} \texttt{eq-subst} \texttt{forall-elim} \texttt{forall-intro} \texttt{forsome-elim} \texttt{iff-elim} \texttt{iff-intro} \texttt{implies-intro} \texttt{macete} \texttt{not-elim} \texttt{not-intro} \texttt{or-elim} \texttt{or-intro-left} \texttt{or-intro-right} \texttt{proof-by-contradiction} \texttt{theorem-assumption} \texttt{tuple-equality-decompose}\\
\texttt{vlet} & \texttt{and-elim} \texttt{cut} \texttt{forsome-elim} \texttt{iff-elim} \texttt{not-elim} \texttt{or-elim}\\
\texttt{wbc} & \texttt{and-elim} \texttt{arith-forsome} \texttt{arith-ground} \texttt{arith-simplify} \texttt{cut} \texttt{detach} \texttt{forall-elim} \texttt{forsome-elim} \texttt{iff-elim} \texttt{not-elim} \texttt{or-elim} \texttt{theorem-assumption}\\
\texttt{wk} & \texttt{weakening}\\
\bottomrule
\end{longtable}
\noindent\footnotesize \textsuperscript{r}: read from the command registry rather
than from the kernel map, which predates the command.\normalsize}

\VNBtacticusestable

\subsection{How an inference is recorded, and checked}\label{sec:inference-checked}

Every kernel operation ends by calling one procedure,
\texttt{dg-apply-rule!}\ (\texttt{deduction-graphs.scm}).  It is handed
the rule, the hypothesis sequents and the node $N$.  It checks the
inference (below); then it posts each hypothesis sequent as a node,
reusing a node the graph already has for the same sequent; then it
creates the inference node and writes the arrows.  No other procedure
writes to the graph.
%% dg-apply-rule!, deduction-graphs.scm:375; make-inference-node :59
The inference node records exactly three things: the rule, the list of
hypothesis nodes, and the conclusion node $N$.  The rule is the name of
the operation, or a list headed by that name when the operation carries
data about the application: an index (\texttt{(cartesian-elim 2)}), the
name and the patterns of a macete (\texttt{(macete \textit{name} $L$
$R$)}), or the certificate of an oracle (\texttt{(ineq $c$)},
\texttt{(sos $c$)}).

\subsubsection*{The construction and the check}

An operation has two parts.  The \emph{builder} is a procedure
\texttt{pi-\ldots!}\ (in \texttt{primitive-inferences.scm},
\texttt{macetes.scm} and the oracle files).  It reads $N$ and the
operation's argument and \emph{constructs} the hypothesis sequents: for
\texttt{and-intro}, \texttt{pi-direct-inference!} takes the two conjuncts
of $G$ and makes $\Gamma \Longrightarrow A$ and
$\Gamma \Longrightarrow B$; for \texttt{forall-intro} it also chooses the
eigenvariables; for an oracle it runs the decision procedure.  The
\emph{checker} is a separate procedure, registered for the operation's
name in one of four files (\texttt{rule-checkers-logic.scm},
\texttt{-schema.scm}, \texttt{-rewrite.scm}, \texttt{-oracle.scm}).
\texttt{dg-apply-rule!}\ hands it the finished result: the rule, the
hypothesis sequents and the sequent of $N$.  The checker does not call
the builder and does not compute hypotheses from $N$.  It decides whether
the sequents it is given stand in the relation of the operation's figure
(Section~\ref{sec:kernel-ops-detail}), and answers yes, or no with a
reason.

The checkers of the logic and schema operations use four things from the
expression layer: substitution, the free-variable computation, equality
of formulas up to renaming of bound variables ($\aeq$), and one-way
matching.  The rewrite checkers add two: a theorem is read as a rewrite
rule by the same procedure the rewriter uses
(\texttt{extract-rewrite-patterns}), because that reading is what the
name of a macete denotes; and a closed arithmetic condition is evaluated
by the arithmetic evaluator.  The oracle checkers carry their own exact
arithmetic, linear forms and polynomials, written separately from the
oracles.
%% headers of rule-checkers-logic.scm (1-13), rule-checkers-rewrite.scm (18-26),
%% rule-checkers-oracle.scm (1-14)

Two of the relations, written out as the checker tests them:
%% and-intro checker, rule-checkers-logic.scm:347; forall-elim checker :702
\begin{itemize}
\item \texttt{and-intro} accepts hypotheses $H_1, H_2$ and conclusion
  $\Gamma \Longrightarrow G$ exactly when $G$ is $A \wedge B$, there are
  two hypotheses, the context of each equals $\Gamma$ as a set, the
  assertion of $H_1$ is $\aeq A$ and that of $H_2$ is $\aeq B$.
\item \texttt{forall-elim} accepts exactly when there are one or two
  hypotheses, the assertion of $H_1$ is $\aeq G$, and either the context
  of $H_1$ is included in $\Gamma$, or $\Gamma$ contains a formula
  $\forall x.\, B$ and the context of $H_1$ is $\Gamma \cup \{D\}$ for a
  formula $D$ that matching shows to be $B[x{:=}t]$ for some term $t$;
  and, if there is a second hypothesis, it is
  $\Gamma \Longrightarrow u = u$ for some term $u$.
\end{itemize}

The eigenvariable condition of \texttt{forall-intro} is tested as
follows.  The checker counts the length of
the chain of leading universal quantifiers of $G$, passing through each
guard $x \in A$ on the variable just bound, and tries that length $k$
first and then each smaller one.  For a given $k$ it reads the
eigenvariables $y_1, \ldots, y_k$ off the hypothesis
$\Gamma' \Longrightarrow B'$ by matching: $G$ with its $k$ bound
variables replaced by holes is matched, guard by guard, against the
formulas of $\Gamma'$ that are not in $\Gamma$, and its body against
$B'$.  A hole that matching does not reach belongs to a variable absent
from the body, and is given the bound variable's own name.  It then
walks the chain forwards with a set $V$ of excluded variables, which
starts as $\FV(\Gamma)$.  At step $i$ the formula is
$\forall x_i.\, C_i$, and $y_i$ must be a variable, not in $V$, and not
free in $\forall x_i.\, C_i$.  If $C_i[x_i{:=}y_i]$ is
$(y_i \in A_i) \Rightarrow D$, the guard $y_i \in A_i$ is collected, its
free variables ($y_i$ and those of $A_i$) are added to $V$, and the walk
continues with $D$; otherwise $y_i$ is added to $V$ and the walk
continues with $C_i[x_i{:=}y_i]$.  The inference is accepted when the
formula reached after $k$ steps is $\aeq B'$ and $\Gamma'$ is $\Gamma$
together with the collected guards.
%% chk-forall-intro-at, rule-checkers-logic.scm:514; chk-forall-eigen-ok? :491;
%% chk-forall-fill :477; chk-vacuous-name :540; the k-loop in the checker at :555

If no checker is registered for the name, or the checker refuses,
\texttt{dg-apply-rule!}\ raises an error and the graph is unchanged.  A
refusal is printed even when output is suppressed, and the load ends by
reporting the number of inferences verified and the number refused.
%% dg-check-inference!, dg-check-report-refusal!, deduction-graphs.scm

An error in a builder goes unnoticed only if the checker's relation also
admits the inference the builder produced.  Where a figure of
Section~\ref{sec:kernel-ops-detail} carries a condition labelled
\emph{Not checked}, the relation is weaker than the operation, and at
that point the operation is the only safeguard.  There are four such
places: definedness in \texttt{forall-elim} and \texttt{reflexivity}, by
policy, since the certificate decides when an obligation may be omitted
and the checker accepts the inference either way; capture in
\texttt{eq-subst}; the variable of the second hypothesis of
\texttt{big-union-sethood}; and the two refusals of
\texttt{macete-hyp}, which are not needed for soundness.
%% policy: header of rule-checkers-logic.scm, lines 29-43

\subsubsection*{Oracles: a certificate check or a second decision}

An oracle has no proof to replay, so what its checker can establish
depends on what the oracle records.  There are two cases.
%% header of rule-checkers-oracle.scm, lines 22-62

\emph{Certificate check} (\texttt{ineq}, \texttt{sos}).  The oracle
records, with its name, the object that shows the assertion true, and
the checker verifies that object by arithmetic of its own, without
elimination, search or any code of the oracle.  For \texttt{ineq} the
object is a Farkas certificate: the checker reads the assertion and the
named premises as linear forms, multiplies each by its multiplier, adds,
and confirms that every atom cancels and that the constant left is an
absurd inequality.  For \texttt{sos} it is the list of squares and
weights: the checker expands $\lambda_1 c_1^2 + \cdots + \lambda_n c_n^2$
and $b - a$ as polynomials and compares them monomial by monomial.  Either
the arithmetic confirms the certificate, or the inference is refused.
%% rco--farkas-ok? :515, rco--check-sos :956, rule-checkers-oracle.scm

\begin{sloppypar}
\emph{Second decision} (\texttt{arith-ground}, \texttt{arith-forsome},
\texttt{arith-simplify}, \texttt{comm-ring-simplify},
\texttt{ring-simplify}).  These oracles record their name only, so the
checker decides the same question again with a second implementation of
the same mathematics: exact evaluation for \texttt{arith}, expansion of
both sides to a sum of monomials for the two ring oracles.  The two
implementations share no code beyond the expression layer.  A second
decision catches an error in how the oracle reads its input or reports
its answer: an atom misread, a typing it failed to require, a premise
dropped, a case it mishandles that the second implementation handles
correctly.  It does not catch an error that both implementations make
the same way, for instance a mistake in the mathematics of the normal
form itself.
\end{sloppypar}
%% rco--check-arith-ground :242, rco--check-crs :929, rco--check-rs :939

Neither kind of check adds proving power.  A checker only judges an
inference that an oracle has already produced; an assertion the oracle
declines is not proved by its checker.

\subsubsection*{Where a certificate is kept, and how it is shown}

The certificate of \texttt{ineq} is part of the rule recorded in the
inference node, beside the name of the operation.  The rule has the form
\texttt{(ineq $c$)}, where $c$ is a list of pairs
\texttt{($\mathit{id}$ .\ $m$)}: $m$ is an exact nonnegative rational,
and $\mathit{id}$ is \texttt{goal} (the negation of the assertion),
\texttt{h}$I$ (the $I$-th assumption of $N$, counted from one), or
\texttt{h}$I$\texttt{-rev} (the $I$-th assumption, an equation, read from
right to left).  For an assertion that is an equation, $c$ is a pair of
two such lists, one for each direction.  For example
\texttt{(ineq ((h1 .\ 1) (h3 .\ 2) (goal .\ 1)))} says that the first
assumption, twice the third, and the negated assertion add up to a
false inequality between constants.  The rule of \texttt{sos} is
\texttt{(sos (($c_1$ .\ $\lambda_1$) \ldots))}.
%% pi-ineq!, structure-library/ineq-oracle.scm:218 (tag built at :308);
%% ineq-hyp-constraints :167; sos-oracle.scm:120

The certificate therefore belongs to the deduction graph of the proof,
and lasts as long as that graph.  What \texttt{qed} keeps of a proof is
the statement, the script and the bill; the script records the command
as it was typed, for instance \texttt{(ineq 1 3)}, not the certificate.
Replaying the script, as the page audit does, runs the oracle again, and
the checker verifies the certificate recorded by that run.
%% proof-debt.scm:205 (*proof-oracles* holds oracle names, not certificates)

When it succeeds the oracle prints the certificate:
\begin{VNBcode}
;; ineq: closed by Farkas certificate ((h1 . 1) (h3 . 2) (goal . 1))
\end{VNBcode}
\texttt{sos} prints its decomposition in the same way.  \texttt{proof-tex}
prints, for each step, the command and the sequent it leaves, read from
the command trace recorded while the proof ran; it does not print the
certificate.  The command \texttt{(proof-rules)} reads the graph of the
current proof and lists the rules its inferences recorded.
%% ineq-report, ineq-oracle.scm:215; sos-report, sos-oracle.scm:49;
%% proof-tex--steps-from-live, proof-tex.scm:94; proof-rules, suggest.scm:5795

\subsubsection*{Controls}

Each checker has two controls in the test suite: the proofs of the
library, which it must accept, and an inference that is wrong in a
stated way, which it must refuse.  A checker that accepted everything
would pass the first and fail the second.

\begin{remark}
Checking is on by default.  Setting the environment variable
\texttt{VNB\_NO\_RULE\_CHECK} turns it off.  On the library as loaded on
2026-09-20, 197{,}677 inferences were verified and none refused, and the
load took 30\% longer with checking than without.
\end{remark}

\subsection{The boundary is checked, not merely documented}\label{sec:kernel-checked}

The restriction would be worth little if it were a convention.  The eight
files permitted to record inferences are named in a list in the source
(\texttt{*kernel-caller-files*}).  Every time the system starts, it reads
all of its own source with the Scheme reader and counts, in each file,
the calls to \texttt{dg-apply-rule!}\ and any use of that procedure as a
value --- the second being the way a file could otherwise hand the
capability to code outside the list.  If any file not on the list records
inferences, or if any file passes the procedure on, the system reports
the offending files and \emph{refuses to finish starting}.  A reader is
used rather than a text search so that a mention inside a comment, a
string or quoted data is not mistaken for code.

On a clean tree the check prints one line:

\begin{VNBcode}
;; kernel-callers-audit: ok (69 call site(s) to dg-apply-rule!
;;   in 8 file(s), all allowed; 0 value use(s))
\end{VNBcode}

A companion check, \texttt{kernel-rules-audit}, compares the rules
actually recorded during the load against the documented list in both
directions, so a rule that is trusted but undocumented, or documented but
unreachable, is reported rather than left to be noticed.

Widening the trusted base is possible --- a new decision procedure is a
legitimate thing to add --- but it cannot happen quietly: the new file
must be added to the list by hand, which is a deliberate act that leaves
a record.  The complete inventory (every entry point, the rules it may
record, and which proof commands reach it) is
\texttt{reference/KERNEL-MAP.md}.

%% =========================================================
\section{Starting and inspecting a proof}

\begin{VNBcode}
; Start a new proof of a formula
(define ps (start-proof (make-wff-from-string
  "n in NN and m in NN implies n in NN")))

; Display the current state
(print-proof-state ps)

; Check if done
(proof-done? ps)   ; => #f or #t

; List open goals
(proof-open-goals ps)
\end{VNBcode}

The variable \texttt{ps} holds the \emph{proof state}: a deduction graph
plus the current \emph{focus} sequent node (the subgoal currently being
worked on).  Every proof command returns a new proof state with an
updated focus.

\texttt{start-proof} takes a \texttt{wff} (produced by \texttt{make-wff} or \texttt{make-wff-from-string}).
The root sequent has no assumptions: a proof begins from the bare goal,
and every hypothesis arrives through an inference.  All sequent assumptions and assertions are
\texttt{wff} objects throughout the proof; the library name propagates
automatically to every subformula produced by inference rules.

%% =========================================================
\section{Proof commands}
\label{sec:commands}

Commands are Scheme procedures of the form
\[
  \texttt{(cmd-}\mathit{name}\;\; \mathit{ps}\;\; \mathit{arg}_1 \ldots\texttt{)}
\]
Each command applies one primitive inference to the focus node, then
advances the focus to the first new ungrounded subgoal.

\subsection{Logical connectives}

\begin{itemize}

\item \texttt{(cmd-direct-inference ps)}
  Decompose the goal by its top-level connective:
  \begin{itemize}
  \item \texttt{p and q} $\to$ two subgoals $p$, $q$.
  \item \texttt{p implies q} $\to$ assume $p$, prove $q$.
  \item \texttt{p iff q} $\to$ two subgoals (each direction).
  \item \texttt{not(p)} $\to$ assume $p$, prove \texttt{falsity}.
  \item \texttt{forall([\ldots], p)} $\to$ peel the entire leading
    \texttt{forall} chain in one step, introducing a fresh variable for
    each bound variable.  If the substituted body is
    \texttt{$y$ in $A$ implies \ldots}, absorb \texttt{$y$ in $A$} into
    the assumptions and continue peeling.  Stop at the first
    non-\texttt{forall} head.
    Example: goal \texttt{forall([x in A, y], body)} yields one subgoal
    $x_0 \in A \;\Longrightarrow\; \texttt{body}[x/x_0,\,y/y_0]$.
  \end{itemize}

\item \texttt{(cmd-antecedent-inference ps $f$)}
  Decompose assumption $f$ by its top-level connective:
  \begin{itemize}
  \item \texttt{p and q} $\to$ replace with $p$ and $q$ separately.
  \item \texttt{p or q} $\to$ case split: branch on $p$, branch on $q$.
  \item \texttt{not(p)} $\to$ derive \texttt{falsity} if $p$ also in context.
  \item \texttt{forsome([\ldots], p)} $\to$ introduce a fresh witness for $x$.
  \end{itemize}

\item \texttt{(cmd-assumption ps)}
  Close goal immediately if it matches an assumption or is \texttt{truth}.

\item \texttt{(cmd-arith ps)}
  Close the current goal by ground arithmetic evaluation.  The goal must
  be a closed formula (no free variables) built from numeric literals,
  $+$, $*$, $-$, \texttt{recip}, \texttt{abs}, \texttt{succ},
  \texttt{conjugate}, $\leq$, $=$,
  $\in \mathit{NN}/\mathit{ZZ}/\mathit{QQ}/\mathit{RR}/\mathit{CC}$,
  and the propositional connectives.  Exact (rational) arithmetic only;
  no axiom chain is constructed.
  Related: \texttt{(vnb-do-arith \textit{f})} returns \texttt{\#t}/\texttt{\#f}
  for use outside a proof.

\item \texttt{(cmd-ring-simplify ps)}
  Close a goal by reducing both sides to their canonical form in the free
  associative $\mathbb{Z}$-algebra.  Addition is commutative (both sides
  are brought to the same sorted canonical order); multiplication is
  non-commutative.  Generators are detected automatically and no generator
  list need be supplied: \texttt{(rs)} treats any \emph{term} that
  \texttt{arith-eval} cannot reduce to a number as a generator of the free
  ring --- a bare symbol such as \texttt{x}, or a compound term such as
  \texttt{f(x)} whose head is not one of $+\,-\,*$.  A compound generator is
  a single opaque indeterminate: \texttt{f(x)} is atomic and its argument is
  not examined.

  The goal may be in either of two forms:
  \begin{itemize}
  \item \textbf{Bare equality} \texttt{$e_1$ = $e_2$}: generators are
    certified by \texttt{(IN $v$ $D$)} assumptions already in the sequent
    context.  Typically reached after \texttt{(di)} peels the leading
    \texttt{forall} chain.
  \item \textbf{Universally quantified identity}
    \texttt{forall([$x_1$ in $D_1$, \ldots], $e_1$ = $e_2$)}: the
    quantifier typings themselves certify the generators.  \texttt{(di)}
    is not needed; \texttt{(rs)} closes the goal in one step.
  \end{itemize}

  In both cases every generator must be certified as an element of one of
  the five ring domains $\mathit{NN}$, $\mathit{ZZ}$, $\mathit{QQ}$,
  $\mathit{RR}$, $\mathit{CC}$ --- a symbol \texttt{x} by an
  \texttt{(IN x $D$)} typing, a compound generator \texttt{f(x)} by an
  \texttt{(IN f(x) $D$)} assumption.  This certification is also what keeps
  the rule sound on partial terms.  Since \texttt{=} is strict (it requires
  both sides defined), \texttt{3 + f(x) = 1 + f(x) + 2}
  holds only when \texttt{f(x)} is defined; membership entails definedness,
  so demanding \texttt{(IN f(x) $D$)} supplies exactly that guarantee.  A
  generator that is not certified leaves the goal untouched --- the command
  never makes an unsound step.

  \begin{VNBcode}
;;; Directly on a universally quantified goal -- no (di) needed:
(sp (make-wff-from-string
  "forall([a in ZZ, b in ZZ, c in ZZ], a + b + c = c + b + a)"))
(rs)    ; closes in one step

;;; After (di), a compound term certified in context becomes a generator:
(sp (make-wff-from-string
  "forall([x in ZZ], (f(x) in ZZ) implies 3 + f(x) = 1 + f(x) + 2)"))
(di) (di) (rs)    ; di peels the forall then the hypothesis; rs then closes

;;; Coefficient arithmetic:
(sp (make-wff-from-string "forall([x in ZZ], 2*x + 3*x = 5*x)"))
(rs)
  \end{VNBcode}

  The primitive inference is \texttt{pi-ring-simplify!}, in
  \texttt{structure-library/ring-simplify.scm}; it normalizes each side to a
  canonical sum-of-monomials polynomial (\texttt{vnb->poly}) and compares.  The interactive short form is \texttt{(rs)}.

\item \texttt{(cmd-comm-ring-simplify ps)}
  The commutative counterpart of \texttt{(cmd-ring-simplify ps)}.  Where
  \texttt{(rs)} fixes the order of factors within a monomial,
  \texttt{(crs)} treats multiplication as commutative: a monomial is a
  \emph{multiset} of generators, canonicalised as a sorted generator
  list, so \texttt{x*y} and \texttt{y*x} share one normal form and
  commutativity lives in the representation rather than in any rewrite
  step.  Each side reduces to a sum of such monomials over \texttt{ZZ},
  and the goal closes iff the two normal forms coincide.  Use it for any
  identity whose proof needs reordering of factors, which \texttt{(rs)}
  cannot perform.

  The goal may take the same two shapes as for \texttt{(rs)} (a bare
  equality with generators certified by context typings, or a typed
  universally quantified identity), and additionally a \textbf{generic}
  form over an arbitrary ring: \texttt{forall($R$, IS-COMMUTATIVE-RING($R$)
  implies \ldots)} built from \texttt{(ADD $R$)}, \texttt{(MUL $R$)},
  \texttt{(NEG $R$)}, \texttt{(ZERO $R$)}, \texttt{(ONE $R$)}, with the
  generators ranging over the carrier \texttt{(CARR $R$)}.

  Generators are detected and certified exactly as for \texttt{(rs)}: a
  symbol or an opaque compound term, certified by an \texttt{(IN $v$ $D$)}
  typing in the concrete form or an \texttt{(IN $v$ (CARR $R$))} assumption in
  the generic form.  The same partiality caveat applies --- a compound
  generator fires only once certified, which in VNB guarantees it is
  defined.

  \begin{VNBcode}
;;; Reordering of factors -- (rs) cannot close this, (crs) does:
(sp (make-wff-from-string
  "forall([x in ZZ, y in ZZ], x*y*x = x*x*y)"))
(crs)   ; closes in one step

;;; Commutative expansion:
(sp (make-wff-from-string
  "forall([x in ZZ, y in ZZ], (x+y)*(x+y) = x*x + 2*x*y + y*y)"))
(crs)
  \end{VNBcode}

  The procedure is a trusted oracle (it compares canonical normal forms
  rather than mechanising the ring properties), carrying the warrant
  \texttt{comm-ring-simplify} of kind \texttt{well-known}.  The primitive
  inference is \texttt{pi-comm-ring-simplify!} in
  \texttt{comm-ring-simplify.scm}.  The interactive short form is
  \texttt{(crs)}.

\item \texttt{(cmd-reflexivity ps)}
  Close goal \texttt{a = a} when \texttt{a} is demonstrably defined.
  Fails if \texttt{a} may be undefined.

\item \texttt{(cmd-quasi-reflexivity ps)}
  Close goal \texttt{a == a} unconditionally, for any term \texttt{a}.

\item \texttt{(cmd-cut ps $L$)}
  Introduce lemma $L$: subgoal~1 proves $L$; subgoal~2 has $L$ as assumption.

\item \texttt{(cmd-or-intro-left ps)}
  To prove \texttt{p or q}, prove $p$ instead.

\item \texttt{(cmd-or-intro-right ps)}
  To prove \texttt{p or q}, prove $q$ instead.

\item \texttt{(cmd-proof-by-contradiction ps)}
  To prove $p$, assume \texttt{not(p)} and derive \texttt{falsity}.

\item \texttt{(cmd-instantiate ps $F$\ $t$)}
  From \texttt{forall([$x$], $p$)} $= F$ in context, add $p[x/t]$.

\item \texttt{(cmd-exists-witness ps $t$)}
  To prove \texttt{forsome([$x$], $p$)}, prove $p[x/t]$.

\item \texttt{(cmd-weaken ps $f$)}
  Remove formula $f$ from the assumptions.

\item \texttt{(cmd-backchain ps $F$)}
  From \texttt{p implies goal} $= F$ in context, generate subgoal~$p$.

\item \texttt{(cmd-theorem-assumption ps $\mathit{name}$)}
  Add the named theorem (or axiom) to the assumptions.

\item \texttt{(cmd-apply-macete ps $\mathit{name}$)}
  Apply named macete: rewrite the goal using the associated theorem.

\item \begin{sloppypar}\texttt{(cmd-eq-subst ps $F$)}
  Equality substitution (Leibniz).  From an equality \texttt{$s$ = $t$} $= F$
  in the assumptions (either orientation), rewrite every occurrence of $s$
  in the goal to $t$.  Short form \texttt{(subst $F$)}.  This realises the
  schema \texttt{$s$ = $t$} $\wedge$ \texttt{$P[t]$} $\Rightarrow$
  \texttt{$P[s]$} for an arbitrary context $P$; sound for partial equality
  because \texttt{$s$ = $t$} already certifies both sides
  defined.\end{sloppypar}

\item \texttt{(cmd-qed ps $\mathit{name}$)}
  After a proof completes, install its root assertion as a named theorem.
  Errors if the proof is not complete.  This is the only command that
  records a result as \emph{proven}.

\end{itemize}

\subsection{N-ary constructors}
\label{sec:nary-proof-commands}

\begin{itemize}

\item \begin{sloppypar}\texttt{(cmd-cartesian-intro ps)}
  To prove \texttt{[$a_1$,\ldots,$a_n$] in cartesian($A_1$,\ldots,$A_n$)},
  generate $n$ subgoals \texttt{$a_i$ in $A_i$}.\end{sloppypar}

\item \begin{sloppypar}\texttt{(cmd-cartesian-elim ps $F$\ $k$)}
  From \texttt{[$a_1$,\ldots,$a_n$] in cartesian($A_1$,\ldots,$A_n$)} $= F$
  in context, add \texttt{$a_k$ in $A_k$}.\end{sloppypar}

\item \texttt{(cmd-tuples-intro ps)}
  To prove \texttt{[$a_1$,\ldots,$a_n$] in tuples($A$)},
  generate $n$ subgoals \texttt{$a_i$ in $A$}.
  For the empty sequence \texttt{[]} the goal closes immediately (0 subgoals).

\item \begin{sloppypar}\texttt{(cmd-tuples-elim ps $F$\ $k$)}
  From \texttt{[$a_1$,\ldots,$a_n$] in tuples($A$)} $= F$
  in context, add \texttt{$a_k$ in $A$}.\end{sloppypar}

\item \texttt{(cmd-nth-reduce ps)}
  Rewrite every occurrence of \texttt{nth($k$, [$\ldots$])} in the goal to
  the $k$-th element (1-indexed).

\item \begin{sloppypar}\texttt{(cmd-functoid-beta ps)}
  Rewrite every occurrence of
  \texttt{lambda([$x_1$ in $A_1$,\ldots], $\mathit{body}$)($v_1$,\ldots)}
  (or \texttt{lambdoid}) in the goal to
  $\mathit{body}[x_1 := v_1]\cdots$ (capture-avoiding).\end{sloppypar}

\item \texttt{(cmd-union-intro ps $k$)}
  To prove \texttt{$x$ in union($A_1$,\ldots,$A_n$)},
  prove \texttt{$x$ in $A_k$} instead.

\item \texttt{(cmd-union-elim ps $F$)}
  From \texttt{$x$ in union($A_1$,\ldots,$A_n$)} $= F$ in context,
  case-split: one branch per $A_i$ with \texttt{$x$ in $A_i$} assumed.

\item \begin{sloppypar}\texttt{(cmd-intersection-intro ps)}
  To prove \texttt{$x$ in intersection($A_1$,\ldots,$A_n$)},
  generate $n$ subgoals \texttt{$x$ in $A_i$}.\end{sloppypar}

\item \begin{sloppypar}\texttt{(cmd-intersection-elim ps $F$\ $k$)}
  From \texttt{$x$ in intersection($A_1$,\ldots,$A_n$)} $= F$ in context,
  add \texttt{$x$ in $A_k$}.\end{sloppypar}

\end{itemize}

%% =========================================================
\section{Two surfaces, one script}
\label{sec:two-surfaces}

Every proof step exists in two forms.

The \emph{builder} \texttt{(cmd-\textit{name}\ \textit{ps}\ \textit{arg}
\ldots)} (Section~\ref{sec:commands}) takes a proof state explicitly and
\emph{returns} a new one; it has no side effects and never prints.  It is
the form to call from Scheme code that manipulates proof states as values.

The \emph{interactive short form} \texttt{(\textit{name}\ \textit{arg}
\ldots)} (Section~\ref{sec:interactive}) is a thin wrapper around the
builder: it reads the global proof state \texttt{*ps*}, applies the
builder, stores the result back in \texttt{*ps*}, and calls \texttt{(show)}
so the new state is displayed.  This is the form typed in a live session.

\begin{center}
\begin{tabular}{@{}lll@{}}
\toprule
Builder & Short form & \\
\midrule
\texttt{(cmd-direct-inference ps)}        & \texttt{(di)}    & no argument \\
\texttt{(cmd-antecedent-inference ps $f$)}& \texttt{(ai $f$)}& formula argument \\
\texttt{(cmd-backchain ps $F$)}           & \texttt{(bc $F$)}& formula argument \\
\texttt{(cmd-instantiate ps $F$ $t$)}     & \texttt{(inst $F$ $t$)} & formula + term \\
\texttt{(cmd-apply-macete ps $n$)}        & \texttt{(mac $n$)} & theorem name \\
\bottomrule
\end{tabular}
\end{center}

\noindent
The wrapper does one more thing: it \emph{records itself}.  Each short form
appends an entry \texttt{(\textit{name}\ .\ \textit{args})} to the global
accumulator \texttt{*proof-script*}, which \texttt{(sp \ldots)} clears at
the start of a proof.  So there is no separate scripting language to learn:
\textbf{the script is the transcript} --- the very sequence of short forms
you typed.  \texttt{(qed \textit{name})} and \texttt{(save-proof
\textit{name})} store that sequence; \texttt{(replay-proof \textit{name})}
re-runs it (Section~\ref{sec:replay}).

\subsection{Supplying arguments}
\label{sec:arg-collection}

A short form takes one of a few kinds of argument, and each kind is
collected the same way wherever it appears.

\begin{itemize}
\item \textbf{Nothing.}  The structural steps --- \texttt{di}, \texttt{ass},
  \texttt{pbc}, \texttt{rs}, \texttt{rfl}, \texttt{ci}, \texttt{ni},
  \ldots\ --- read the focus goal and need no argument.

\item \textbf{A formula} (\texttt{ai}, \texttt{bc}, \texttt{wk},
  \texttt{cut}, \texttt{ue}, the assumption cited by \texttt{inst} and
  \texttt{mac-h}).  It may be supplied three interchangeable ways:
  \begin{itemize}
  \item a \textbf{surface string}, parsed for you ---
    \texttt{"n in NN and m in NN"};
  \item a \textbf{raw S-expression} ---
    \texttt{'(and (in n nn) (in m nn))};
  \item for a command that cites an existing \emph{assumption}, a
    \textbf{1-based assumption index}: a positive integer $k$ selects the
    $k$-th assumption of the focus sequent, exactly as numbered in the
    display.  \texttt{(ai 1)} stands in for retyping the first assumption.
  \end{itemize}

\item \textbf{A term} (\texttt{ew}, the witness slot of \texttt{inst},
  \texttt{if-true}/\texttt{if-false}).  A string or an S-expression; there
  is no index form, since a term is not an assumption to point at.

\item \textbf{A name} (\texttt{ta}, \texttt{mac}, \texttt{bc*},
  \texttt{qed}).  A symbol naming an installed theorem or macete.

\item \textbf{A small integer} (\texttt{ui $k$}, \texttt{ce $f$ $k$},
  \texttt{focus $n$}).  Selects a branch, a tuple component, or an open
  goal --- not an assumption index.
\end{itemize}

\begin{VNBcode}
(sp (make-wff-from-string "n in NN and m in NN implies n in NN"))
(di)        ; assume the conjunction (now assumption #1)
(ai 1)      ; crack assumption #1 -- the index, not the retyped formula
(ass)
(qed 'and-elim-demo)

(lookup-proof 'and-elim-demo)
; => ((di) (ai 1) (ass))     -- the index is recorded verbatim
\end{VNBcode}

\begin{remark}
The assumption-index form is accepted by the commands that \emph{cite} an
assumption --- \texttt{ai}, \texttt{bc}, \texttt{wk}, \texttt{inst},
\texttt{ce}, \texttt{te}, \texttt{ie}, \texttt{mac-h}.  It is recorded
verbatim in the script, so a replay re-resolves ``assumption~1'' against
the focus sequent it actually reaches --- robust to the formula's exact
text, but dependent on the assumption \emph{order} being the same on
replay.  When that is in doubt, give the formula itself.
\end{remark}

\subsection{Argument collection in a front end: \texttt{bc*}, \texttt{B}, \texttt{B+}}
\label{sec:bc-star-collection}

\texttt{bc*} (Section~\ref{sec:interactive}) matches a theorem's conclusion
against the goal and instantiates whatever that match determines.  A schema
variable the match leaves \emph{undetermined} must be supplied --- as a
binding list \texttt{((\textit{v}\ \textit{val})\ \ldots)} of Scheme
expressions.  The pure query \texttt{(bc*-undetermined '\textit{name})}
reports exactly which variables those are, without touching the proof.

A front end uses that query to collect the values \emph{by name} instead of
making the user write the quoted binding list: it asks for each undetermined
variable's value as an ordinary \emph{term string} (\texttt{"RAN(f)"},
\texttt{"nn"}) and parses it.  In the Emacs Focus workspace this is the
\texttt{B} key --- the argument \emph{builder}.  Its companion key
\texttt{B+} is a saturating closer: it only commits a citation it
can first \emph{prove} will close the focus goal
(\texttt{bc*-can-close?} matches the conclusion \emph{and} checks every
hypothesis is already an assumption), so it never leaves the proof worse
than it found it.

%% =========================================================
\section{Interactive short-form commands}\label{sec:tactics-ref}
\label{sec:interactive}

\texttt{interactive.scm} provides the short-form commands --- one- or two-letter
wrappers around the \texttt{cmd-*} proof procedures --- that you type in a live
proof session, at the bare REPL or through the Emacs workspace.  Each mutates the
global proof state \texttt{*ps*} and calls \texttt{(show)} automatically; the
exceptions, which return a value rather than mutating state, are \texttt{sp},
\texttt{show}, \texttt{pp}, \texttt{fa}, \texttt{fs}, \texttt{vnb-unwind},
\texttt{vnb-wind}, \texttt{parse-string}, and \texttt{make-wff-from-string}.

The reference below is \emph{generated} from the live \texttt{*tactic-help*}
registry (through \texttt{reference/TACTICS.md}), so it lists the current commands
with their signatures and trust \emph{kind} and cannot drift from the prover.  At
the REPL, \texttt{(tactics)} prints the same menu and \texttt{(tactics 'name)} the
full per-command help.

%% AUTO-GENERATED by docs/gen-tactics-ref.py from reference/TACTICS.md -- do not edit.
%% Regenerate:  python3 docs/gen-tactics-ref.py   (after a prover load refreshes TACTICS.md)

\DeclareRobustCommand{\tactkind}[1]{{\footnotesize\textsc{#1}}}

Every tactic is tagged with a \emph{kind} grounded in the \texttt{dg-apply-rule!} tag it emits, so the trusted base is legible: \texttt{rule} is a single primitive kernel inference (the fixed trusted base); \texttt{oracle} is a trusted decision procedure run as a black box; \texttt{composite} only chains kernel rules, adding no new inference; \texttt{meta} performs no deduction (session, search, navigation).  A proof's trust surface is exactly its \texttt{rule} steps plus whichever \texttt{oracle}s and asserted premises it cites.

The same reference is available at the REPL with \texttt{(tactics 'name)} and in the browser file \texttt{reference/TACTICS.md}; all three are generated from one registry, so they never disagree.

An entry's \emph{Uses} line names every kernel operation the command's own code can cause to be recorded, those recorded by the procedures it calls included. It is the bound on what the command can do, whether or not any proof in the library exercises it, and it is measured by reading the source rather than declared alongside the command. One further set is common to every command and is stated once instead of in every entry: the hook that runs after any command which changed the proof, to close the definedness sequents an instantiation left owed. Section~\ref{sec:command-uses} tabulates the same figures for every command at once.

\subsection*{Starting \& finishing a proof}

\par\smallskip{\raggedright\noindent\texttt{(sp goal)}\quad\tactkind{meta}\index{sp@\texttt{sp} (tactic)}\par}\nobreak
\noindent Start a proof of `goal' (a ``string'', a raw S-expr, or a wff); clears the script.\par
\noindent\textit{Uses:} no kernel operation: it records no inference.\par
\noindent\textit{When useful:} starting a new proof\par

Begin a new proof: state what you want to prove.  The goal may be a ``string'' in surface syntax, a raw S-expr, or an already-built wff (e.g. (wff ``...'') or (make-wff '(...)));  sp coerces a string/S-expr for you.  This becomes your one open goal and clears any previous script.  (Technically: set the proof goal.)

\par\smallskip{\raggedright\noindent\texttt{(qed 'name)}\quad\tactkind{meta}\index{qed@\texttt{qed} (tactic)}\par}\nobreak
\noindent Install the finished proof as theorem `name' and save its replayable script.\par
\noindent\textit{Uses:} no kernel operation: it records no inference.\par
\noindent\textit{When useful:} no open goals remain -- record the result\par

Finish: once there are no open goals, this records the result as a named theorem you can cite later, and saves the replayable script.  (Technically: install-theorem! plus script capture; reports the asserted facts the proof still rests on, `proven modulo \{...\}'.)

\par\smallskip{\raggedright\noindent\texttt{(save-proof 'name)}\quad\tactkind{meta}\index{save-proof@\texttt{save-proof} (tactic)}\par}\nobreak
\noindent Save the current script under a name without finishing.\par
\noindent\textit{Uses:} no kernel operation: it records no inference.\par

Save the current (possibly unfinished) proof script under a name, to resume or replay later.  (Technically: snapshot the script without installing a theorem.)

\par\smallskip{\raggedright\noindent\texttt{(replay-proof 'name [subst])}\quad\tactkind{meta}\index{replay-proof@\texttt{replay-proof} (tactic)}\par}\nobreak
\noindent Re-run a saved script on the current goal, optionally renaming free vars.\par
\noindent\textit{Uses:} every one of the 64 kernel operations.\par

The optional subst is an alist ((old . new) ...) applied to every command argument before replay -- this is how one proof is reused at fresh eigenvariables.


\subsection*{Decomposing the goal}

\par\smallskip{\raggedright\noindent\texttt{(di)}\quad\tactkind{rule}\index{di@\texttt{di} (tactic)}\par}\nobreak
\noindent Direct inference: split an AND goal, move an IMPLIES antecedent into the assumptions, or introduce a fresh eigenvariable for a leading FORALL.\par
\noindent\textit{Uses:} \texttt{and-intro} \texttt{forall-intro} \texttt{iff-intro} \texttt{implies-intro} \texttt{not-intro}.\par
\noindent\textit{When useful:} the goal head is AND / IMPLIES / FORALL -- decompose before deciding\par

Break the goal into its pieces.  A conjunction `A and B' splits into two goals, A and B.  An implication `P implies Q' assumes P (P becomes a hypothesis you may use) and leaves you to prove Q.  A universal `for all x, ...' fixes an arbitrary x (a new constant standing for `any x') and asks for the statement at that x.  Apply it repeatedly until the goal is a single atomic statement.  (Technically: the goal-side introduction rules for AND / IMPLIES / FORALL; the constant introduced for a FORALL is an eigenvariable.  Pairs with mac, which unfolds a definition in the goal.)

\par\smallskip{\raggedright\noindent\texttt{(pbc)}\quad\tactkind{rule}\index{pbc@\texttt{pbc} (tactic)}\par}\nobreak
\noindent Proof by contradiction: assume the goal's negation, prove FALSITY.\par
\noindent\textit{Uses:} \texttt{proof-by-contradiction}.\par
\noindent\textit{When useful:} a direct argument stalls and assuming the negation gives something concrete to work with\par

Argue by contradiction: assume the goal is false and derive an absurdity.  (Technically: classical reductio -- adds the negation of the goal as a hypothesis and changes the goal to FALSITY.)

\par\smallskip{\raggedright\noindent\texttt{(oi-l)}\quad\tactkind{rule}\index{oi-l@\texttt{oi-l} (tactic)}\par}\nobreak
\noindent OR-intro left: reduce an (OR a b) goal to a.\par
\noindent\textit{Uses:} \texttt{or-intro-left}.\par
\noindent\textit{When useful:} the goal is (OR a b) and the LEFT alternative is provable\par

Prove a disjunction `A or B' by proving the LEFT alternative, A.  (Technically: or-introduction, left.)

\par\smallskip{\raggedright\noindent\texttt{(oi-r)}\quad\tactkind{rule}\index{oi-r@\texttt{oi-r} (tactic)}\par}\nobreak
\noindent OR-intro right: reduce an (OR a b) goal to b.\par
\noindent\textit{Uses:} \texttt{or-intro-right}.\par
\noindent\textit{When useful:} the goal is (OR a b) and the RIGHT alternative is provable\par

Prove a disjunction `A or B' by proving the RIGHT alternative, B.  (Technically: or-introduction, right.)

\par\smallskip{\raggedright\noindent\texttt{(ew term)}\quad\tactkind{rule}\index{ew@\texttt{ew} (tactic)}\par}\nobreak
\noindent Existential witness: discharge a FORSOME goal by supplying the witness term.\par
\noindent\textit{Uses:} \texttt{forsome-intro}.\par
\noindent\textit{When useful:} the goal is FORSOME and you have a witness term in mind\par

Prove `there exists an x with property P' by exhibiting a specific witness: you supply the term, and the goal becomes `P holds of that term'.  The everyday `take x = ...' step.  (Technically: existential introduction for a FORSOME goal.)

\par\smallskip{\raggedright\noindent\texttt{(ci)}\quad\tactkind{rule}\index{ci@\texttt{ci} (tactic)}\par}\nobreak
\noindent Cartesian intro: prove a CARTESIAN-product membership component-wise.\par
\noindent\textit{Uses:} \texttt{cartesian-intro}.\par
\noindent\textit{When useful:} the goal is membership in a CARTESIAN product\par

Prove that a pair (or tuple) belongs to a Cartesian product A x B by proving each coordinate lies in its factor.  (Technically: Cartesian-membership introduction, component-wise.)

\par\smallskip{\raggedright\noindent\texttt{(ti)}\quad\tactkind{rule}\index{ti@\texttt{ti} (tactic)}\par}\nobreak
\noindent Tuple intro: prove a tuple/LIST membership component-wise.\par
\noindent\textit{Uses:} \texttt{tuples-intro}.\par
\noindent\textit{When useful:} the goal is membership in a set of tuples / LIST\par

Prove that a list belongs to the set of tuples over A by proving each entry lies in A.  (Technically: tuple-membership introduction, component-wise.)

\par\smallskip{\raggedright\noindent\texttt{(ii)}\quad\tactkind{rule}\index{ii@\texttt{ii} (tactic)}\par}\nobreak
\noindent Intersection intro: prove (IN x (INTERSECTION ...)) for each branch.\par
\noindent\textit{Uses:} \texttt{intersection-intro}.\par
\noindent\textit{When useful:} the goal is membership in an INTERSECTION\par

Prove that x lies in an intersection of sets by proving it lies in each of them.  (Technically: intersection-membership introduction.)

\par\smallskip{\raggedright\noindent\texttt{(ui k)}\quad\tactkind{rule}\index{ui@\texttt{ui} (tactic)}\par}\nobreak
\noindent Union intro: reduce a goal (IN x (UNION ...)) to membership of x in the k-th set (1-based).\par
\noindent\textit{Uses:} \texttt{union-intro}.\par
\noindent\textit{When useful:} the goal is membership in a UNION and you know which branch\par

Prove that x lies in a union of sets by proving it lies in the k-th one (you choose which).  (Technically: union-membership introduction at branch k.)

\par\smallskip{\raggedright\noindent\texttt{(ni)}\quad\tactkind{rule}\index{ni@\texttt{ni} (tactic)}\par}\nobreak
\noindent Natural-number induction on the goal's leading FORALL over NN.\par
\noindent\textit{Uses:} \texttt{nn-induction}.\par
\noindent\textit{When useful:} the goal is `forall n in NN, P(n)'\par

Prove `for all natural numbers n, P(n)' by induction: show P(0), and that P(n) implies P(n+1).  (Technically: natural-number induction on the goal's leading FORALL over NN.)


\subsection*{Using a hypothesis}

\par\smallskip{\raggedright\noindent\texttt{(ai hyp)}\quad\tactkind{rule}\index{ai@\texttt{ai} (tactic)}\par}\nobreak
\noindent Antecedent inference: decompose a cited assumption -- AND-split, OR-into-cases, or FORSOME-elimination to a fresh eigenvariable.\par
\noindent\textit{Uses:} \texttt{and-elim} \texttt{forsome-elim} \texttt{iff-elim} \texttt{not-elim} \texttt{or-elim}.\par
\noindent\textit{When useful:} a hypothesis is an AND / OR / FORSOME to break apart\par

Break a hypothesis into its pieces -- the mirror image of di, but on the assumptions instead of the goal.  From a hypothesis `A and B' you get both A and B; from `A or B' you split into two cases (prove the goal in each); from `there exists x, P(x)' you get a fresh name for such an x together with P(x).  Cite the hypothesis by its number in the display, its formula, or a ``string''.  (Technically: the hypothesis-side elimination rules for AND / OR / FORSOME; the fresh name is an eigenvariable.)

\par\smallskip{\raggedright\noindent\texttt{(ass)}\quad\tactkind{rule}\index{ass@\texttt{ass} (tactic)}\par}\nobreak
\noindent Close the goal by an assumption alpha-equivalent to it.\par
\noindent\textit{Uses:} \texttt{assumption}.\par
\noindent\textit{When useful:} the goal already appears (up to bound-var renaming) among the hypotheses\par

Finish the goal because it is already one of your hypotheses -- the `that is exactly what we assumed' step.  Renaming of dummy/bound variables doesn't matter (`there exists x, P(x)' closes `there exists y, P(y)').  (Technically: the assumption rule, matching up to alpha-equivalence.)

\par\smallskip{\raggedright\noindent\texttt{(inst forall-hyp term)}\quad\tactkind{rule}\index{inst@\texttt{inst} (tactic)}\par}\nobreak
\noindent Instantiate a universally-quantified assumption at `term', adding the instance to context.\par
\noindent\textit{Uses:} \texttt{forall-elim}.\par
\noindent\textit{When useful:} you have a forall-hypothesis and a specific value to use it at\par

Use a `for all x, ...' hypothesis at a particular value: you supply the term, and the statement with x replaced by that term is added to your hypotheses.  (Technically: universal instantiation of an assumption.)

\par\smallskip{\raggedright\noindent\texttt{(inst+ forall-hyp term)}\quad\tactkind{composite}\index{inst+@\texttt{inst+} (tactic)}\par}\nobreak
\noindent Instantiate an in-context universal at `term', then forward-detach any guards whose antecedents are in context.\par
\noindent\textit{Uses:} \texttt{arith-forsome} \texttt{arith-ground} \texttt{arith-simplify} \texttt{cut} \texttt{detach} \texttt{forall-elim}.\par
\noindent\textit{When useful:} a GUARDED forall-hyp whose guards are dischargeable from context (e.g. the metric laws post-grind)\par

inst followed by detach: from `forall x. (x in S) => P(x)' and `t in S' already known, this lands `P(t)' directly (peeling nested guards level by level), instead of leaving the guarded implication for you to detach by hand.  The hypothesis-side analogue of `fact' (which assembles a THEOREM); this assembles an in-context UNIVERSAL.  scout's inst lane emits these to close witness-needing goals like the metric laws.  (Technically: pi-instantiate! then pi-detach! while the consequent stays a guard with an in-context antecedent.)

\par\smallskip{\raggedright\noindent\texttt{(mp)}\quad\tactkind{composite}\index{mp@\texttt{mp} (tactic)}\par}\nobreak
\noindent Close a goal that is an INSTANCE of a universal already in your context.  No arguments: the term is DERIVED by matching, not supplied.\par
\noindent\textit{Uses:} \texttt{arith-forsome} \texttt{arith-ground} \texttt{arith-simplify} \texttt{assumption} \texttt{cut} \texttt{detach} \texttt{forall-elim}.\par
\noindent\textit{When useful:} the goal IS an instance of a universal you already have -- no term to supply, it is derived by matching\par

The syllogism, with nothing to fill in.  From `forall([thing in human], thing in mortal)' together with `socrates in human', it closes `socrates in mortal'.  Where inst+ makes you name the term, mp works it out: it matches the universal's conclusion against the goal, which either determines the bound variable or fails outright, and then checks that the universal's guard at that term is something you already have.  Nothing is searched for and nothing is ranked -- if the goal is an instance there is exactly one term it can be.  It DECLINES, naming what it found, when no universal in the context has the goal as an instance, and when more than one does: a closer that silently picks between two universals is one you cannot predict.  One bound variable in this version; for `forall([a in A, b in B], ...)' use inst+.  (Technically: composite -- inst+ at the derived term, then ass.  It adds no kernel rule and no debt; the bill is what typing the two steps by hand would bill, which is nothing.)

\par\smallskip{\raggedright\noindent\texttt{(grind-and-mp)}\quad\tactkind{composite}\index{grind-and-mp@\texttt{grind-and-mp} (tactic)}\par}\nobreak
\noindent Tidy the goal with grind, then close it by modus ponens -- the one-button form of the syllogism, for a sentence typed exactly as it reads.\par
\noindent\textit{Uses:} \texttt{and-elim} \texttt{and-intro} \texttt{arith-forsome} \texttt{arith-ground} \texttt{arith-simplify} \texttt{assumption} \texttt{cut} \texttt{detach} \texttt{forall-elim} \texttt{forall-intro} \texttt{forsome-elim} \texttt{iff-elim} \texttt{iff-intro} \texttt{implies-intro} \texttt{macete-hyp} \texttt{not-elim} \texttt{not-intro} \texttt{or-elim}.\par
\noindent\textit{When useful:} a sentence typed as it reads -- `all men are mortal and socrates is a man implies socrates is mortal' -- closing in one move\par

Runs the two steps a syllogism actually takes on a sentence you have just typed.  `forall([thing in human], thing in mortal) and socrates in human implies socrates in mortal' is not yet a syllogism: it is an implication whose antecedent is a conjunction.  grind does the BOOKKEEPING -- split the AND, move the antecedents into the assumptions -- which leaves the goal `socrates in mortal' with the universal and `socrates in human' as hypotheses; mp then does the LOGIC (universal instantiation, then modus ponens).  They stay separate tactics because they are different kinds of step, and mp on its own is what to reach for once the goal is already tidy.  The recorded script keeps both steps, so the page reads as the two moves rather than as this wrapper.  (Technically: composite -- grind, quietly since on a tidy goal it legitimately declines, then mp.  Adds no kernel rule and no debt.)

\par\smallskip{\raggedright\noindent\texttt{(detach! impl)}\quad\tactkind{rule}\index{detach"!@\texttt{detach"!} (tactic)}\par}\nobreak
\noindent Forward modus ponens: from an in-context (IMPLIES A B) whose A is also in context, leave B in context.\par
\noindent\textit{Uses:} \texttt{detach}.\par
\noindent\textit{When useful:} you have both P and (P => Q) in context and want Q\par

If you have both `P' and `P implies Q' among your hypotheses, this adds `Q' to them.  It grows what you KNOW (the hypotheses) rather than changing the goal -- the forward counterpart of bc.  (Technically: forward modus ponens, the kernel rule pi-detach!; sound since P and P=>Q give Q.)

\par\smallskip{\raggedright\noindent\texttt{(fact 'thm term ...)}\quad\tactkind{composite}\index{fact@\texttt{fact} (tactic)}\par}\nobreak
\noindent Forward APPLICATION of a theorem: bring it in, instantiate its leading universals with the terms, and auto-detach every antecedent already in context, landing the consequent as a hypothesis.\par
\noindent\textit{Uses:} \texttt{arith-forsome} \texttt{arith-ground} \texttt{arith-simplify} \texttt{cut} \texttt{detach} \texttt{forall-elim} \texttt{theorem-assumption}.\par
\noindent\textit{When useful:} a library law `forall x. H(x) => P(x)' whose P you want landed as a hypothesis\par

Handles interleaved forall/implies (e.g. forall s. IS-X(s) => forall a. a in CARR(s) => P): consumes one term per FORALL, detaches each IMPLIES whose antecedent is in context.  The forward-assembly workhorse -- a law `forall x. H(x) => P(x)' becomes the usable fact P in one call, instead of ta + inst* + cut/backchain.  See theorem-library/module-zero-act.scm.

\par\smallskip{\raggedright\noindent\texttt{(ce hyp k)}\quad\tactkind{rule}\index{ce@\texttt{ce} (tactic)}\par}\nobreak
\noindent Cartesian elim: project the k-th component out of a CARTESIAN-membership assumption.\par
\noindent\textit{Uses:} \texttt{cartesian-elim}.\par
\noindent\textit{When useful:} a hypothesis is CARTESIAN-product membership\par

From a hypothesis that a tuple lies in a Cartesian product, extract that its k-th coordinate lies in the k-th factor.  (Technically: Cartesian-membership elimination, k-th projection.)

\par\smallskip{\raggedright\noindent\texttt{(te hyp k)}\quad\tactkind{rule}\index{te@\texttt{te} (tactic)}\par}\nobreak
\noindent Tuple elim: project the k-th component out of a tuple-membership assumption.\par
\noindent\textit{Uses:} \texttt{tuples-elim}.\par
\noindent\textit{When useful:} a hypothesis is tuple membership\par

From a hypothesis that something is a tuple over A, extract that its k-th entry lies in A.  (Technically: tuple-membership elimination, k-th projection.)

\par\smallskip{\raggedright\noindent\texttt{(ie hyp k)}\quad\tactkind{rule}\index{ie@\texttt{ie} (tactic)}\par}\nobreak
\noindent Intersection elim: extract the k-th branch of an INTERSECTION-membership assumption.\par
\noindent\textit{Uses:} \texttt{intersection-elim}.\par
\noindent\textit{When useful:} a hypothesis is INTERSECTION membership\par

From a hypothesis that x lies in an intersection, extract that x lies in the k-th set.  (Technically: intersection-membership elimination.)

\par\smallskip{\raggedright\noindent\texttt{(ue hyp)}\quad\tactkind{rule}\index{ue@\texttt{ue} (tactic)}\par}\nobreak
\noindent Union elim: split a cited (IN x (UNION ...)) membership assumption into one subgoal per set -- union-side case analysis, the dual of `ui'.\par
\noindent\textit{Uses:} \texttt{union-elim}.\par
\noindent\textit{When useful:} a hypothesis is UNION membership -- case-split on it\par

From a hypothesis that x lies in a union, split into cases -- one for each set x might belong to -- and prove the goal in each.  (Technically: union-membership elimination, case analysis; the dual of ui.)

\par\smallskip{\raggedright\noindent\texttt{(cut formula)}\quad\tactkind{rule}\index{cut@\texttt{cut} (tactic)}\par}\nobreak
\noindent Cut: prove `formula' as a side subgoal, then continue the main goal with `formula' added as an assumption (Gentzen cut).\par
\noindent\textit{Uses:} \texttt{cut}.\par
\noindent\textit{When useful:} you want to prove a lemma on the spot, then use it\par

Introduce a lemma you prove on the spot.  You state a formula; it becomes a side goal to prove, and on the main line you may then use it as a hypothesis.  The standard `we first show that ..., and now using it ...' move.  (Technically: Gentzen's cut -- nothing is left assumed-but-unproved, the side goal discharges it.)

\par\smallskip{\raggedright\noindent\texttt{(wk hyp)}\quad\tactkind{rule}\index{wk@\texttt{wk} (tactic)}\par}\nobreak
\noindent Weaken: drop a cited assumption from the context to tidy the hypothesis list.  `hyp' may be a formula, a ``string'', or a 1-based assumption index.\par
\noindent\textit{Uses:} \texttt{weakening}.\par
\noindent\textit{When useful:} the hypothesis list is cluttered with something no longer needed\par

Discard a hypothesis you no longer need, to keep the assumption list readable.  (Technically: weakening -- removing a hypothesis is always sound.)

\par\smallskip{\raggedright\noindent\texttt{(keep hyp ...)}\index{keep@\texttt{keep} (tactic)}\par}\nobreak
\noindent Keep only the cited assumptions and drop every other one, as ONE recorded step.  Each `hyp' may be a formula, a ``string'', or a 1-based assumption index; they are recorded as formulas.\par
\noindent\textit{When useful:} the context is deep and only a few hypotheses matter for the next step (prop, ineq)\par

Prune the context down to what the next step needs -- what a driver does before `prop' or `ineq' in a deep context.  (Technically: one weakening per dropped hypothesis, on the same leaf, recorded once; the driver kit's dk-only! is this command.)


\subsection*{Rewriting}

\par\smallskip{\raggedright\noindent\texttt{(mac 'name)}\quad\tactkind{rule}\index{mac@\texttt{mac} (tactic)}\par}\nobreak
\noindent Rewrite the GOAL with an equivalence macete (unfold a definition, apply an iff/=/== law).  Fires only where the macete's side-conditions already hold in context.\par
\noindent\textit{Uses:} \texttt{cartesian-decompose} \texttt{macete} \texttt{tuple-equality-decompose}.\par
\noindent\textit{When useful:} the GOAL has a defined predicate to unfold / an identity to apply, side-conditions already met\par

Rewrite the goal using a definition or a known equivalence/equality (a `macete').  For instance, replace a defined predicate by what it stands for, or apply an identity.  It fires only where the law's side-conditions already hold in your hypotheses; where they don't, it leaves that spot untouched and looks deeper inside.  (Technically: goal-side rewriting by an equivalence/equality macete; all-or-nothing -- it does NOT spawn unmet side-conditions as goals.  Its hypothesis-side cousin is mac-h; the minor-premise-spawning variant is macm.)

\par\smallskip{\raggedright\noindent\texttt{(slot 'acc)}\index{slot@\texttt{slot} (tactic)}\par}\nobreak
\noindent Reduce a structure ACCESSOR to its projection: (CARR s) becomes (NTH 1 s).  The one door for accessor reductions -- use it, never mac, on an accessor name.\par
\noindent\textit{Uses:} \texttt{cartesian-decompose} \texttt{macete} \texttt{tuple-equality-decompose}.\par

Replace an accessor by the tuple position it stands for: (CARR s) is slot 1, so it becomes (NTH 1 s).  Pair it with nth-r to compute on a concrete structure -- (MUL ZZ-RING) reduces to bintimes.  It refuses anything that is not an accessor.  (Technically: fires the accessor's macete, which def-structure installs at declaration.  It exists so that ALL accessor reductions go through ONE procedure: the reduction is currently global and unconditional -- index k for every argument -- and making it structure-relative later, guarded by IS-X(s), should change this door and not its callers.  accessor-callsite-audit fails the suite if any file fires an accessor macete by name.)

\par\smallskip{\raggedright\noindent\texttt{(slot-h 'acc hyp)}\index{slot-h@\texttt{slot-h} (tactic)}\par}\nobreak
\noindent Reduce a structure ACCESSOR to its projection inside a cited HYPOTHESIS -- `slot', hypothesis-side.\par
\noindent\textit{Uses:} \texttt{macete-hyp}.\par

The hypothesis-side door of slot: replace (PTS RR-MS) by the tuple position it stands for inside an assumption you cite, instead of in the goal.  Before it existed, a proof needing that reduction in a hypothesis reached for mac-h with the accessor macete's own name (rr-ms@pts) -- firing an accessor macete BY NAME, which is what the accessor pin (accessor-callsite-audit) exists to prevent, and which kept the suite red.  It refuses anything that is not an accessor, and it ERRORS rather than guess when the hypothesis mentions two different instances of the same accessor: after the first rewrite the assumption is a different formula, so a silent half-rewrite would be the failure mode.  (Technically: slot's projection-macete lookup run over the cited assumption, then cmd-apply-macete-to-assumption.)

\par\smallskip{\raggedright\noindent\texttt{(macm 'name)}\quad\tactkind{rule}\index{macm@\texttt{macm} (tactic)}\par}\nobreak
\noindent Like mac, but a CONDITIONAL macete fires even when its side-conditions are not yet in context: each unmet condition is left as a new subgoal.\par
\noindent\textit{Uses:} \texttt{cartesian-decompose} \texttt{macete} \texttt{tuple-equality-decompose}.\par
\noindent\textit{When useful:} the GOAL has a CONDITIONAL identity/definition to apply whose side-conditions you are willing to leave as subgoals\par

The minor-premises variant of mac.  mac only rewrites where the law's side-conditions ALREADY hold; macm rewrites regardless, spawning each unmet condition as a fresh goal to discharge -- the goal-side mirror of what mac-h already does on hypotheses.  This is what makes conditional finite-sum lemmas (finsum-*, sum-*, ...) usable as goal rewrites instead of only forward via fact.  (Technically: goal-side rewriting by a conditional macete; it emits the SAME single `macete' kernel rule as mac -- the unmet conditions become that one step's minor-premise subgoals.  It introduces NO new inference rule.)

\par\smallskip{\raggedright\noindent\texttt{(mac-h 'name hyp)}\quad\tactkind{rule}\index{mac-h@\texttt{mac-h} (tactic)}\par}\nobreak
\noindent Rewrite a cited ASSUMPTION in place with an equivalence macete; any side-condition not already in context is spawned as a new subgoal.\par
\noindent\textit{Uses:} \texttt{macete-hyp}.\par
\noindent\textit{When useful:} a HYPOTHESIS has a definition to unfold / an equivalence to apply\par

Like mac, but rewrites inside one of your HYPOTHESES instead of the goal -- replacing it by an equivalent statement (e.g. unfolding a definition in a hypothesis).  If the rewrite law has a side-condition you haven't established, that side-condition becomes a new goal to prove.  (Technically: hypothesis-side rewriting by Leibniz substitution of equivalents; sound because the macete is a genuine equivalence under its side-conditions, which are spawned as goals.)

\par\smallskip{\raggedright\noindent\texttt{(mac-h*)}\quad\tactkind{composite}\index{mac-h*@\texttt{mac-h*} (tactic)}\par}\nobreak
\noindent Saturating mac-h: repeatedly unfold every defined predicate in the hypotheses and split the conjunctions they expose, until nothing is left folded.  One step in the trace instead of a dozen.\par
\noindent\textit{Uses:} \texttt{and-elim} \texttt{forsome-elim} \texttt{iff-elim} \texttt{macete-hyp} \texttt{not-elim} \texttt{or-elim}.\par
\noindent\textit{When useful:} several hypotheses are folded definitions / conjunctions to open at once\par

The hands-free version of mac-h.  Instead of naming each definition and each conjunction by hand -- (mac-h 'IS-METRIC-SPACE A1), split, (mac-h 'is-metric A2), split, ... -- (mac-h*) keeps unfolding any assumption whose head is a defined predicate and splitting any AND assumption, restarting until a full pass changes nothing.  It records as a SINGLE step, so the proof (and its PDF) is far shorter.  It adds no kernel rule: every unfold is the same kernel-checked mac-h, every split the same ai -- just driven to a fixpoint.  Ask (what-now) to see which hypotheses it would touch first.

\par\smallskip{\raggedright\noindent\texttt{(grind)}\quad\tactkind{composite}\index{grind@\texttt{grind} (tactic)}\par}\nobreak
\noindent Saturate the no-choice moves at the focus: decompose the goal connective (di) and break open the hypotheses (mac-h*), repeatedly, until neither fires.  The whole introduce-everything prefix in one step.\par
\noindent\textit{Uses:} \texttt{and-elim} \texttt{and-intro} \texttt{forall-intro} \texttt{forsome-elim} \texttt{iff-elim} \texttt{iff-intro} \texttt{implies-intro} \texttt{macete-hyp} \texttt{not-elim} \texttt{not-intro} \texttt{or-elim}.\par
\noindent\textit{When useful:} right after sp, or any time the focus has connective/definitional structure -- normalize first\par

The deterministic normalizer.  It keeps applying di -- which strips a leading FORALL/IMPLIES/AND, introducing the bound variables and moving antecedents into your hypotheses -- and mac-h*, which unfolds defined predicates and splits conjunctions in the hypotheses, until the goal head is no longer a connective and no hypothesis is foldable.  None of these moves involves CHOOSING a lemma, so there is never anything to reconsider: grind just puts the focus into normal form (e.g. a metric-law goal becomes `(d s)(x,y) = (d s)(y,x)' with the metric's defining properties sitting unfolded in context).  It collapses the di...di mac-h* prefix that bloats proofs into a single step, and adds no kernel rule.

\par\smallskip{\raggedright\noindent\texttt{(scout [depth [branch [nodes]]])}\quad\tactkind{meta}\index{scout@\texttt{scout} (tactic)}\par}\nobreak
\noindent Speculatively search to a bounded depth across INDEPENDENT scratch branches; RETURNS a nested list (number-of-branches (d b) goal best-partials closing-branches) rather than printing.  Never touches the live proof.\par
\noindent\textit{Uses:} every one of the 64 kernel operations.\par
\noindent\textit{When useful:} you are stuck and want the machine to TRY move-sequences (returns data)\par

The copilot's deep lane.  Where (what-now)/(tt) suggest one move, scout actually TRIES sequences: it clones the focus into a fresh deduction graph per branch (so backtracking is free -- a dead branch is just discarded), and does a BEST-FIRST search (frontier ordered by open subgoals left, ties broken deeper-first so it dives toward a closure) whose alphabet is (grind), the closers (ass/rfl/crs/arith), the top rewrite-index (mac) rules, the parameterless backchain (bc*) lemmas, and the INST LANE -- (inst+ <universal hyp> <context-typed term>), instantiating a `forall' hypothesis at a term the context already types and detaching the guard.  It RETURNS the result as data -- (examined-count (depth branch) goal best-partials closing-branches), where best-partials is ((open-goals-left (form ...)) ...) most-reduced first and closing-branches is ((form ...) ...) shortest first -- so you can pick it apart programmatically; use (scout-show ...) for the readable REPL report.  Every move is the real, kernel-checked tactic, so a closing branch is a genuine proof when you adopt it with (scout-run k).  A CITATION GUARD suppresses circular closures: a branch that discharges the goal by citing a library theorem that is the goal -- DIRECTLY (statement alpha-equal to the goal, `P proved by P') or TRANSITIVELY (the cited theorem was itself proven using the goal, i.e. a name alpha-equal to the goal is in its proof-debt closure -- the `compact=>complete by compact=>bongo + bongo=>complete' trap) -- is dropped from closing-branches, the theorem name recorded in *scout-citations*, and scout-show notes `the goal is already theorem X'.  (Loops built purely from independent ASSERTED facts have no live dependency edge to detect; that is warrant/debt hygiene, not scout's to check.)  Bounds: depth (default 6), per-node fan-out branch (default 3), total nodes (default 600).  El cheapo: it does not reason about which move is wise (only the witness terms are guessed, from what the context already types -- a bad guess just dies on the clone); it brute-forces a small tree and you eyeball the survivors.  The inst lane closes the metric laws past grind; goals that need a witness scout can't type (a fresh existential, a constructed term) still surface under best-partials.

\par\smallskip{\raggedright\noindent\texttt{(scout-show [depth [branch [nodes]]])}\quad\tactkind{meta}\index{scout-show@\texttt{scout-show} (tactic)}\par}\nobreak
\noindent Like (scout), but PRINTS the human-readable report (examined count, goal, closing branches or best partials) to the REPL.  Returns the same nested list (scout) does.\par
\noindent\textit{Uses:} every one of the 64 kernel operations.\par
\noindent\textit{When useful:} same as scout but you want the readable report at the REPL\par

The eyeball version of scout: same search, same return value, but it also prints the numbered CLOSING branches (or, when none close, the best partials with how many goals each leaves open).  Use scout-show interactively, (scout) when you want to consume the result as data.

\par\smallskip{\raggedright\noindent\texttt{(scout-run k)}\quad\tactkind{composite}\index{scout-run@\texttt{scout-run} (tactic)}\par}\nobreak
\noindent Adopt closing branch k from the last (scout)/(scout-show) onto the live proof, running its tactics for real (they record and display normally).\par
\noindent\textit{Uses:} no kernel operation: it records no inference.\par
\noindent\textit{When useful:} scout found a closing branch you want to adopt onto the live proof\par

After scout finds CLOSING branches [1], [2], ..., (scout-run k) replays branch k's tactic forms through the real tactics on your live *ps*, from the same focus scout cloned -- so the proof advances and the steps are recorded for the script / PDF exactly as if you had typed them.  Works after either (scout) or (scout-show); both stash the closing branches.

\par\smallskip{\raggedright\noindent\texttt{(prep 'ineq)}\index{prep@\texttt{prep} (tactic)}\par}\nobreak
\noindent Why will a tactic not fire here, and what must be done first?  Runs the tactic's OWN preconditions one at a time, marks each ok / PREP / STOP, names the library lemma that repairs each unmet one, and prints a PLAN it has checked by running it on a scratch clone.  READ-ONLY: the live proof is never touched.  (prep 'ineq) is the only method so far.\par
\noindent\textit{Uses:} no kernel operation: it records no inference.\par

A tactic reports failure as a boolean, so every unmet precondition comes out as the same \#f and the same warning -- (ineq) says `goal not a linear-RR consequence of the named assumptions' whether your goal merely needs a di, or an atom needs a typing fact the library already proves, or the goal is simply false.  prep runs the same predicates separately and tells you WHICH failed.  For ineq the obligations are: the goal must BE an order relation (repair: di, counted by measuring, since one di consumes a typed FORALL and its guard together); every atom must be certified in RR by a LITERAL (IN t RR) scan -- (IN k NN) does not count, the oracle does no subtype reasoning (repair: a coercion lemma, e.g. nn-in-rr); some assumption must be order-shaped, since a fact like not(k=0) is invisible to Farkas (repair: a bridge lemma, e.g. nn-pos-of-nonzero); and the goal must actually follow, which only Fourier-Motzkin decides -- a STOP there means no prepping will help, which is the useful answer.  The repair search is by SHAPE over the whole library and is verified by SPECULATION: a candidate counts only if ineq then fires, not merely because it landed something order-shaped.  Every repair that works is reported, not just the first -- they come out alphabetical, which is no order of merit, and if the goal is itself a library theorem then citing IT is one of the closers.  Worked case: |- forall k in NN. \textasciitilde{}(k=0) => k < 2k.  prep returns (di) (di) (fact 'nn-in-rr 'k) (fact 'nn-pos-of-nonzero 'k) (ineq 4) -- six steps, where the hand-built cut/crs route through k = k+0 < k+k = 2k took thirteen and billed two extra transitivity lemmas.  A tactic joins the table with (prep-method! 'FUBA proc); crs and ass are the obvious next two.

\par\smallskip{\raggedright\noindent\texttt{(subst '(= s t))  |  (subst '(== s t))}\quad\tactkind{rule}\index{subst@\texttt{subst} (tactic)}\par}\nobreak
\noindent Rewrite s -> t throughout the goal, using an equation s = t -- or a quasi-equation s == t -- that is in context, in EITHER orientation (Leibniz substitution).  Reaches every position, operator (head) position included; a binder whose variable the equation mentions is not entered.\par
\noindent\textit{Uses:} \texttt{eq-subst}.\par
\noindent\textit{When useful:} you have an equation s = t (or a quasi-equation s == t) in context and want to rewrite the goal by it -- wherever the term occurs, operator position included\par

Use an equation among your hypotheses to replace s by t everywhere in the goal.  BOTH equalities license it: `=' is VNB's partial equality, and `s = t' already entails s = s and t = t, so t is defined wherever s was and the rewrite never replaces a defined term by an undefined one; `==' is quasi-equality (same definedness, equal where defined), which is a congruence and so substitutes exactly as `=' does -- which you need, since the partial-op recursion and bridge facts are stated with `=='.  The head you pass need not match the head in context, and neither need the orientation: all four of (= s t), (= t s), (== s t), (== t s) are searched for, and the goal is always rewritten s -> t.  LIMIT: the rewrite walk reaches argument positions only, so it is a silent no-op on a term in OPERATOR position -- (subst '(= (VADD md) ...)) will not touch the goal ((VADD md) x y).  Structure accessors are almost always in operator position; use the equation as a macete there (mac on the goal, mac-h on a hypothesis).  (Technically: Leibniz substitution from an in-context equality or quasi-equality.)

\par\smallskip{\raggedright\noindent\texttt{(beta)}\quad\tactkind{rule}\index{beta@\texttt{beta} (tactic)}\par}\nobreak
\noindent Beta-reduce a functoid application in the goal.\par
\noindent\textit{Uses:} \texttt{functoid-beta}.\par
\noindent\textit{When useful:} the goal has a functoid applied to an argument\par

Simplify a function-expression applied to an argument by substituting the argument into its body.  (Technically: beta-reduction of a functoid application in the goal.)

\par\smallskip{\raggedright\noindent\texttt{(nth-r)}\quad\tactkind{rule}\index{nth-r@\texttt{nth-r} (tactic)}\par}\nobreak
\noindent Reduce an NTH applied to a literal LIST in the goal.\par
\noindent\textit{Uses:} \texttt{nth-reduce}.\par
\noindent\textit{When useful:} the goal has NTH of a literal LIST\par

Simplify `the k-th entry of an explicit list [a1, a2, ...]' to that entry.  (Technically: reduce NTH applied to a literal LIST.)

\par\smallskip{\raggedright\noindent\texttt{(len-r)}\quad\tactkind{rule}\index{len-r@\texttt{len-r} (tactic)}\par}\nobreak
\noindent Reduce a LENGTH applied to a literal LIST in the goal.\par
\noindent\textit{Uses:} \texttt{length-reduce}.\par
\noindent\textit{When useful:} the goal has LENGTH of a literal LIST\par

Simplify `the length of an explicit list [a1, ..., an]' to the count n.  Structural (the dual of nth-r): the spine of a list literal is total, so its length is n regardless of whether the entries are defined.

\par\smallskip{\raggedright\noindent\texttt{(rfl)}\quad\tactkind{rule}\index{rfl@\texttt{rfl} (tactic)}\par}\nobreak
\noindent Close a reflexive equality goal (t = t).\par
\noindent\textit{Uses:} \texttt{reflexivity}.\par
\noindent\textit{When useful:} the goal is `t = t' with t manifestly defined\par

Close a goal `t = t' -- a thing equals itself.  (Technically: reflexivity of equality.)

\par\smallskip{\raggedright\noindent\texttt{(qrfl)}\quad\tactkind{rule}\index{qrfl@\texttt{qrfl} (tactic)}\par}\nobreak
\noindent Quasi-reflexivity: close t = t under the partial-equality definedness reading.\par
\noindent\textit{Uses:} \texttt{quasi-reflexivity}.\par
\noindent\textit{When useful:} the goal is `t = t' but t might be UNDEFINED (partial =)\par

Close `t = t' under the partial-equality reading, where asserting t = t also asserts that t is DEFINED.  Use this rather than rfl when t might be undefined.  (Technically: quasi-reflexivity; see the partial-equality convention, where `t = t' is the definedness predicate.)


\subsection*{Propositional logic}

\par\smallskip{\raggedright\noindent\texttt{(prop)}\index{prop@\texttt{prop} (tactic)}\par}\nobreak
\noindent Decide the goal by PROPOSITIONAL logic from the context, and close it if it follows.  Every non-connective formula -- `x in a', an equation, a whole `forall(...)' -- is one opaque atom; AND/OR/NOT/IMPLIES/IFF are read.  On a goal that does NOT follow it prints the countermodel (which atom must be true, which false) and leaves the proof untouched.  Adds no trust: it decides semantically, then discharges through di/ai/oi/ass/use-em/have!/detach!, so the bill is unchanged and the recorded script is the ordinary step-by-step proof.\par
\noindent\textit{Uses:} \texttt{and-elim} \texttt{and-intro} \texttt{arith-forsome} \texttt{arith-ground} \texttt{arith-simplify} \texttt{assumption} \texttt{cartesian-decompose} \texttt{cut} \texttt{detach} \texttt{eq-subst} \texttt{forall-elim} \texttt{forall-intro} \texttt{forsome-elim} \texttt{iff-elim} \texttt{iff-intro} \texttt{implies-intro} \texttt{macete} \texttt{not-elim} \texttt{not-intro} \texttt{or-elim} \texttt{or-intro-left} \texttt{or-intro-right} \texttt{proof-by-contradiction} \texttt{theorem-assumption} \texttt{tuple-equality-decompose}.\par
\noindent\textit{When useful:} the goal follows from the hypotheses by AND/OR/NOT/IMPLIES/IFF alone -- no quantifier or equality reasoning needed\par

Close a goal that follows from the hypotheses by pure propositional reasoning -- the leaves where you can SEE the answer and still have to pick between (oi-l)(ass), (oi-r) plus a conjunction split, and (ai) on a negation.  It treats each atomic statement as a black box, so it will not instantiate a quantifier or reason about equality: if the goal needs an instance, land the instance first (fact / inst+) and run it again.  When it declines it names a falsifying assignment, which usually tells you exactly which hypothesis is missing.  (Technically: three-valued evaluation over the atoms decides the entailment; the proof is then replayed through the kernel rules, case-splitting with excluded middle where a disjunctive goal needs it.)


\subsection*{Arithmetic \& ring oracles}

\par\smallskip{\raggedright\noindent\texttt{(arith)}\quad\tactkind{oracle}\index{arith@\texttt{arith} (tactic)}\par}\nobreak
\noindent Discharge a ground arithmetic goal by evaluation.\par
\noindent\textit{Uses:} \texttt{arith-forsome} \texttt{arith-ground} \texttt{arith-simplify} \texttt{cartesian-decompose} \texttt{eq-subst} \texttt{macete} \texttt{tuple-equality-decompose}.\par
\noindent\textit{When useful:} the goal is ground arithmetic -- no variables\par

Close a goal that is a concrete numerical fact with no variables -- e.g. 2 + 3 = 5, or 7 in NN -- by just computing it.  (Technically: decision by ground arithmetic evaluation.)

\par\smallskip{\raggedright\noindent\texttt{(rs)}\quad\tactkind{oracle}\index{rs@\texttt{rs} (tactic)}\par}\nobreak
\noindent Ring-simplify the goal (normal form over the ambient ring).\par
\noindent\textit{Uses:} \texttt{ring-simplify}.\par
\noindent\textit{When useful:} the goal is a (possibly non-commutative) ring-word identity\par

Simplify the goal to a normal form in the ambient ring (which need not be commutative).  The non-commutative companion of crs.  (Technically: ring-simplify to a canonical word form.)

\par\smallskip{\raggedright\noindent\texttt{(crs)}\quad\tactkind{oracle}\index{crs@\texttt{crs} (tactic)}\par}\nobreak
\noindent Commutative-ring decision procedure: prove a polynomial identity over ZZ[generators].  Expands literal powers, so (x+y)\textasciicircum{}2 = ... closes directly.\par
\noindent\textit{Uses:} \texttt{comm-ring-simplify}.\par
\noindent\textit{When useful:} the goal is a commutative-ring identity, e.g. (x+y)\textasciicircum{}2 = x\textasciicircum{}2+2xy+y\textasciicircum{}2\par

Prove a polynomial identity that holds in EVERY commutative ring -- e.g. (x+y)\textasciicircum{}2 = x\textasciicircum{}2 + 2xy + y\textasciicircum{}2 -- by reducing both sides to a canonical sum-of-monomials form and checking they agree.  A genuine decision procedure: a true commutative-ring identity closes, a non-identity is refused.  Literal powers are expanded for you.  (Technically: normal form over the free commutative ring ZZ[generators].)

\par\smallskip{\raggedright\noindent\texttt{(simp [target])}\quad\tactkind{oracle}\index{simp@\texttt{simp} (tactic)}\par}\nobreak
\noindent Rewrite a commutative-ring SUBTERM of the goal to canonical form, IN PLACE (e.g. (x+y)\textasciicircum{}2 inside a larger goal becomes x\textasciicircum{}2 + 2*x*y + y\textasciicircum{}2).  Works on BOTH surfaces: concrete number domains (+ * - \textasciicircum{} over NN/ZZ/QQ/RR/CC) and a generic ring s ((ADD s)/(MUL s)/(NEG s), carrier (CARR s)).  No arg = outermost ring subterm; ``term'' targets a specific one.  Sound by cut + crs + eq-subst (no new kernel rule); needs the subterm's generators typed in context (true post-di), else refuses and names them.\par
\noindent\textit{Uses:} \texttt{comm-ring-simplify} \texttt{cut} \texttt{eq-subst}.\par
\noindent\textit{When useful:} a ring SUBTERM of a larger goal should be put in normal form in place\par
\par\smallskip{\raggedright\noindent\texttt{(supply)}\index{supply@\texttt{supply} (tactic)}\par}\nobreak
\noindent Close an inequality goal the way a hand proof would: beta-reduce any applied lambda, land the typing certificates and the standard bounds the linear oracle cannot see -- including the TRIANGLE inequality at the summands of an abs of a sum -- then call ineq over every premise.  Rehearses the whole sequence on a throwaway copy and does NOTHING unless it closes; on a miss it prints the sequence it tried.  Adds no trust: it drives lam-b, fact and ineq, each of which records itself.\par
\noindent\textit{Uses:} every one of the 64 kernel operations.\par

The committing form of what-now's SUPPLY THE ORACLE lane.  `ineq' reads abs(...), max(...) and (dist(s))(x,y) as opaque atoms and refuses any premise whose atoms are not certified real, so an inequality that is obviously true can fail for want of a typing nobody mentioned.  This lands them.  The one real idea is subadditivity: where the goal bounds abs(u + v), the useful intermediate term is abs(u) + abs(v), and the decomposition is in the term itself -- nothing is searched for.  Because a lambda reduction cannot be undone and can owe an unprovable leaf, the whole sequence is rehearsed before any of it is run.  (Technically: certificate closure to depth 2 over the forward-citation lane, plus the curated bound table, then Fourier-Motzkin.)

\par\smallskip{\raggedright\noindent\texttt{(ineq i1 i2 ...)}\quad\tactkind{oracle}\index{ineq@\texttt{ineq} (tactic)}\par}\nobreak
\noindent Close a linear-inequality goal over RR (<= < > >= = between RR terms) as a consequence of the named assumptions (1-based indices), via the Fourier-Motzkin/Farkas oracle.  Linearizes over + - * and the binplus/binneg/bintimes aliases; every MAXIMAL non-arithmetic subterm is an atom that must be certified in RR.  (Does NOT see through a generic ring's (ADD s)/(MUL s) -- those become opaque atoms.)\par
\noindent\textit{Uses:} \texttt{ineq}.\par
\noindent\textit{When useful:} a LINEAR <= / < over RR that follows from inequalities you can cite\par

Close a LINEAR inequality over the reals that follows from inequalities you cite (by their hypothesis numbers) -- by chaining them, adding them, and scaling by positive constants.  Anything that is not built from + - and multiplication-by-constants is treated as an opaque quantity, so it handles e.g. the triangle inequality where the distances are unknowns.  For NONLINEAR (polynomial) inequalities use sos instead.  (Technically: Fourier-Motzkin / Farkas over the ordered field RR; every atom must be certified real.)

\par\smallskip{\raggedright\noindent\texttt{(sos ``c1'' ``c2'' ...)}\quad\tactkind{oracle}\index{sos@\texttt{sos} (tactic)}\par}\nobreak
\noindent Sum-of-squares closer for a nonstrict polynomial inequality a <= b over RR (the nonlinear companion of (ineq)).  You supply the terms to be SQUARED; it finds the nonnegative coefficients.  (tactics 'sos) for the worked example.\par
\noindent\textit{Uses:} \texttt{sos}.\par
\noindent\textit{When useful:} a NONSTRICT POLYNOMIAL <= over RR (the nonlinear cousin of ineq)\par

A sum-of-squares closer.  To prove  a <= b  it is enough to exhibit b - a  as a sum of squares (each obviously >= 0), reducing the inequality to an algebraic identity.  You hand sos the terms to be SQUARED -- the c\_i, NOT the squares -- and it solves for nonnegative rationals lambda\_i with

\begin{VNBsnip}
      b - a  =  lambda_1 c_1^2 + ... + lambda_n c_n^2
\end{VNBsnip}

(matched coefficient-by-coefficient in crs's commutative-ring normal form), then PRINTS the lambda\_i it found.  You do not supply them.

Worked example.  Goal  forall([x in rr, y in rr], x*y <= x\textasciicircum{}2 + y\textasciicircum{}2).  Here b - a = x\textasciicircum{}2 - x*y + y\textasciicircum{}2 = 1/2(x-y)\textasciicircum{}2 + 1/2 x\textasciicircum{}2 + 1/2 y\textasciicircum{}2,  so the things squared are  x-y, x, y:

\begin{VNBsnip}
      (sos "x - y" "x" "y")
          ;; closes, printing  1/2(x-y)^2 + 1/2 x^2 + 1/2 y^2
\end{VNBsnip}

2*x*y <= x\textasciicircum{}2 + y\textasciicircum{}2 needs only one square:  (sos ``x - y'').  The three-variable a*b+b*c+c*a <= a\textasciicircum{}2+b\textasciicircum{}2+c\textasciicircum{}2  wants  (sos ``a - b'' ``b - c'' ``c - a'').

Notes. - Write the c\_i with the goal's own variable names.  di preserves them, so sos works before OR after di; skip di and sos peels the typed forall([... in rr]) wrapper itself. - Every generator (variable) must be certified in RR -- a goal binder x in rr  or an  (IN g RR)  assumption. - A wrong or insufficient certificate refuses cleanly: nothing is closed, you simply try other squares. - Strict (<) goals are refused -- a square can be 0, so squares alone never force a strict inequality.


\subsection*{Chained reasoning}

\par\smallskip{\raggedright\noindent\texttt{(calc L0 (rel1 L1 [just1]) (rel2 L2 [just2]) ...)}\quad\tactkind{composite}\index{calc@\texttt{calc} (tactic)}\par}\nobreak
\noindent Ground the focus goal (REL L0 Ln) by a CHAIN of intermediaries L0 rel1 L1 rel2 L2 ... reln Ln: proves each link and composes them into the endpoint.  A link no lane can close is LEFT OPEN as a leaf -- the refinement point, and the PSS candidate.  Relations: = == (folded by transitivity), < <= (folded through the co-*-trans lemmas, with = steps riding along), IFF (chained per direction).  RETURNS the links, the open ones, and their PSS candidates as an alist.\par
\noindent\textit{Uses:} every one of the 64 kernel operations.\par
\noindent\textit{When useful:} the goal is (REL L0 Ln) and you can WRITE the chain of intermediaries that gets there -- or you want the one link that will not close isolated as a PSS candidate\par

Write the argument the way you would on paper -- as a chain through intermediate terms -- and let the machine prove each link:

\begin{VNBsnip}
      Goal  k < 2*k   (with k in NN, ~(k=0))
      (calc 'k '(= (+ k 0)) '(< (+ k k)) '(= (* 2 k)))
          ;; i.e.  k = k+0 < k+k = 2*k
\end{VNBsnip}

Each step names a RELATION and the next LINE; calc cuts the link (rel L(i-1) Li), dispatches a lane at it, and on success composes the whole chain into the goal.  It is pure bookkeeping over the trusted tactics -- no kernel rule, no axiom of its own -- so a calc proof bills exactly the lemmas its lanes and composers cite.

The COMPOSER is forced by the relations you used, and there are three families: `cong' for = and ==, folded by eq-trans; `order' for < and <=, folded through the co-lt-trans / co-le-lt-trans / co-eq-lt-trans ... lemmas (a = step rewrites an endpoint of the running relation and rides along, which is why the chain above needs no separate arithmetic); and `iff', which di's the goal into its two directions and chains each with ai/detach!, because VNB cannot quantify over propositions and so has NO first-order iff-trans lemma to fold with.

The JUSTIFICATION of a link is optional and defaults to 'auto -- try the relation's lanes: crs then arith for = / ==; for < / <= the NN->RR bridge then ineq over the order-shaped premises; grind for IFF.  Otherwise pass 'scout (a small scout search), a macete/theorem NAME (mac it, then close), an explicit tactic form like '(fact 'nn-pos-of-nonzero 'k) eval'd at the link's focus, or 'open / \#f to leave the link open ON PURPOSE.

Sanity first: calc ERRORS before touching the proof if the goal's head is not the composed relation, or its LHS is not L0, or its RHS is not Ln.  A link already in the main context is taken as a GIVEN and skipped -- cutting it would be an alpha self-loop that opens no leaf.

What it is FOR is the open links.  A chain whose links are all closed is a proof; a chain with one open link has isolated your obstacle to a single formula, printed with its free variables and their context typing -- which is the PSS candidate to state as a lemma.  Refine by inserting more intermediaries until each link is discoverable.  (See also (prep 'ineq), which answers the other question: why a lane will not fire.)

\par\smallskip{\raggedright\noindent\texttt{(have! CLAIM [THUNK])}\quad\tactkind{composite}\index{have"!@\texttt{have"!} (tactic)}\par}\nobreak
\noindent Assert CLAIM as an intermediate step and carry on with it as a hypothesis: cut CLAIM, discharge the side goal from context (or by THUNK), and stay on the main branch.\par
\noindent\textit{Uses:} \texttt{and-intro} \texttt{arith-forsome} \texttt{arith-ground} \texttt{arith-simplify} \texttt{assumption} \texttt{cartesian-decompose} \texttt{cut} \texttt{detach} \texttt{eq-subst} \texttt{forall-elim} \texttt{forall-intro} \texttt{iff-intro} \texttt{implies-intro} \texttt{macete} \texttt{not-intro} \texttt{theorem-assumption} \texttt{tuple-equality-decompose}.\par
\noindent\textit{When useful:} you want to state an intermediate fact that follows immediately from context and continue with it -- the `we have X' step\par

The `we have X' step.  You state an intermediate fact CLAIM that follows immediately from what is already known; have! cuts it, proves the resulting side goal automatically, and returns you to the main branch with CLAIM now available as a hypothesis.  The automatic discharge (`from-context!') closes the everyday cases -- a conjunction, splitwise; an (IN (a*b) NN) by nn-mul-closed on the factors; a numeral membership by arith; otherwise the goal is already a context assumption up to alpha.  When the side goal needs more than that, pass a THUNK -- any tactic sequence -- as the second argument to discharge it your way.  It ERRORS, never silently no-ops, if CLAIM is already in context up to alpha: the cut would self-loop, opening one child and no main branch.  So a script reads as the argument does -- `we have q0*q0 = 3*(k*k); we have k =/= 0; ...'.  (Technically: composite -- cut then from-context! or THUNK; adds no kernel rule.)


\subsection*{Backchaining with a theorem}

\par\smallskip{\raggedright\noindent\texttt{(ta 'name)}\quad\tactkind{rule}\index{ta@\texttt{ta} (tactic)}\par}\nobreak
\noindent Theorem-assumption: bring the named installed theorem into context as an assumption.\par
\noindent\textit{Uses:} \texttt{theorem-assumption}.\par
\noindent\textit{When useful:} you want a library theorem available as a hypothesis\par

Bring an already-proved theorem into your current hypotheses so you can use it (instantiate it, detach from it, ...).  (Technically: adds the named installed theorem as an assumption.)

\par\smallskip{\raggedright\noindent\texttt{(wbc ['name])}\quad\tactkind{composite}\index{wbc@\texttt{wbc} (tactic)}\par}\nobreak
\noindent Witness-shape backchain: on an `exists v. ...' goal, cite a PSS lemma that MANUFACTURES a witness of that shape, so you can finish with inst+/grind/ew.\par
\noindent\textit{Uses:} \texttt{and-elim} \texttt{arith-forsome} \texttt{arith-ground} \texttt{arith-simplify} \texttt{cut} \texttt{detach} \texttt{forall-elim} \texttt{forsome-elim} \texttt{iff-elim} \texttt{not-elim} \texttt{or-elim} \texttt{theorem-assumption}.\par
\noindent\textit{When useful:} an EXISTENCE goal whose witness a PSS lemma manufactures (diagonalization / block-family)\par

The discovery move for existence proofs.  Faced with `there exists phi with ...', it looks up the produces-witness index for lemmas whose conclusion builds the RIGHT KIND of object -- the same existential-witness shape -- even when the rest of the statement differs (which is why plain bc* misses them: bc* demands the WHOLE conclusion unify).  E.g. on `exists phi. STRICTLY-MONO-NN(phi) and IS-CAUCHY-SEQ(SUBSEQ f phi)' it reaches for `diagonalization' (which makes a STRICTLY-MONO-NN phi from a nested family); on a nested-family goal it reaches for `block-family'.  With no argument it cites the top retrieved producer; (wbc 'name) cites the one you name.  It brings the lemma in (via ta); you then discharge its hypotheses (inst+), skolemize its existential (grind), and supply its witness to your goal (ew).  Circular `headline' lemmas that already conclude your exact goal are filtered out.  (See (witness-producers goal) for the raw retrieval; what-now lists the (wbc 'name) candidates on any existential goal.)

\par\smallskip{\raggedright\noindent\texttt{(bc impl)}\quad\tactkind{rule}\index{bc@\texttt{bc} (tactic)}\par}\nobreak
\noindent Backchain the goal through an (IMPLIES A B) already in context: if the goal matches B, the new goal is A.\par
\noindent\textit{Uses:} \texttt{backchain}.\par
\noindent\textit{When useful:} you have (P => Q) in context and the goal is Q\par

Work backwards through an implication you already have.  If `P implies Q' is among your hypotheses and your goal is Q, this reduces the goal to proving P.  The everyday `to get Q it's enough to show P'.  (Technically: backchaining the goal through an in-context implication whose conclusion matches.)

\par\smallskip{\raggedright\noindent\texttt{(bc* 'thm [((v val)...)] h1 h2 ...)}\quad\tactkind{composite}\index{bc*@\texttt{bc*} (tactic)}\par}\nobreak
\noindent Matching backchain: unify the theorem's conclusion against the goal, then leave its (instantiated) antecedents as subgoals -- optional handlers hk run on the k-th subgoal.\par
\noindent\textit{Uses:} \texttt{assumption} \texttt{backchain} \texttt{cut} \texttt{forall-elim} \texttt{theorem-assumption}.\par
\noindent\textit{When useful:} a library theorem's CONCLUSION matches your goal\par

Apply a known theorem to your goal.  If you have a theorem `if A and B then C' and your goal is its conclusion C -- matching the theorem's variables to yours, and the names of any dummy/bound variables don't matter (`there exists a net N' applies to a goal `there exists a net F') -- this replaces `prove the goal' with `prove A' and `prove B', the theorem's hypotheses.  The everyday `by Theorem X it suffices to show A and B'.  You may attach a handler to each hypothesis to dispatch it; values the match can't determine you supply as ((v val) ...).  (Technically: peels the theorem's leading universals and implications, matches the conclusion -- alpha-aware on bound variables -- and replays ta/inst/cut/bc automatically.)


\subsection*{Lambda, comprehension \& description}

\par\smallskip{\raggedright\noindent\texttt{(lam-t)}\quad\tactkind{rule}\index{lam-t@\texttt{lam-t} (tactic)}\par}\nobreak
\noindent VNB-LAMBDA typing: reduce (IN (VNB-LAMBDA ...) (FUN A B)) to TWO obligations -- the body, and that the domain is a set.\par
\noindent\textit{Uses:} \texttt{lambda-type}.\par
\noindent\textit{When useful:} the goal types a VNB-LAMBDA into FUN(A,B)\par

Show that a function defined by a formula (`x |-> ...') maps A into B -- reduces to showing that, for an arbitrary input in A, the value lies in B, AND that A is a set.  The second is not a formality: carrying the domain in the term stops one lambda from typing into FUN(A,B) for every A, but says nothing about A being a set, and a lambda over a proper class is not a function.  The lambda's declared domain must also be the FUN's, or the rule does not apply at all.  (Technically: VNB-LAMBDA typing into FUN A B.)

\par\smallskip{\raggedright\noindent\texttt{(lam-b)}\quad\tactkind{rule}\index{lam-b@\texttt{lam-b} (tactic)}\par}\nobreak
\noindent VNB-LAMBDA beta: reduce an applied lambda to its substituted body.\par
\noindent\textit{Uses:} \texttt{lambda-beta}.\par
\noindent\textit{When useful:} the goal has a VNB-LAMBDA applied to an argument\par

Simplify a function `x |-> e(x)' applied to an argument a to e(a) -- the body with the argument substituted in.  A lambda carries its DOMAIN, and it is only defined on that domain, so the reduction is licensed only where a is known to be in it: if the membership is neither in the context nor supplied by an enclosing binder, the step still fires but leaves (IN a A) as an extra subgoal.  Land the typing fact BEFORE the lam-b and no subgoal appears.  (Technically: VNB-LAMBDA beta-reduction.)

\par\smallskip{\raggedright\noindent\texttt{(lam-b-h hyp)}\quad\tactkind{rule}\index{lam-b-h@\texttt{lam-b-h} (tactic)}\par}\nobreak
\noindent VNB-LAMBDA beta in a cited ASSUMPTION -- what mac-h is to mac.\par
\noindent\textit{Uses:} \texttt{lambda-beta-hyp}.\par

Beta-reduce an applied lambda inside a hypothesis, in place.  Needed because a `fact' that instantiates a theorem's function variable at a lambda lands the APPLIED lambda in the CONTEXT, where the goal-side `lam-b' cannot reach it: union-of-opens-open at the identity family g := x |-> x lands is-open(md, big-union(i, fam, (x |-> x)(i))).  Without this the proof must detour through a cut beta-equation and a subst.  Cites nothing, so it adds no debt.

\par\smallskip{\raggedright\noindent\texttt{(sep-set)}\quad\tactkind{rule}\index{sep-set@\texttt{sep-set} (tactic)}\par}\nobreak
\noindent Separation sethood: the separation set \{x in A | p\} is a set.\par
\noindent\textit{Uses:} \texttt{sep-sethood}.\par
\noindent\textit{When useful:} the goal asserts a separation set \{ x in A | p \} is a set\par

Show that a set-builder set \{x in A | p(x)\} is genuinely a set.  (Technically: separation sethood -- a subclass of a set is a set.)

\par\smallskip{\raggedright\noindent\texttt{(sep-mi)}\quad\tactkind{rule}\index{sep-mi@\texttt{sep-mi} (tactic)}\par}\nobreak
\noindent Separation membership intro: prove (IN t \{x in A | p\}).\par
\noindent\textit{Uses:} \texttt{sep-mem-intro}.\par
\noindent\textit{When useful:} the goal is membership in a separation set \{ x in A | p \}\par

Prove that a term t belongs to \{x in A | p(x)\} by showing t lies in A and satisfies the condition p.  (Technically: separation-membership introduction.)

\par\smallskip{\raggedright\noindent\texttt{(sep-me hyp)}\quad\tactkind{rule}\index{sep-me@\texttt{sep-me} (tactic)}\par}\nobreak
\noindent Separation membership elim: split a separation-membership assumption into A-membership and the predicate.\par
\noindent\textit{Uses:} \texttt{sep-mem-elim}.\par
\noindent\textit{When useful:} a hypothesis is separation-set membership -- unpack it\par

From a hypothesis that t belongs to \{x in A | p(x)\}, extract the two facts it packs: t lies in A, and p(t) holds.  (Technically: separation-membership elimination.)

\par\smallskip{\raggedright\noindent\texttt{(comp-mi)}\quad\tactkind{rule}\index{comp-mi@\texttt{comp-mi} (tactic)}\par}\nobreak
\noindent Comprehension membership intro.\par
\noindent\textit{Uses:} \texttt{comp-mem-intro}.\par
\noindent\textit{When useful:} the goal is comprehension-set membership\par

Prove that something belongs to a comprehension set by establishing the set's defining condition for it.  (Technically: comprehension-membership introduction.)

\par\smallskip{\raggedright\noindent\texttt{(comp-me hyp)}\quad\tactkind{rule}\index{comp-me@\texttt{comp-me} (tactic)}\par}\nobreak
\noindent Comprehension membership elim.\par
\noindent\textit{Uses:} \texttt{comp-mem-elim}.\par
\noindent\textit{When useful:} a hypothesis is comprehension-set membership\par

From a comprehension-membership hypothesis, extract its defining condition.  (Technically: comprehension-membership elimination.)

\par\smallskip{\raggedright\noindent\texttt{(iota-d term)}\quad\tactkind{rule}\index{iota-d@\texttt{iota-d} (tactic)}\par}\nobreak
\noindent Definite-description: discharge the IOTA uniqueness obligation for `term'.\par
\noindent\textit{Uses:} \texttt{iota-def}.\par
\noindent\textit{When useful:} the goal involves IOTA and you must discharge its uniqueness obligation\par

Justify `the unique x such that p(x)' (a definite description) by proving that exactly one such x exists.  (Technically: IOTA -- discharge the uniqueness obligation for the described term.)

\par\smallskip{\raggedright\noindent\texttt{(bu-set)}\quad\tactkind{rule}\index{bu-set@\texttt{bu-set} (tactic)}\par}\nobreak
\noindent Big-union sethood.\par
\noindent\textit{Uses:} \texttt{big-union-sethood}.\par
\noindent\textit{When useful:} the goal asserts a BIG-UNION is a set\par

Show that a big union (the union of an indexed family of sets) is itself a set.  (Technically: big-union sethood.)

\par\smallskip{\raggedright\noindent\texttt{(bu-mi w)}\quad\tactkind{rule}\index{bu-mi@\texttt{bu-mi} (tactic)}\par}\nobreak
\noindent Big-union membership intro via the index witness w.\par
\noindent\textit{Uses:} \texttt{big-union-mem-intro}.\par
\noindent\textit{When useful:} the goal is BIG-UNION membership -- you have the index witness\par

Prove that an element lies in a big union by exhibiting one index w whose set already contains it.  (Technically: big-union-membership introduction via the index witness.)

\par\smallskip{\raggedright\noindent\texttt{(bu-me hyp)}\quad\tactkind{rule}\index{bu-me@\texttt{bu-me} (tactic)}\par}\nobreak
\noindent Big-union membership elim.\par
\noindent\textit{Uses:} \texttt{big-union-mem-elim}.\par
\noindent\textit{When useful:} a hypothesis is BIG-UNION membership -- obtain its index\par

From a hypothesis that an element lies in a big union, obtain an index whose set contains it (a fresh name for that index).  (Technically: big-union-membership elimination.)


\subsection*{Conditional terms}

\par\smallskip{\raggedright\noindent\texttt{(if-true t)}\quad\tactkind{rule}\index{if-true@\texttt{if-true} (tactic)}\par}\nobreak
\noindent Reduce an (IF p a b) term on the p branch: spawns p as a subgoal; the continuation gains (= (IF p a b) a).\par
\noindent\textit{Uses:} \texttt{if-true}.\par
\noindent\textit{When useful:} the goal has an (IF p a b) term to evaluate on the p branch\par

Evaluate a conditional term `if p then a else b' on the assumption that p holds: you take on p as a side goal, and may then use that the conditional equals a.  (Technically: if-true reduction, spawning p.)

\par\smallskip{\raggedright\noindent\texttt{(if-false t)}\quad\tactkind{rule}\index{if-false@\texttt{if-false} (tactic)}\par}\nobreak
\noindent Reduce an (IF p a b) term on the not-p branch: spawns (NOT p); the continuation gains (= (IF p a b) b).\par
\noindent\textit{Uses:} \texttt{if-false}.\par
\noindent\textit{When useful:} the goal has an (IF p a b) term to evaluate on the not-p branch\par

Evaluate a conditional term `if p then a else b' on the assumption that p fails: you take on `not p' as a side goal, and may then use that the conditional equals b.  (Technically: if-false reduction, spawning not p.)


\subsection*{Navigation, display \& tacticals}

\par\smallskip{\raggedright\noindent\texttt{(focus n)}\index{focus@\texttt{focus} (tactic)}\par}\nobreak
\noindent Switch the focus to the n-th open goal (1-based).\par
\noindent\textit{Uses:} no kernel operation: it records no inference.\par
\par\smallskip{\raggedright\noindent\texttt{(backup-one)}\quad\tactkind{meta}\index{backup-one@\texttt{backup-one} (tactic)}\par}\nobreak
\noindent Take back the last recorded command: restore the goal, the deduction graph and the proof script to the state before it.  Repeat to walk further back; (sp) clears the stack.  `undo' is an alias.\par
\noindent\textit{Uses:} no kernel operation: it records no inference.\par
\noindent\textit{When useful:} the last command was a mistake -- a greedy `di' that ate the induction, a `cut' you did not mean, a branch you want back\par

The one undo in VNB.  It restores the deduction graph -- not merely the focus -- by rolling back the journalled arrow and grounding writes and DROPPING every sequent and inference node posted since, so no node from the abandoned branch survives to be counted as an open leaf and `qed' cannot be handed a phantom obligation.  *proof-script* and the live trace are rewound with it, so the saved script and the printed proof are the proof you actually kept.  It backs up over the LAST RECORDED command: a command that changed nothing was never recorded and is not a step to back up over.  One thing is deliberately not restored -- *fresh-counter*, the eigenvariable source -- so backing up over an `ai' or `ew' that minted `u\_4' and re-running it mints `u\_5'; the proof is the same, the witness has a different name.  Depth is capped at *vnb-undo-depth* (64).

\par\smallskip{\raggedright\noindent\texttt{(undo)}\quad\tactkind{meta}\index{undo@\texttt{undo} (tactic)}\par}\nobreak
\noindent Alias for (backup-one).\par
\noindent\textit{Uses:} no kernel operation: it records no inference.\par
\noindent\textit{When useful:} the last command was a mistake -- alias for backup-one\par
\par\smallskip{\raggedright\noindent\texttt{(show)}\index{show@\texttt{show} (tactic)}\par}\nobreak
\noindent Redisplay the current proof state.\par
\noindent\textit{Uses:} no kernel operation: it records no inference.\par
\par\smallskip{\raggedright\noindent\texttt{(dial notch)}\index{dial@\texttt{dial} (tactic)}\par}\nobreak
\noindent The PRESENTATION RESOLUTION DIAL: how much of the sequent to show.  r0 kernel S-expression, r1 surface syntax (the default, and the last INVERTIBLE notch -- what is printed re-parses to what was printed), r2 guards and typing hypotheses elided, r3 also structures destructured to their declared slot names and deep subterms elided.  (dial) reports the notch, (dial n) sets it, (dial 'auto) reads the current goal and suggests one, (dial 'why) shows the score behind the suggestion.  For ONE FORMULA rather than the proof state, (dial-wff x [n]) -- x a wff, a ``surface string'' or an S-expression, n defaulting to the dial's current notch: (dial-wff (lookup-theorem 'matmul-assoc) 3).\par
\noindent\textit{Uses:} no kernel operation: it records no inference.\par

It is a FILTER, not a set of modes, and two constraints make it one.  The rungs NEST: turning it up only ever SUBTRACTS, so everything readable at a higher notch is readable below it -- checked in the suite, atom by atom.  And the APERTURE IS ON SCREEN: every notch above r1 names itself on the line and counts what it hid, because a lossy render that does not announce its own lossiness is the defect class of `(f)' printing as `f'.  The auto-notch SUGGESTS a starting notch and then STAYS PUT: it is consulted when you ask, never per sequent, so the lens cannot move while you are reading a proof.

\par\smallskip{\raggedright\noindent\texttt{(dial-wff x notch)}\index{dial-wff@\texttt{dial-wff} (tactic)}\par}\nobreak
\noindent Present ONE FORMULA at a resolution notch, instead of the current proof state: x is a wff, a ``surface string'' or an S-expression, and notch defaults to the dial's current setting.  (dial-wff (lookup-theorem 'matmul-assoc) 3).  Same rungs, same aperture line, no proof required.\par
\noindent\textit{Uses:} no kernel operation: it records no inference.\par

`sequent->string' is the display path for a proof STATE, so the dial reaches a formula only while a proof is holding it.  A theorem looked up, a statement being drafted, a hypothesis pasted from somewhere wants the same dial, and this is it.

\par\smallskip{\raggedright\noindent\texttt{(goal-status)}\index{goal-status@\texttt{goal-status} (tactic)}\par}\nobreak
\noindent One-line summary: done / N open goals.\par
\noindent\textit{Uses:} no kernel operation: it records no inference.\par
\par\smallskip{\raggedright\noindent\texttt{(repeat thunk [cap])}\index{repeat@\texttt{repeat} (tactic)}\par}\nobreak
\noindent Run thunk until it stops changing the proof state (LCF REPEAT).\par
\noindent\textit{Uses:} no kernel operation: it records no inference.\par
\par\smallskip{\raggedright\noindent\texttt{(orelse t1 t2 ...)}\index{orelse@\texttt{orelse} (tactic)}\par}\nobreak
\noindent Run thunks in order, stop at the first that makes progress (LCF ORELSE).\par
\noindent\textit{Uses:} no kernel operation: it records no inference.\par
\par\smallskip{\raggedright\noindent\texttt{(quietly thunk)}\index{quietly@\texttt{quietly} (tactic)}\par}\nobreak
\noindent Run thunk with state-dump output suppressed; returns its value.\par
\noindent\textit{Uses:} no kernel operation: it records no inference.\par

\subsection*{Choosing and naming witnesses}

\par\smallskip{\raggedright\noindent\texttt{(choose ex witness [prover])}\index{choose@\texttt{choose} (tactic)}\par}\nobreak
\noindent Choose eps such that FUBA(eps): cut the existential ex = `forsome([eps], FUBA(eps))', discharge it by exhibiting WITNESS (proving FUBA(WITNESS) by PROVER, or from the context), eliminate it, and RETURN the fresh eigenvariable with FUBA of it landed.\par
\noindent\textit{Uses:} \texttt{and-elim} \texttt{and-intro} \texttt{arith-forsome} \texttt{arith-ground} \texttt{arith-simplify} \texttt{assumption} \texttt{cartesian-decompose} \texttt{cut} \texttt{detach} \texttt{eq-subst} \texttt{forall-elim} \texttt{forall-intro} \texttt{forsome-elim} \texttt{forsome-intro} \texttt{iff-elim} \texttt{iff-intro} \texttt{implies-intro} \texttt{macete} \texttt{not-elim} \texttt{not-intro} \texttt{or-elim} \texttt{theorem-assumption} \texttt{tuple-equality-decompose}.\par

The `we may assume without loss of generality that FUBA(eps)' step: a hypothesis `forall([eps in rr], GUBA(eps))' is about to be used at some eps satisfying FUBA, so choose one first and then instantiate the universal at the name this returns (use-at / obtain-at).  You pay for the existence on the spot -- WITNESS is the value and PROVER (a thunk) shows FUBA holds of it; with no PROVER the side goal is closed from the context.  The rest of the proof then depends on FUBA(eps) alone and never on which witness paid for it, which is what makes it read as `without loss of generality'.  If the existential is already in context the cut is skipped and this is plain elimination; a formula that is not an existential is an error.  choose-pos is this tactic with FUBA frozen at pos-rr and witness 1: `let eps > 0 be given' when nothing gives it.  (Technically: composite -- cut / ew / ai and the AND-splitting of what lands; no kernel rule, no choice principle, and the recorded script is the expansion.)

\par\smallskip{\raggedright\noindent\texttt{(minimize! '(v1 ... vk) GUARD MEASURE)}\quad\tactkind{composite}\index{minimize"!@\texttt{minimize"!} (tactic)}\par}\nobreak
\noindent Choose v1..vk satisfying GUARD with the NN-valued MEASURE as small as possible: lands the witnesses and their minimality, and opens the two obligations any minimisation owes (MEASURE lands in NN; GUARD is satisfiable).\par
\noindent\textit{Uses:} \texttt{and-elim} \texttt{and-intro} \texttt{arith-forsome} \texttt{arith-ground} \texttt{arith-simplify} \texttt{assumption} \texttt{cartesian-decompose} \texttt{cut} \texttt{detach} \texttt{eq-subst} \texttt{forall-elim} \texttt{forall-intro} \texttt{forsome-elim} \texttt{forsome-intro} \texttt{iff-elim} \texttt{iff-intro} \texttt{implies-intro} \texttt{macete} \texttt{not-elim} \texttt{not-intro} \texttt{or-elim} \texttt{reflexivity} \texttt{sep-mem-elim} \texttt{sep-mem-intro} \texttt{theorem-assumption} \texttt{tuple-equality-decompose}.\par
\noindent\textit{When useful:} the goal falls to a `least such' / minimal-counterexample argument -- descent proofs (sqrt 2, sqrt 3 irrational), least-degree or least-pivot witnesses\par

The `least such' / minimal-counterexample step, mechanised.  You give three things: the variables to choose, a GUARD formula they must satisfy, and a natural-number-valued MEASURE term to make small.  On the main branch it hands you fresh eigenconstants w1..wk with two hypotheses -- GUARD holds of them, and nothing satisfying GUARD has a strictly smaller MEASURE -- and it opens exactly the two side goals a minimisation genuinely owes: that MEASURE lands in NN wherever GUARD holds (the TYPE goal), and that GUARD holds somewhere (the NONEMPTY goal).  Everything in between -- forming the value set, well-ordering it, unpacking the least witness, and restating minimality in terms of your variables rather than the set -- is done for you.  Every vi must occur in MEASURE (a variable the measure ignores is not one you are minimising over).  Because it quantifies over the vk at the META level, no lambda and no FUN(U,NN) typing obligation ever enters the logic -- it takes a formula and a term, which is what a driver has in hand, and never forms the set U at all.  Returns (list (w1 ... wk) TYPE-node NONEMPTY-node); either node is \#f when that obligation was already in context up to alpha, so test before refocusing.  (Technically: composite -- it drives cut / di / ai / ew / sep-mi / sep-me / fact / inst / detach! / subst / ass / rfl and adds no kernel rule; its one mathematical appeal is `nn-least-element', the well-ordering of NN, proven from ord-well-ordered.  NO choice principle is used: well-ordering returns a MEMBER of a separation set, and sep-me recovers the witness.)

\par\smallskip{\raggedright\noindent\texttt{(obtain LANE)}\quad\tactkind{composite}\index{obtain@\texttt{obtain} (tactic)}\par}\nobreak
\noindent Run LANE -- a thunk whose forward step lands a `there exists' into context -- then eliminate that existential and RETURN the fresh witness's name, read off the proof state rather than guessed.\par
\noindent\textit{Uses:} every one of the 64 kernel operations.\par
\noindent\textit{When useful:} a forward step yields `there exists ...' and you want to name the witness for later use, without guessing the engine's eigenvariable\par

The `let w be such an x' step, with the eigenvariable named for you.  LANE is a zero-argument procedure whose effect is to land some `there exists v. P(v)' among your hypotheses -- e.g. (lambda () (fact 'nn-3-div-square w)).  obtain runs it, spots the existential that newly appeared, eliminates it (introducing a fresh eigenvariable and landing P at that witness), unpacks any conjunction in the body, and hands back the eigenvariable's NAME -- identified by free-variable set-difference (the symbol now in context that was not there before), so the driver never guesses what the engine called it.  This is the robust way to name a descended witness: (define k (obtain (lambda () ...))), then use k.  Returns \#f (with a note) if the lane landed no existential.  (Technically: composite -- a guarded forward discharge followed by forsome-elim (ai) and AND-splitting.  Plain existential elimination, so it owes nothing and uses no choice.)

\par\smallskip{\raggedright\noindent\texttt{(vlet (n1 ...) FORMER)}\quad\tactkind{composite}\index{vlet@\texttt{vlet} (tactic)}\par}\nobreak
\noindent Bind proof-object names from the current state.  FORMER is (match PATTERN) -- bind each name to its slot in a context formula matching PATTERN -- or (choice [v body]) -- eliminate an in-context existential and bind its witness.\par
\noindent\textit{Uses:} \texttt{and-elim} \texttt{cut} \texttt{forsome-elim} \texttt{iff-elim} \texttt{not-elim} \texttt{or-elim}.\par
\noindent\textit{When useful:} you need to NAME a witness or a matched subterm from the proof state, rather than navigate to it by shape\par

A binding form for the pieces of a proof, so a driver NAMES what it needs instead of navigating to it by shape.  Two FORMERs.  (match PATTERN): the listed names ARE the holes (wildcards) in PATTERN; everything else is literal.  vlet searches the context for a formula or subterm matching PATTERN and binds each name to what filled its slot -- pure selection, no proof step, no obligation; a hard error if nothing matches (you named a piece of something absent).  (choice): eliminate the sole existential in context and bind its witness.  (choice v body): present-else-debt -- if `there exists v. body' is already in context up to alpha, eliminate it; otherwise cut it, LEAVE the existence side-goal open as a debt leaf, and eliminate on the main branch -- either way the witness is bound and body[witness] is landed.  vlet's (choice) and `obtain' are the same underlying mechanism -- eliminate an existential, name the witness by free-variable difference -- differing only in what they take: obtain runs a lane and names what it just produced, vlet names an existential already (or, with a body, about to be) in context.  No choice AXIOM is used; a single witness is plain existential-elimination.  (Technically: a define-syntax expanding to (define n ...) over vlet--match / vlet--choice!; composite, no kernel rule beyond forsome-elim and, for present-else-debt, cut.)




%% =========================================================
\section{Focus management}

After each command the focus advances automatically to the first
ungrounded subgoal just introduced.  If there are no new subgoals, it
moves to any remaining open goal.

\begin{remark}
Every proof command operates exclusively on the focus node; the other
open goals are untouched.  There can be several open goals simultaneously
(for example, after \texttt{(ci)} on a conjunction goal, or after
\texttt{(cut f)} which always creates two subgoals).  \texttt{(show)}
displays them all, numbering the non-focus goals so you can use those
numbers with \texttt{(focus $n$)}.  Until you call \texttt{(focus $n$)},
the system always works on whichever goal it chose automatically.
\end{remark}

In the interactive session use:
\begin{VNBcode}
(show)       ; list all open goals (numbered) and the current focus
(focus n)    ; switch focus to the n-th open goal (1-based)
\end{VNBcode}

The underlying API (for use outside the interactive session) is:
\begin{VNBcode}
(focus-on ps sqn)            ; focus on a specific sequent node
(focus-on-first-open ps)     ; focus on the first open goal
(proof-open-goals ps)        ; list all open goals
\end{VNBcode}

%% =========================================================
\section{Macetes}
\label{sec:macetes}

A \emph{macete} is a named, composable rewrite strategy.  When a
theorem is installed, its macete is automatically derived by extracting a
source pattern and replacement pattern from the theorem's logical form:

\begin{center}
\begin{tabular}{ll}
\toprule
Theorem form & Rewrite \\
\midrule
$\forall\bar{x}.\;(s = r)$ & $s \mapsto r$ \\
$\forall\bar{x}.\;(s \Leftrightarrow r)$ & $s \mapsto r$ \\
$\forall\bar{x}.\;(\text{conds} \Rightarrow (s = r))$ & $s \mapsto r$ when conds hold \\
$\forall\bar{x}.\;p$ & $p \mapsto \top$ \\
\bottomrule
\end{tabular}
\end{center}

The quantifiers need not all be leading.  A \texttt{forall} occurring in
positive position --- under an \texttt{implies} consequent or inside an
\texttt{and} --- is pulled to the front before the source and replacement
patterns are extracted, so a theorem such as
\texttt{forall([a,b], $H$ implies forall(x, $s$ iff $r$))} still yields the
rewrite $s \mapsto r$ (conditional on $H$).

Macetes can be composed:
\begin{VNBcode}
(macete-series m1 m2 m3)   ; apply in sequence
(macete-repeat m)           ; apply until no change
(macete-try m)              ; apply m, succeed even if m fails
(macete-choice m1 m2 m3)   ; first one that applies
\end{VNBcode}

To install a custom macete:
\begin{VNBcode}
(install-theorem! 'eq-symm
  (make-wff-from-string "forall([x, y], x = y iff y = x)"))

; Use it:
(cmd-apply-macete ps 'eq-symm)
\end{VNBcode}

\subsection{Applying macetes to wffs outside a proof}

\begin{sloppypar}The function \texttt{apply-macete} applies a named macete directly to a
\texttt{wff} without creating a proof state:\end{sloppypar}

\begin{itemize}
\item \texttt{(apply-macete \textit{name} \textit{w})} --- rewrites \textit{w};
  returns a new \texttt{wff} if any rewrite fires, or \texttt{\#f} if the
  pattern matches nowhere in \textit{w}.
\end{itemize}

This is purely syntactic: no deduction graph is involved and no
theorem is proved.  It is useful for inspecting rewrite behavior,
preprocessing formulas, or building higher-level tools outside of proofs.

\paragraph{Unconditional equation macete.}
\begin{VNBcode}
(install-theorem! 'nn-zero-succ
  (make-wff-from-string "forall([n in NN], succ(n) = n + 1)"))

(apply-macete 'nn-zero-succ (make-wff-from-string "succ(3) = x"))
; => wff for "3 + 1 = x"

(apply-macete 'nn-zero-succ (make-wff-from-string "3 in nn"))
; => #f               -- macete doesn't fire
\end{VNBcode}

\paragraph{Biconditional (\texttt{iff}) macete.}
\begin{VNBcode}
(install-theorem! 'eq-symm
  (make-wff-from-string "forall([x, y], x = y iff y = x)"))

(apply-macete 'eq-symm (make-wff-from-string "n + m = 0"))
; => wff for "0 = n + m"

(apply-macete 'eq-symm (make-wff-from-string "n = m implies truth"))
; => wff for "m = n implies truth"

(apply-macete 'eq-symm (make-wff-from-string "n in nn"))
; => #f
\end{VNBcode}

\paragraph{Closure axiom macete (conditional).}
\begin{VNBcode}
;; nn-add-closed: forall([n in NN, m in NN], n + m in NN)
;; Condition (n in NN and m in NN) must be present in local context.

(apply-macete 'nn-add-closed
  (make-wff-from-string "n in nn and m in nn implies n + m in nn"))
; => wff for "n in nn and m in nn implies truth"

(apply-macete 'nn-add-closed (make-wff-from-string "n + m in nn"))
; => #f  -- conditions not in context; rewrite correctly suppressed
\end{VNBcode}

\begin{remark}
\texttt{apply-macete} applies the rewrite depth-first throughout the
formula (outermost first within each subexpression).  Conditions of a
conditional macete are checked against the \emph{local context}
accumulated as the rewriter descends: when descending into the
consequent of \texttt{implies}, the antecedent is added to the local
context; when descending into the right conjunct of \texttt{and}, the
left conjunct is added; and so on (Monk~\cite{Monk1988}).  A condition that
matches a formula already present in the local context is considered
discharged; if any condition is unmet at a candidate site, the rewrite
is suppressed at that site.
\end{remark}

%% =========================================================
\section{Example proofs}

\subsection{Identity}

\begin{VNBcode}
(sp (make-wff-from-string "x in NN implies x in NN"))
(di)    ; assume (x in nn), prove (x in nn)
(ass)   ; it is in context
; => Proof complete.
\end{VNBcode}

\subsection{AND elimination}

\begin{VNBcode}
(sp (make-wff-from-string "n in NN and m in NN implies n in NN"))
(di)                          ; assume (n in nn) and (m in nn)
(ai "n in NN and m in NN")   ; split into (n in nn) and (m in nn)
(ass)                         ; prove (n in nn)
\end{VNBcode}

\subsection{VNB axiom instance}

Proving \texttt{"a in b implies a in SET"} from the base axiom:

\begin{VNBcode}
(sp (make-wff-from-string "a in b implies a in SET"))
(di)
(ta 'membership-implies-sethood)
(inst "forall([a, b], a in b implies a in set)" 'a)
(inst "forall([b], a in b implies a in set)" 'b)
(bc "a in b implies a in set")
(ass)
\end{VNBcode}

\subsection{Double negation (classical)}

\begin{VNBcode}
(sp (make-wff-from-string "not(not(x in NN)) implies x in NN"))
(di)                              ; assume not(not((x in nn))), prove (x in nn)
(pbc)                             ; assume not((x in nn)), prove falsity
(ai "not(not(x in NN))")         ; not(not(P)) and not(P) yield falsity
\end{VNBcode}

\subsection{Ordered tuples and Cartesian products}

\begin{VNBcode}
; From x in A, y in B, z in C prove [x,y,z] in cartesian(A,B,C)
(sp (make-wff-from-string
      "x in A and y in B and z in C implies [x, y, z] in cartesian(A, B, C)"))
(di)
(ai "x in A and y in B and z in C")
(ai "y in B and z in C")
(ci)              ; cartesian-intro generates 3 subgoals
(ass) (ass) (ass)

; NTH projection: nth(2, [x, y, z]) = y
(sp (make-wff-from-string "nth(2, [x, y, z]) = y"))
(nth-r)           ; reduces goal to (y = y)
(rfl)             ; closes by reflexivity
\end{VNBcode}

\subsection{Mathematical structures}

Algebraic structures are declared with \texttt{declare-structure}.
No quoting is needed: carrier names and operation signatures are written
directly.

\begin{VNBcode}
(declare-structure GROUP
  (carriers CARR)
  (op OPR (CARTESIAN CARR CARR) CARR)  ; operation: CARR x CARR -> CARR
  (constant IDEN CARR)                 ; identity: IDEN in CARR
  (op INV CARR CARR))                  ; inverse: CARR -> CARR
; Accessor indices: CARR -> 1, OPR -> 2, IDEN -> 3, INV -> 4
; (the group laws are folded in by additional (property ...) clauses)
\end{VNBcode}

\begin{sloppypar}The \texttt{(op name domain range)} clause declares
$\texttt{name}(s) \in \texttt{FUN}(\texttt{domain}, \texttt{range})$.
Higher-arity operations write the Cartesian product explicitly:
\texttt{(op OPR (CARTESIAN CARR CARR) CARR)} for a binary operation.
A \texttt{(constant name S)} clause declares an element of set $S$.
Accessor indices are assigned sequentially: carriers first (starting
at~1), then operations in declaration order.\end{sloppypar}

A \texttt{declare-structure} call installs two things automatically:
\begin{itemize}
\item \emph{Accessor macetes}: \texttt{CARR}$(s)$ rewrites to
  \texttt{nth(1,$s$)}, \texttt{OPR}$(s)$ to \texttt{nth(2,$s$)}, etc.
\item \emph{Definitional biconditional} \texttt{IS-NAME}: the axiom
  $\forall s.\;\texttt{IS-GROUP}(s) \Leftrightarrow
   (\texttt{CARR}(s)\in\mathrm{SET} \wedge \texttt{OPR}(s)\in\cdots)$.
\end{itemize}

Given \texttt{IS-GROUP(g)} in context, the \texttt{IS-GROUP} macete
unfolds to:
\begin{VNBsnip}
CARR(g) in SET
and OPR(g) in fun(cartesian(CARR(g), CARR(g)), CARR(g))
and IDEN(g) in CARR(g)
and INV(g) in fun(CARR(g), CARR(g))
\end{VNBsnip}

\paragraph{Pre-loaded structures.}
The following structures are defined in \texttt{structure-library/} (one
file per structure) and loaded automatically.

\begin{center}\scriptsize
\begin{tabular}{@{}p{2.3cm}p{4.3cm}p{5.3cm}@{}}
\toprule
\textbf{Structure} & \textbf{Accessor slots (index)} & \textbf{Characteristic axioms} \\
\midrule
\texttt{SEMIGROUP} &
  \texttt{CARR}\,(1), \texttt{OPR}\,(2) &
  \texttt{semigroup-assoc} \\[4pt]
\texttt{MONOID} &
  \texttt{CARR}\,(1), \texttt{OPR}\,(2), \texttt{IDEN}\,(3) &
  \texttt{monoid-assoc},
  \texttt{monoid-left-id},
  \texttt{monoid-right-id} \\[4pt]
\texttt{GROUP} &
  \texttt{CARR}\,(1), \texttt{OPR}\,(2), \texttt{IDEN}\,(3), \texttt{INV}\,(4) &
  \texttt{group-assoc},
  \texttt{group-left-id},
  \texttt{group-left-inv} \\[4pt]
\texttt{RING} &
  \texttt{CARR}\,(1), \texttt{ADD}\,(2), \texttt{MUL}\,(3),
  \texttt{NEG}\,(4), \texttt{ZERO}\,(5), \texttt{ONE}\,(6) &
  \texttt{ring-add-assoc}, \texttt{ring-add-comm},
  \texttt{ring-add-left-id}, \texttt{ring-add-left-inv},
  \texttt{ring-mul-assoc}, \texttt{ring-mul-left-id},
  \texttt{ring-mul-right-id},
  \texttt{ring-left-dist}, \texttt{ring-right-dist} \\[4pt]
\texttt{METRIC-SPACE} &
  \texttt{PTS}\,(1), \texttt{DIST}\,(2) &
  \texttt{metric-pos},
  \texttt{metric-self-zero},
  \texttt{metric-zero-eq},
  \texttt{metric-sym},
  \texttt{metric-triangle} \\
\bottomrule
\end{tabular}
\end{center}

\begin{remark}
Each entry above is citable by name with \texttt{(ta 'name)}, which transports
it into the current proof.  They do \emph{not} all have the same standing,
however, and the column heading is a convenience rather than a claim.

A structure declaration generates an \texttt{IS-X} defining biconditional, and
each operation-property conjunct of it, unfolded, \emph{is} one of the
structure's laws.  Stating such a law separately therefore adds no mathematical
content, and the question is only how the system records that fact.  Three
answers are in use, and they cost a citing \texttt{qed} the same --- nothing ---
by three different routes:

\begin{itemize}
\item \emph{Proved.}  The \textsc{Metric-Space} laws (\texttt{metric-pos},
  \texttt{metric-self-zero}, \texttt{metric-zero-eq}, \texttt{metric-sym},
  \texttt{metric-triangle}) are established by \texttt{prove-metric-law!} in
  \texttt{structure-library/metric-laws.scm}; the \textsc{Group} laws
  (\texttt{group-assoc}, \texttt{group-left-id}, \texttt{group-left-inv},
  \texttt{group-identity-in}) and \texttt{abelian-group-opr-comm} by
  \texttt{stl--project!} in \texttt{structure-library/subtype-laws.scm}.  The
  unfold is \emph{checked}, not asserted to exist.
\item \emph{Installed and stamped \texttt{definitional}.}  The \textsc{Ring}
  laws are installed by \texttt{add-axiom!} in
  \texttt{structure-library/ring.scm} and then swept with
  \texttt{register-provenance!\ \ldots\ 'definitional} in the same file.  A
  \texttt{definitional} fact contributes the empty set to every bill, exactly as
  a \texttt{primitive} one does.
\item \emph{Wrapped at installation.}  \texttt{structure-library/module.scm}
  reaches the same place by the other door, installing its projections inside
  \texttt{(fluid-let ((*current-provenance* 'definitional)) \ldots)}.
\end{itemize}

Two laws in the table are genuinely \emph{derived} rather than projected, and
they do carry a warrant: \texttt{ring-mul-zero-left} and
\texttt{ring-mul-zero-right} ($0\cdot a = 0$, which needs distributivity and
cancellation, not merely the definition of a ring).
\texttt{reference/PROOF-DEBT.md} is the authority on which is which; this table
is not.  For example, to apply
\texttt{group-left-inv} (which asserts
$\forall s,a.\;\texttt{IS-GROUP}(s)\Rightarrow a\in\texttt{CARR}(s)
 \Rightarrow \texttt{OPR}(s)(\texttt{INV}(s)(a),a)=\texttt{IDEN}(s)$),
supply the group and the element as instantiation witnesses.
\end{remark}

If you prefer to work with raw tuples directly, the n-ary constructors
are still available:
\begin{VNBcode}
; Prove a tuple is in a Cartesian product
(sp (make-wff-from-string
  "G in power(U) and op in fun(cartesian(G,G),G) and e in G and inv in fun(G,G)
   implies [G, op, e, inv]
   in cartesian(power(U), fun(cartesian(G,G),G), G, fun(G,G))"))
(di)
(ai "G in power(U) and op in fun(cartesian(G,G),G) and e in G and inv in fun(G,G)")
(ai "op in fun(cartesian(G,G),G) and e in G and inv in fun(G,G)")
(ai "e in G and inv in fun(G,G)")
(ci)                              ; cartesian-intro: 4 subgoals
(ass) (ass) (ass) (ass)
\end{VNBcode}

\subsection{Applying a macete inside a proof}

The short-form command \texttt{(mac 'name)} applies the named macete to
the \emph{current goal}.  It unifies the macete's source pattern against
the goal formula (outermost-first, depth-first) and replaces the first
matched subterm with the replacement pattern.

\paragraph{Biconditional macete.}
\texttt{pairing-membership} states
$\forall a\,b\,x.\;x \in \mathrm{pair}(a,b) \Leftrightarrow (x=a \vee x=b)$.
Its source pattern is \verb|x in pair(a, b)| and its replacement is
\verb|x = a or x = b|.

\begin{VNBcode}
(sp (make-wff-from-string "2 in pair(2, 3)"))
; proof state:  []  |-  2 in pair(2, 3)
(mac 'pairing-membership)
; source: x in pair(a, b)   replacement: x = a or x = b
; unifies: x=2, a=2, b=3
; goal rewritten:  2 in pair(2, 3)  ->  2 = 2 or 2 = 3
; proof state:  []  |-  2 = 2 or 2 = 3
(oi-l)      ; or-intro-left: new goal  2 = 2
(rfl)       ; reflexivity
; => Proof complete.
\end{VNBcode}

\paragraph{Rewrite fires inside a larger formula.}
The rewrite finds its source pattern at the first matching position,
which may be nested inside a connective.

\begin{VNBcode}
(sp (make-wff-from-string "x in pair(a, b) implies x = a or x = b"))
; proof state:  []  |-  (x in pair(a,b)) implies (x = a or x = b)
(mac 'pairing-membership)
; pattern "x in pair(a, b)" matched at the left of implies
; goal rewritten:  (x = a or x = b) implies (x = a or x = b)
(di)        ; discharge implies: push antecedent to context
; context: ["x = a or x = b"]   goal: "x = a or x = b"
(ass)
; => Proof complete.
\end{VNBcode}

\begin{remark}
Macetes rewrite the goal; they do not rewrite the context.  If the
matching pattern appears only in an assumption, \texttt{(mac 'name)}
will not fire.  To use an assumption as a rewrite source, retrieve it
with \texttt{(ai ...)} so that it appears in the goal.
\end{remark}

\paragraph{Conditional macete.}
For a conditional theorem
$\forall\bar{x}.\;(\mathit{cond} \Rightarrow s = r)$, the macete still
fires unconditionally: it rewrites $s \mapsto r$ without generating
side-goals for $\mathit{cond}$.  The goal $s$ is replaced by~$r$
(or by $\top$ when the conclusion is not an equation or biconditional),
and the resulting $\top$ goal is closed with \texttt{(ass)}.

\begin{VNBcode}
; nn-add-closed:
;   forall([a in NN, b in NN], a + b in NN)
; As a macete: source = "a + b in nn"
;              replacement = truth
(sp (make-wff-from-string "2 + 3 in NN"))
; proof state:  []  |-  2 + 3 in NN
(mac 'nn-add-closed)
; unifies a=2, b=3
; conditions (a in NN, b in NN) not discharged
; goal rewritten:  2 + 3 in NN  ->  truth
(ass)       ; truth is trivially grounded
; => Proof complete.
\end{VNBcode}

\begin{remark}
The conditions of a conditional macete are silently dropped.  This
makes \texttt{(mac 'name)} fast for goals whose conditions are obviously
satisfied, but it is the user's responsibility to verify them; the
proof sketch above is technically unsound without separate proofs that
$2 \in \mathit{NN}$ and $3 \in \mathit{NN}$.  For fully verified
condition discharge, use \texttt{(ta 'name)} followed by explicit
instantiation and back-chaining.
\end{remark}

\subsection{Ground arithmetic with \texttt{(arith)}}

The \texttt{(arith)} command (primitive \texttt{cmd-arith}) closes any
goal whose formula reduces to true under exact arithmetic evaluation:

\begin{VNBcode}
(sp (make-wff-from-string "1/3 + 1/6 = 1/2"))
(arith)                              ; closes by exact rational arithmetic

(sp (make-wff-from-string "2 ^ 10 = 1024"))
(arith)                              ; 2^10 evaluates to 1024

(sp (make-wff-from-string "3 + 4 <= 2 * 5"))
(arith)                              ; 7 <= 10
\end{VNBcode}

\subsection{The calculator pattern}

\texttt{(arith)} also recognises a goal of the form
\begin{VNBsnip}
forsome([n in CLASS], EXPR = n)
\end{VNBsnip}
where \texttt{EXPR} is a ground arithmetic expression and \texttt{CLASS}
is one of \texttt{NN}, \texttt{ZZ}, \texttt{QQ}, \texttt{RR},
\texttt{CC}.  It evaluates \texttt{EXPR}, checks that the result belongs
to \texttt{CLASS}, and closes the goal with that value as the witness.
This lets the prover double as a calculator:

\begin{VNBcode}
(sp (make-wff-from-string "forsome([n in NN], 2^10 = n)"))
(arith)   ; closes with witness 1024

(sp (make-wff-from-string "forsome([q in QQ], 1/3 + 1/4 = q)"))
(arith)   ; closes with witness 7/12
\end{VNBcode}

\begin{remark}
The symmetric form \texttt{n = EXPR} is also accepted.
The pattern is intentionally narrow: \texttt{EXPR} must be fully ground
(no free variables), and the existential variable must appear alone on
one side of the equality.
\end{remark}

\subsection{\texorpdfstring{$\beta$}{beta}-reduction with \texttt{(beta)}}

Applied \texttt{lambda} (or \texttt{lambdoid}) terms are reduced
in-place, then the underlying numeric or logical content is closed
by another rule:

\begin{VNBcode}
(sp (make-wff-from-string "lambda([x in nn], 2 * x)(5) = 10"))
(beta)   ; goal becomes (2 * 5 = 10)
(arith)  ; closes
\end{VNBcode}

\subsection{Auto-installing arithmetic theorems}

\texttt{vnb-create-num-theorem} verifies a closed arithmetic wff by
ground evaluation and installs it as a named theorem.  The theorem is
then available via \texttt{(ta \textit{name})}:

\begin{VNBcode}
(vnb-create-num-theorem 'pythag-3-4-5 "3 ^ 2 + 4 ^ 2 = 5 ^ 2")
; => pythag-3-4-5

(sp (make-wff-from-string "truth implies 3 ^ 2 + 4 ^ 2 = 5 ^ 2"))
(di) (ta 'pythag-3-4-5) (ass)
; => Proof complete.
\end{VNBcode}

False or undecidable formulas are rejected:

\begin{VNBcode}
(vnb-create-num-theorem 'bogus "1 + 1 = 3")
; ERROR: formula evaluates to FALSE

(vnb-create-num-theorem 'bogus2 "x = 5")
; ERROR: not a decidable closed arithmetic sentence
\end{VNBcode}

\subsection{Installing a proven theorem}
\label{sec:qed-and-subst}

A completed proof does not by itself add anything to the library.  Use
\texttt{(qed \textit{name})} to install the root assertion as a named
theorem:

\begin{VNBcode}
(sp (make-wff-from-string "n in NN and m in NN implies m in NN and n in NN"))
(di) (ai "n in NN and m in NN") (di) (ass) (ass)
(qed 'and-comm)
; => and-comm
; (lookup-theorem 'and-comm) returns the formula S-expression.
\end{VNBcode}

\paragraph{What \texttt{subst-wff} is and is not.}
\texttt{subst-wff} applies capture-avoiding substitution to a wff's
underlying S-expression and returns a new \texttt{wff} (re-validated
through \texttt{make-wff}).  It is \emph{not} a logical inference;
after calling it you must reprove the result.  It is useful for
generating related goals from a known formula:

\begin{VNBcode}
(define w
  (subst-wff '((n . (in x nn)) (m . (in y rr)))
             (lookup-theorem 'and-comm)))
; w is the wff for (x in nn) and (y in rr) implies (y in rr) and (x in nn)
\end{VNBcode}

\subsection{Saving and replaying proof scripts}
\label{sec:replay}

Every interactive command you issue between \texttt{(sp\ \ldots)} and
\texttt{(qed\ \ldots)} is appended to a global accumulator
\texttt{*proof-script*}.  \texttt{(qed \textit{name})} installs the
proven theorem and saves the script under the same name; you can also
call \texttt{(save-proof \textit{name})} on demand.

\begin{VNBcode}
(sp (make-wff-from-string "n in NN and m in NN implies m in NN and n in NN"))
(di) (ai "n in NN and m in NN") (di) (ass) (ass)
(qed 'and-comm)
; theorem AND script installed under the same name

(lookup-proof 'and-comm)
; => ((di) (ai (and (in n nn) (in m nn))) (di) (ass) (ass))
\end{VNBcode}

\paragraph{Replay verbatim.}
\texttt{(replay-proof \textit{name})} re-runs the saved commands
sequentially against the current proof state.  Useful for redoing a
proof in a fresh session, or applying an identical script to a sequent
of the same shape.

%% =========================================================
\subsection{NN induction for type properties}
\label{sec:ex-ni-type}

The \texttt{(ni)} command reduces a goal
$\forall n \in \mathit{NN}.\; P(n)$ to two subgoals:
a base case $P(0)$ (focused first) and a step case
$\forall n \in \mathit{NN}.\; P(n) \Rightarrow P(\mathrm{succ}(n))$.
The following proof shows that $\mathit{PROD\text{-}ORD}(m, f, n)$
stays in the carrier $\mathit{CARR}(m)$ for every $n \in \mathit{NN}$,
deriving \texttt{prod-ord-type} from first principles.

\begin{VNBcode}
(sp (make-wff '(IMPLIES (IS-MONOID m)
               (IMPLIES (IN f (FUN NN (CARR m)))
                 (FORALL n (IMPLIES (IN n NN)
                   (IN (PROD-ORD m f n) (CARR m))))))))
(di) (di)   ; IS-MONOID(m) and f in FUN(NN,CARR(m)) enter context
(ni)        ; NN-induction; BASE gets focus first

;; --- BASE: PROD-ORD(m,f,0) in CARR(m) ---
;; Strategy: PROD-ORD(m,f,0) = IDEN(m) (prod-ord-zero),
;;           IDEN(m) in CARR(m)            (monoid-identity-in),
;;           so by eq-subst-membership the goal closes.
(ta 'eq-subst-membership)
(inst ... '(PROD-ORD m f 0)) (inst ... '(IDEN m)) (inst ... '(CARR m))
(bc '(IMPLIES (AND (= (PROD-ORD m f 0) (IDEN m)) (IN (IDEN m) (CARR m)))
              (IN (PROD-ORD m f 0) (CARR m))))
(di)    ; AND-split: focus = equality goal
(ta 'prod-ord-zero) (inst ... 'm) (inst ... 'f) (ass)
        ; focus jumps: manually refocus on membership goal
(ta 'monoid-identity-in) (inst ... 'm)
(bc '(IMPLIES (IS-MONOID m) (IN (IDEN m) (CARR m)))) (ass)
        ; BASE complete; focus moves to STEP

;; --- STEP: forall([n in NN], IH => PROD-ORD(m,f,succ n) in CARR(m)) ---
(di)    ; peel FORALL n -> n_k; (IN n_k NN) enters context
;; n_k is the bound name introduced by (di); fetch it via:
;;   (cadr (wff-formula
;;           (car (sequent-node-assumptions
;;                  (proof-state-focus *ps*)))))
(di)    ; peel IH: PROD-ORD(m,f,n_k) in CARR(m) enters context
;; Strategy:
;;   PROD-ORD(m,f,succ n_k)
;;     = OPR(m)(PROD-ORD(m,f,n_k), f(n_k))   (prod-ord-succ)
;;   then close via monoid-carrier-closed-opr.
(ta 'eq-subst-membership) (inst ...) (inst ...) (inst ...)
(bc ...)
(di)    ; AND-split: focus = equality goal
(ta 'prod-ord-succ)
(inst ...) (inst ...) (inst ...)
(bc ...) (ass)   ; n_k in NN closes it
;; focus moves to membership goal: OPR(m)(...) in CARR(m)
(ta 'monoid-carrier-closed-opr)
(inst ...) (inst ...) (inst ...)
(bc ...)
(di) (ass)   ; IS-MONOID(m)
(di) (ass)   ; PROD-ORD(m,f,n_k) in CARR(m) [IH]
;; f(n_k) in CARR(m) via fun-apply-type:
(ta 'fun-apply-type)
(inst ...) (inst ...) (inst ...) (inst ...)
(bc ...) (di) (ass) (ass)
;; closes f in FUN(NN,CARR(m)) and n_k in NN
\end{VNBcode}

\begin{remark}
The \texttt{inst} chains are elided above; each peels one leading
\texttt{forall} from a theorem schema.  Two subtleties arise in practice.
First, \texttt{(di)} on an AND-goal always focuses on the left (equality)
conjunct first.  After \texttt{(ass)} closes that conjunct, automatic
focus advancement jumps to the earliest ungrounded goal in the deduction
graph --- which may be the STEP node, not the right (membership) conjunct.
The fix is to capture the membership node
immediately after the AND-split:
\begin{VNBcode}
(di)   ; AND-split
(let* ((dg       (proof-state-dg *ps*))
       (mem-node (car (reverse (dg-sequent-nodes dg)))))
  ... prove equality ... (ass)   ; focus jumps
  (set-proof-state-focus! *ps* mem-node)
  ... prove membership ...)
\end{VNBcode}

Second, \texttt{alpha-equiv?} --- used by \texttt{(inst)} and
\texttt{(bc)} when searching assumptions --- relies on structural
equality of formula heads.  Compound heads such as
\texttt{((OPR m) a b)} (a curried two-argument application; see
Section~\ref{sec:apply-tupling}) are compared correctly only when
the head pair is the same \emph{object}.  If \texttt{alpha-equiv?}
fails to locate a formula you know is in context, the likely cause
is a compound-head mismatch; the fix is in
\texttt{expressions.scm}'s \texttt{alpha-equiv-under?}.
\end{remark}

\begin{remark}[The curried/tupled apply convention]
\label{sec:apply-tupling}
VNB's S-expression form \texttt{(f a$_1$ \ldots a$_n$)} is what the
string-syntax parser emits from \texttt{f(a$_1$,\,\ldots,\,a$_n$)} ---
verbatim, no implicit tupling.  An operation typed
\texttt{f $\in$ fun(cartesian(A$_1$,\,\ldots,\,A$_n$),~B)}, however,
takes one argument from the Cartesian product.  The two views are
reconciled by the library axioms
\begin{equation*}
\texttt{apply-tupling-$n$:} \qquad
(f\ a_1\ \ldots\ a_n)\ \mathrel{==}\ (f\ (\texttt{LIST}\ a_1\ \ldots\ a_n))
\quad\text{for each arity }n
\end{equation*}
(installed for $n=1,2,3$ in \texttt{theorem-library/axioms.scm};
add more arities by declaration as needed).  Quasi-equality
$\mathrel{==}$, not partial $=$: the convention holds vacuously when
either side is undefined.  Combined with \texttt{fun-apply-type} and
\texttt{quasi-eq-def}, these give the $n$-ary apply closure
\((f\ a_1\ \ldots\ a_n)\in B\) from the operation's signature and
the membership of each $a_i$ in its declared component.
\end{remark}

\subsection{Infinite descent: \texorpdfstring{$\sqrt3$}{sqrt3} is irrational}
\label{sec:ex-sqrt3}

Stated without reals or division --- the form the checker proves, entirely
inside $\mathit{NN}$ --- $\sqrt3$ irrational is
\[
  \forall p,q \in \mathit{NN}.\quad
  q \neq 0 \;\Longrightarrow\; p\cdot p \neq 3\cdot(q\cdot q) ,
\]
that is, the goal
\begin{VNBsnip}
(FORALL p (IMPLIES (IN p NN) (FORALL q (IMPLIES (IN q NN)
  (IMPLIES (NOT (= q 0)) (NOT (= (* p p) (* 3 (* q q)))))))))
\end{VNBsnip}

\paragraph{The argument.}
Infinite descent.  Suppose the equation has a solution with $q \neq 0$, and
among all such solutions choose one whose numerator $w$ is \emph{least}.  Since
$3 \mid w\cdot w$ we have $3 \mid w$, say $w = 3k$; substituting and cancelling
the factor $3$ gives $q_0\cdot q_0 = 3\cdot(k\cdot k)$, whence $3 \mid q_0$, say
$q_0 = 3m_0$; cancelling once more gives $k\cdot k = 3\cdot(m_0\cdot m_0)$.  Thus
$(k,m_0)$ is again a solution, with $k \neq 0$ and $k < 3k = w$ --- a numerator
strictly below the least one.  Contradiction.

The entire mathematical content is the single implication
\begin{equation}\label{eq:3divsq}
  3 \mid p\cdot p \;\Longrightarrow\; 3 \mid p ;
\end{equation}
everything else --- clearing denominators, the two cancellations, the order
step $k < 3k$ --- is bookkeeping.

\paragraph{Why parity is not special.}
The $\sqrt2$ proof is the same skeleton with \emph{even} for \emph{divisible by
$3$}: there \eqref{eq:3divsq} reads ``$p\cdot p$ even $\Rightarrow p$ even'' and
follows from parity, every natural being $2k$ or $2k+1$.  For $\sqrt3$ it follows
from the trichotomy --- every natural is $3k$, $3k+1$, or $3k+2$ --- together
with the fact that the three residues square to $0,1,1$ modulo $3$, so a square
is divisible by $3$ only when its root is.  Parity is the $r=2$ instance of a
uniform statement about residues modulo $r$; the $\sqrt3$ development is the
$\sqrt2$ development with the residue system widened by one, and nothing changes
in kind.

\paragraph{The lemma ladder.}
The mod-$3$ arithmetic (\texttt{theorem-library/nn-mod3-proof.scm}) is built by
induction from the two Peano recursion equations for $+$ and $\cdot$, in exact
parallel with the parity ladder (\texttt{nn-parity-proof.scm}).  Each lemma is
\emph{proven}; none is asserted.

\begin{center}
\begin{tabular}{lll}
\toprule
parity ($r=2$) & this proof ($r=3$) & statement \\
\midrule
\texttt{nn-two-mul-succ}     & \texttt{nn-three-mul-succ} & $r\cdot\mathrm{succ}(k) = \mathrm{succ}^{\,r}(r\cdot k)$ \\
\texttt{nn-parity}           & \texttt{nn-trichotomy-3}   & every $n$ is $r\cdot k + i$, residue $i<r$ \\
\texttt{nn-parity-exclusive} & \texttt{nn-3-not-succ}     & $r\cdot x$ is never $r\cdot y + 1$ \\
\texttt{nn-even-square}      & \texttt{nn-3-div-square}   & $r \mid p\cdot p \Rightarrow r \mid p$ \\
\bottomrule
\end{tabular}
\end{center}

\noindent For $r=3$ a single residue-inequality suffices: both nonzero residues
square to $r\cdot j + 1$, so only ``not $\equiv 1$'' is ever needed.

\paragraph{The descent, mechanized.}
The driver (\texttt{theorem-library/sqrt3-proof.scm}) states the goal and hands
the descent to three tools; the intermediate steps are elided.

\begin{VNBcode}
(sp (make-wff '(FORALL p (IMPLIES (IN p NN) (FORALL q (IMPLIES (IN q NN)
     (IMPLIES (NOT (= q 0)) (NOT (= (* p p) (* 3 (* q q)))))))))))
(di)(di)(di)(di)(di)(di)                 ; intro p,q, hyps; assume p*p = 3*(q*q)

(define R (minimize! '(pv) (guard 'pv) 'pv))    ; choose the LEAST numerator w
;; ... discharge minimize!'s two obligations: measure in NN; guard satisfiable

(define k (obtain (lambda () (fact 'nn-3-div-square w))))   ; 3|w  =>  w = 3k
;; ... q0*q0 = 3*(k*k) by have! ; 3|q0 => q0 = 3*m0 by obtain ; k*k = 3*(m0*m0)

(inst+ minf k)              ; minimality gives  w <= k
(fact 'nn-lt-triple k)      ; but k < 3k = w -- contradiction
\end{VNBcode}

\begin{itemize}
\item \texttt{minimize!} supplies the least-numerator witness together with its
  two honest obligations --- the measure lands in $\mathit{NN}$, and the guard
  is satisfiable --- and appeals to no choice principle.
\item \texttt{obtain} names each descended witness --- the $k$ with $w = 3k$,
  the $m_0$ with $q_0 = 3m_0$ --- by \emph{reading it off the proof state} (the
  variable left free once the existential is eliminated), never by guessing a
  machine-generated name.  This is what makes the descent survive a re-run.
\item \texttt{have!} discharges each intermediate equality and typing from
  context, so the script reads as the argument does: ``we have
  $q_0\cdot q_0 = 3\cdot(k\cdot k)$; we have $k \neq 0$; \ldots''.
\end{itemize}

\begin{remark}[The \texttt{minimize!} step]
\label{remark:minimize}
The enigmatic-looking \texttt{(minimize! '(pv) (guard 'pv) 'pv)} decodes as:
choose the single number \texttt{pv} satisfying the descent guard, with
\texttt{pv} itself as the quantity to minimise --- that is, take the
\emph{least} numerator meeting the guard.

The second argument to \texttt{minimize!} is the \textbf{guard formula} itself:
an ordinary VNB formula, free in the variables being chosen (here \texttt{pv}),
stating the property those values must satisfy.  \texttt{guard} is \emph{not} a
built-in and is nothing special to \texttt{minimize!} --- it is a helper the
proof author writes (named \texttt{s3-guard} in \texttt{sqrt3-proof.scm}), a
one-argument function that takes a term and returns the guard formula about it.
So \texttt{(guard 'pv)} yields the formula with \texttt{pv} free --- exactly what
\texttt{minimize!} wants --- and the \emph{same} helper, called later at the
chosen witness $w$ as \texttt{(guard $w$)}, restates that property of $w$.  Here
the property is ``\texttt{pv} is a natural number that is the numerator of a
fraction whose square is~$3$''; written out, the helper is just a formula
constructor:

\begin{VNBcode}
(define (s3-guard pv)
  (list 'AND (list 'IN pv 'NN)
        (list 'FORSOME 'q_
          (list 'AND (list 'IN 'q_ 'NN)
                (list 'AND (list 'NOT (list '= 'q_ 0))
                      (list '= (list '* pv pv) (list '* 3 (list '* 'q_ 'q_))))))))
\end{VNBcode}

Passing the formula literally would work just as well; the helper only spares the
author from writing it twice.  In general,
\[
  \texttt{(minimize! '($v_1$ \ldots\ $v_k$) GUARD MEASURE)}
\]
means ``choose $v_1,\ldots,v_k$ satisfying the formula \texttt{GUARD} with the
$\mathit{NN}$-valued term \texttt{MEASURE} as small as possible.''  On the main
branch it lands two hypotheses, for fresh eigenconstants $w_1,\ldots,w_k$: that
\texttt{GUARD} holds of them, and that nothing satisfying \texttt{GUARD} has a
strictly smaller \texttt{MEASURE}.  It leaves exactly the two obligations any
minimisation owes --- that \texttt{MEASURE} lands in $\mathit{NN}$ wherever
\texttt{GUARD} holds, and that \texttt{GUARD} holds somewhere.  Everything
between --- forming the value set, well-ordering it, unpacking the least
witness, restating minimality over the $v_i$ rather than the set --- it does for
you.  Every $v_i$ must occur in \texttt{MEASURE} (a variable the measure ignores
is not one you are minimising over).

It adds no kernel rule --- it chains \texttt{cut}, \texttt{di}, \texttt{ai},
\texttt{ew}, \ldots --- and uses \emph{no choice principle}: its one appeal is
\texttt{nn-least-element}, the well-ordering of $\mathit{NN}$, and well-ordering
returns a genuine \emph{member} of the value set, from which the witness is read
back.  Stated instead as a first-order lemma, the same fact would force every
call site to tuple its variables into a Cartesian product, build a lambda, and
prove that lambda inhabits $\mathit{FUN}(U,\mathit{NN})$ --- a typing obligation
larger than the minimisation it buys.  As a tactic it quantifies over the $v_i$
at the meta level and never forms $U$ at all.
\end{remark}

\begin{remark}[\texttt{obtain} versus \texttt{vlet}]
\texttt{obtain} and the \texttt{choice} former of \texttt{vlet} are the same
mechanism --- eliminate an existential and read its witness off by
free-variable difference --- but they differ in what they take.  \texttt{obtain}
takes a \emph{lane}: it runs a forward step and names the witness that step
\emph{just produced}.  That is the right fit inside a descent, where the relevant
existential is freshly landed and typically sits beside others: at the
``$3 \mid w$'' step the conclusion $\exists k.\,w = 3k$ coexists in context with
the evidence $\exists m.\,w\cdot w = 3m$ cut in to derive it, and \texttt{obtain}
selects the new one automatically.  Concretely, that one step reads
\begin{VNBsnip}
;; as written in the proof:
(define k (obtain (lambda () (fact 'nn-3-div-square w))))
;; the vlet alternative -- land the existential, then name it by its body:
(fact 'nn-3-div-square w)
(define k (vlet--choice! 'k '(AND (IN k NN) (= w (* 3 k)))))
\end{VNBsnip}
The \texttt{vlet} body names the eigenconstant $w$, a value produced at run
time, so this is the procedure \texttt{vlet--choice!}, not the
\texttt{(vlet \ldots)} macro (whose body is a fixed literal pattern).
\texttt{vlet} instead destructures a
\emph{named, stable} formula: its \texttt{match} former binds the holes of a
pattern (``the $a$ and $b$ in $x\cdot b = a$''), and its \texttt{choice} former
eliminates a \emph{designated} existential.  The rule of thumb: \texttt{obtain}
to name whatever a forward step produced, \texttt{vlet} to take apart a formula
you can already name.
\end{remark}

\begin{remark}
The proof adds no asserted step of its own: the mod-$3$ ladder is proven and the
descent cites it, exactly as \texttt{sqrt2-proof.scm} cites the parity ladder.
A new irrationality --- $\sqrt r$ for any non-square $r$ --- is then a matter of
widening the residue ladder, not of inventing an argument.  That division of
labour is the lesson: the author states the intermediate claims, and the machine
supplies the least witness (\texttt{minimize!}), the fresh names
(\texttt{obtain}), and the bookkeeping (\texttt{have!}, \texttt{from-context!}).
\end{remark}

\paragraph{The same proof at the interactive surface.}
By the transcript principle (Section~\ref{sec:two-surfaces}), most of this
script is \emph{already} a sequence of interactive short forms: \texttt{di},
\texttt{ai}, \texttt{ew}, \texttt{fact}, \texttt{cut}, \texttt{subst},
\texttt{crs}, \texttt{ass}, and \texttt{minimize!} are exactly what a user
types.  What the driver adds is the orchestration a human otherwise does by eye,
and it is of two kinds.  \emph{Focus}: after a branching command one selects the
next open leaf from the frontier, where the script calls \texttt{dk-focus!}.
\emph{Names}: each \texttt{(ai\ (FORSOME \ldots))} skolemizes to a machine-chosen
eigenconstant, printed in the resulting sequent, which the user reads and types
into the following commands.  The $3\mid w$ step, done by hand, is
\begin{VNBsnip}
(fact 'nn-3-div-square w)                          ; lands  exists k. w = 3k
(ai '(FORSOME k (AND (IN k NN) (= w (* 3 k)))))    ; skolemize; sequent now shows k_899
;; ... read k_899 off the sequent; type it into the subsequent (subst ...) / (fact ...)
\end{VNBsnip}
\texttt{obtain} is exactly this pair automated --- run the forward step, read the
fresh name back by free-variable difference --- so that a \emph{script}, unlike a
transcript with \texttt{k\_899} hard-coded, replays unchanged when the
eigenvariable counter lands on a different number.

%% =========================================================
